\documentclass[11pt,a4paper]{article}
% The preamble of the exercise book (« Exercices »), exercises/exercices.fr.tex.
%
% Loaded as `% The preamble of the exercise book (« Exercices »), exercises/exercices.fr.tex.
%
% Loaded as `% The preamble of the exercise book (« Exercices »), exercises/exercices.fr.tex.
%
% Loaded as `\input{../transcripts/preamble/exercices}` from `exercises/`. It
% takes from germain.sty what a printed book needs — the fonts, the maths
% packages, French babel, the page geometry — and none of the transcription's
% apparatus: the book restates, it does not record, so \page, \ill, \uncertain
% and the rest have no business in it and are left undefined, so that a stray
% one is an error here as it is in `scripts/render.mjs`.
%
% The four constructs below are the book's own, fixed by the contract between
% the file and the site. Each is recognised by `scripts/render.mjs`, which
% renders the same file to `public/exercises/exercices.html`; a construct
% added here must be added there too, or rendering fails loudly.
%
%   \livretitre                         the title block
%   \chapitre{<group id>}               which mathematician the chapter is (prints nothing)
%   \begin{exercice}{<titre>}{<niveau>}{<étoiles>}
%                                       a numbered exercise; Lycée, L1, L2 or L3,
%                                       and a difficulty of 1, 2 or 3 stars
%   \manuscrit{<volume id>}{<v1>}{<v2>}  the Gallica views it comes from
%   \begin{solution}                    its solution, printed after it

% \ProvidesFile rather than \ProvidesPackage: the file is \input, not
% \usepackage'd, and the latter then warns that no package was requested.
\ProvidesFile{exercices.sty}[2026/09/22 Exercise book on mathematicians' manuscripts in Gallica]

\usepackage{amsmath,amssymb,amsthm}
\usepackage{mathrsfs}
\usepackage[french]{babel}
\usepackage{fontspec}

% Latin Modern lacks Greek and some scribal signs; `glyphs.sty` sets them in
% fallback fonts, as for the transcriptions. The path is from `exercises/`.
\input{../transcripts/preamble/glyphs.sty}

\usepackage{geometry}
\geometry{a4paper,margin=25mm}
\usepackage{xcolor}
\setlength{\emergencystretch}{3em}
\usepackage{hyperref}
\hypersetup{colorlinks=true,linkcolor=black,urlcolor=black,pdfborder={0 0 0}}

\definecolor{exgris}{gray}{0.40}
\definecolor{exstatut}{rgb}{0.63,0.29,0.18}

% --- Title ------------------------------------------------------------------
% The status line says what the book is: a first draft by a machine, like
% every reading on the site, and it is printed on the first page rather than
% left to a colophon nobody reads.

\newcommand{\livretitre}{%
  \begin{center}
    {\LARGE\bfseries Exercices}\\[6pt]
    {\large\itshape d'après les manuscrits de mathématiciens numérisés dans Gallica}\\[10pt]
    {\small\bfseries\color{exstatut}Version de travail --- énoncés et solutions
      rédigés par une machine, non relus}\\[4pt]
    {\footnotesize\color{exgris}Les manuscrits cités sont lisibles sur
      \url{https://germain.commutator.io} et sur gallica.bnf.fr.
      Source gallica.bnf.fr / Bibliothèque nationale de France.}
  \end{center}
  \bigskip}

% --- Chapters ---------------------------------------------------------------
% The mathematician a chapter is, as the catalogue's group id (`maurolico`,
% `germain`). The site filters on it; the page needs nothing from it.

\newcommand{\chapitre}[1]{\ignorespaces}

% --- Shelfmarks -------------------------------------------------------------
% `naf-5176` → « NAF 5176 »: the catalogue's slug back to the holder's
% shelfmark, as the site prints it. An unknown prefix is an error, not a guess.

\ExplSyntaxOn
\cs_new:Npn \__exercices_shelfmark:w #1-#2 \q_stop
  {
    \str_case:nnF {#1}
      {
        {fr} {Français}
        {naf} {NAF}
        {nal} {NAL}
        {latin} {Latin}
        {arsenal} {Arsenal}
        {rothschild} {Rothschild}
        {enpc} {ENPC}
      }
      { \PackageError {exercices} {Unknown~volume~prefix~`#1'~in~`#1-#2'} {} }
    \nobreakspace #2
  }
\NewDocumentCommand \exshelfmark { m } { \__exercices_shelfmark:w #1 \q_stop }

% The four levels, and nothing else: the site's filter knows these four.
\NewDocumentCommand \exniveau { m }
  {
    \str_case:nnF {#1}
      { {Lycée} {} {L1} {} {L2} {} {L3} {} }
      { \PackageError {exercices} {Unknown~level~`#1'~---~use~Lycée,~L1,~L2~or~L3} {} }
    #1
  }

% The difficulty inside a level: 1, 2 or 3 stars.
\NewDocumentCommand \exetoiles { m }
  {
    \str_case:nnF {#1}
      {
        {1} { $\bigstar$ }
        {2} { $\bigstar\bigstar$ }
        {3} { $\bigstar\bigstar\bigstar$ }
      }
      { \PackageError {exercices} {Unknown~difficulty~`#1'~---~use~1,~2~or~3} {} }
  }
\ExplSyntaxOff

% --- Exercises and solutions ------------------------------------------------
% Numbered through the whole book, not per chapter: the site numbers them the
% same way, and « exercice 7 » must name one exercise in both.

\newcounter{exercice}

\newenvironment{exercice}[3]{%
  \par\bigskip\goodbreak
  \refstepcounter{exercice}%
  \noindent{\bfseries Exercice~\theexercice\ (\exniveau{#2}, \exetoiles{#3}) --- #1}\par
  \nopagebreak\smallskip
}{\par\medskip}

\newcommand{\manuscrit}[3]{%
  \par\noindent{\small\color{exgris}Source~: \exshelfmark{#1},
    \ifnum#2=#3\relax vue~#2\else vues~#2--#3\fi}\par
  \nopagebreak\smallskip\ignorespaces}

\newenvironment{solution}{%
  \par\smallskip\noindent{\bfseries\itshape Solution}\par\nopagebreak\smallskip
}{\par\smallskip\noindent\hfill\rule{0.3\linewidth}{0.3pt}\hfill\null\par\medskip}

\endinput
` from `exercises/`. It
% takes from germain.sty what a printed book needs — the fonts, the maths
% packages, French babel, the page geometry — and none of the transcription's
% apparatus: the book restates, it does not record, so \page, \ill, \uncertain
% and the rest have no business in it and are left undefined, so that a stray
% one is an error here as it is in `scripts/render.mjs`.
%
% The four constructs below are the book's own, fixed by the contract between
% the file and the site. Each is recognised by `scripts/render.mjs`, which
% renders the same file to `public/exercises/exercices.html`; a construct
% added here must be added there too, or rendering fails loudly.
%
%   \livretitre                         the title block
%   \chapitre{<group id>}               which mathematician the chapter is (prints nothing)
%   \begin{exercice}{<titre>}{<niveau>}{<étoiles>}
%                                       a numbered exercise; Lycée, L1, L2 or L3,
%                                       and a difficulty of 1, 2 or 3 stars
%   \manuscrit{<volume id>}{<v1>}{<v2>}  the Gallica views it comes from
%   \begin{solution}                    its solution, printed after it

% \ProvidesFile rather than \ProvidesPackage: the file is \input, not
% \usepackage'd, and the latter then warns that no package was requested.
\ProvidesFile{exercices.sty}[2026/09/22 Exercise book on mathematicians' manuscripts in Gallica]

\usepackage{amsmath,amssymb,amsthm}
\usepackage{mathrsfs}
\usepackage[french]{babel}
\usepackage{fontspec}

% Latin Modern lacks Greek and some scribal signs; `glyphs.sty` sets them in
% fallback fonts, as for the transcriptions. The path is from `exercises/`.
% Fallback fonts for the characters Latin Modern does not have.
%
% Input by `germain.sty` and `exercices.sty` right after fontspec — a
% `.sty`, not a `.tex`, because `scripts/pdf.mjs` compiles every `.tex` it
% finds under `transcripts/` as a document.
%
% The text face stays fontspec's default, Latin Modern. What it lacks, XeTeX
% drops with a line in the log (« Missing character ») and nothing on the page:
% that is how the Greek words the manuscripts quote, the barred p of per (ꝑ),
% the looped z (ʒ) and the rest vanished from the PDFs while the HTML reading
% views, set by the browser, showed them. So each fallback font gets a XeTeX
% character class, and a run of characters of that class switches the family
% to the font and back:
%
%   CMU Serif           Greek and Greek Extended, ʒ — the same Computer
%                       Modern design as Latin Modern, with polytonic Greek
%   Latin Modern Math   ∂ — the glyph the PDFs already print for $\partial$
%   Junicode            ꝑ, ꝓ, ǂ — the medievalists' font for scribal signs,
%                       scaled to Latin Modern's x-height; CMU has ǂ but not
%                       in italic, where the notes that cite it are set
%
% Only the family changes: a character in \textit or \textbf keeps its shape
% and series where the font has one. The transitions to and from every other
% class carry the tokens that class has towards an ordinary letter, so French
% spacing before « ; » or inside « » still holds next to a switched character.
% No grouping, so no brace or \endgroup can come out unbalanced; maths is left
% alone, since XeTeX inserts nothing there.
%
% To cover another character: check that one of the fonts has it (or add a
% family with \germain@fallback), then give it a \germain@fallbackchars line.

\ProvidesFile{glyphs.sty}[2026/09/23 Fallback fonts for characters Latin Modern lacks]

\makeatletter

\newfontfamily\germainCMUfont{cmun}[
  Extension=.otf, NFSSFamily=germaincmu,
  UprightFont=*rm, ItalicFont=*ti, SlantedFont=*sl,
  BoldFont=*bx, BoldItalicFont=*bi, BoldSlantedFont=*bl]
% One upright face only: named for every shape so NFSS never substitutes.
\newfontfamily\germainMathfont{latinmodern-math.otf}[
  NFSSFamily=germainmath,
  ItalicFont=latinmodern-math.otf, SlantedFont=latinmodern-math.otf,
  BoldFont=latinmodern-math.otf, BoldItalicFont=latinmodern-math.otf,
  BoldSlantedFont=latinmodern-math.otf]
\newfontfamily\germainJunicodefont{Junicode.ttf}[
  NFSSFamily=germainjunicode, Scale=MatchLowercase,
  ItalicFont=Junicode-Italic.ttf, SlantedFont=Junicode-Italic.ttf,
  BoldFont=Junicode-Bold.ttf, BoldItalicFont=Junicode-BoldItalic.ttf,
  BoldSlantedFont=Junicode-BoldItalic.ttf]

% \germain@fallback{<NFSS family>} gives the family a character class, and
% the class a macro, \germain@in<class>, that switches to it.
\def\germain@fallback#1{%
  \expandafter\newXeTeXintercharclass\csname germain@class@#1\endcsname
  \expandafter\edef\csname germain@in\number\csname germain@class@#1\endcsname\endcsname{%
    \noexpand\edef\noexpand\germain@prev{\noexpand\f@family}%
    \noexpand\fontfamily{#1}\noexpand\selectfont}}
% \germain@fallbackchars{<NFSS family>}{<first>}{<last>}: a range of code points.
\def\germain@fallbackchars#1#2#3{%
  \@tempcnta=#2\relax
  \loop
    \XeTeXcharclass\@tempcnta=\csname germain@class@#1\endcsname
    \ifnum\@tempcnta<#3\advance\@tempcnta\@ne\repeat}

\germain@fallback{germaincmu}
\germain@fallbackchars{germaincmu}{"0370}{"03FF}  % Greek and Coptic
\germain@fallbackchars{germaincmu}{"1F00}{"1FFF}  % Greek Extended
\germain@fallbackchars{germaincmu}{"0292}{"0292}  % ʒ
\germain@fallback{germainmath}
\germain@fallbackchars{germainmath}{"2202}{"2202} % ∂
\germain@fallback{germainjunicode}
\germain@fallbackchars{germainjunicode}{"A751}{"A751} % ꝑ
\germain@fallbackchars{germainjunicode}{"A753}{"A753} % ꝓ
\germain@fallbackchars{germainjunicode}{"01C2}{"01C2} % ǂ

\def\germain@prev{\rmdefault}
\def\germain@out{\fontfamily{\germain@prev}\selectfont}

% For every fallback class f and every other class k: f→k leaves the font and
% then does what an ordinary letter followed by k would; k→f does what k
% followed by an ordinary letter would, then enters it. Between two fallback
% classes, leave one and enter the other.
\def\germain@join{%
  \XeTeXinterchartokenstate=\@ne
  \@tempcntb=\z@
  \loop
    \ifcsname germain@in\number\@tempcntb\endcsname
      \germain@joinclass
    \fi
    \ifnum\@tempcntb<4095\advance\@tempcntb\@ne\repeat}
\def\germain@joinclass{%
  \@tempcnta=\z@
  \germain@joinloop}
\def\germain@joinloop{%
  \ifnum\@tempcnta=\@tempcntb\else
    \ifcsname germain@in\number\@tempcnta\endcsname
      \edef\germain@tmp{\noexpand\germain@out
        \expandafter\noexpand\csname germain@in\number\@tempcnta\endcsname}%
      \XeTeXinterchartoks\@tempcntb\@tempcnta=\expandafter{\germain@tmp}%
    \else
      \toks@=\expandafter{\the\XeTeXinterchartoks\z@\@tempcnta}%
      \edef\germain@tmp{\noexpand\germain@out\the\toks@}%
      \XeTeXinterchartoks\@tempcntb\@tempcnta=\expandafter{\germain@tmp}%
      \toks@=\expandafter{\the\XeTeXinterchartoks\@tempcnta\z@}%
      \edef\germain@tmp{\the\toks@
        \expandafter\noexpand\csname germain@in\number\@tempcntb\endcsname}%
      \XeTeXinterchartoks\@tempcnta\@tempcntb=\expandafter{\germain@tmp}%
    \fi
  \fi
  \ifnum\@tempcnta<4095
    \advance\@tempcnta\@ne
    \expandafter\germain@joinloop
  \fi}
% babel-french sets its own classes, and saves and restores the interchar
% state, each time a language is selected: join again after it does.
\addto\extrasfrench{\germain@join}
\@ifundefined{extrasenglish}{}{\addto\extrasenglish{\germain@join}}
\AtBeginDocument{\germain@join}

\makeatother


\usepackage{geometry}
\geometry{a4paper,margin=25mm}
\usepackage{xcolor}
\setlength{\emergencystretch}{3em}
\usepackage{hyperref}
\hypersetup{colorlinks=true,linkcolor=black,urlcolor=black,pdfborder={0 0 0}}

\definecolor{exgris}{gray}{0.40}
\definecolor{exstatut}{rgb}{0.63,0.29,0.18}

% --- Title ------------------------------------------------------------------
% The status line says what the book is: a first draft by a machine, like
% every reading on the site, and it is printed on the first page rather than
% left to a colophon nobody reads.

\newcommand{\livretitre}{%
  \begin{center}
    {\LARGE\bfseries Exercices}\\[6pt]
    {\large\itshape d'après les manuscrits de mathématiciens numérisés dans Gallica}\\[10pt]
    {\small\bfseries\color{exstatut}Version de travail --- énoncés et solutions
      rédigés par une machine, non relus}\\[4pt]
    {\footnotesize\color{exgris}Les manuscrits cités sont lisibles sur
      \url{https://germain.commutator.io} et sur gallica.bnf.fr.
      Source gallica.bnf.fr / Bibliothèque nationale de France.}
  \end{center}
  \bigskip}

% --- Chapters ---------------------------------------------------------------
% The mathematician a chapter is, as the catalogue's group id (`maurolico`,
% `germain`). The site filters on it; the page needs nothing from it.

\newcommand{\chapitre}[1]{\ignorespaces}

% --- Shelfmarks -------------------------------------------------------------
% `naf-5176` → « NAF 5176 »: the catalogue's slug back to the holder's
% shelfmark, as the site prints it. An unknown prefix is an error, not a guess.

\ExplSyntaxOn
\cs_new:Npn \__exercices_shelfmark:w #1-#2 \q_stop
  {
    \str_case:nnF {#1}
      {
        {fr} {Français}
        {naf} {NAF}
        {nal} {NAL}
        {latin} {Latin}
        {arsenal} {Arsenal}
        {rothschild} {Rothschild}
        {enpc} {ENPC}
      }
      { \PackageError {exercices} {Unknown~volume~prefix~`#1'~in~`#1-#2'} {} }
    \nobreakspace #2
  }
\NewDocumentCommand \exshelfmark { m } { \__exercices_shelfmark:w #1 \q_stop }

% The four levels, and nothing else: the site's filter knows these four.
\NewDocumentCommand \exniveau { m }
  {
    \str_case:nnF {#1}
      { {Lycée} {} {L1} {} {L2} {} {L3} {} }
      { \PackageError {exercices} {Unknown~level~`#1'~---~use~Lycée,~L1,~L2~or~L3} {} }
    #1
  }

% The difficulty inside a level: 1, 2 or 3 stars.
\NewDocumentCommand \exetoiles { m }
  {
    \str_case:nnF {#1}
      {
        {1} { $\bigstar$ }
        {2} { $\bigstar\bigstar$ }
        {3} { $\bigstar\bigstar\bigstar$ }
      }
      { \PackageError {exercices} {Unknown~difficulty~`#1'~---~use~1,~2~or~3} {} }
  }
\ExplSyntaxOff

% --- Exercises and solutions ------------------------------------------------
% Numbered through the whole book, not per chapter: the site numbers them the
% same way, and « exercice 7 » must name one exercise in both.

\newcounter{exercice}

\newenvironment{exercice}[3]{%
  \par\bigskip\goodbreak
  \refstepcounter{exercice}%
  \noindent{\bfseries Exercice~\theexercice\ (\exniveau{#2}, \exetoiles{#3}) --- #1}\par
  \nopagebreak\smallskip
}{\par\medskip}

\newcommand{\manuscrit}[3]{%
  \par\noindent{\small\color{exgris}Source~: \exshelfmark{#1},
    \ifnum#2=#3\relax vue~#2\else vues~#2--#3\fi}\par
  \nopagebreak\smallskip\ignorespaces}

\newenvironment{solution}{%
  \par\smallskip\noindent{\bfseries\itshape Solution}\par\nopagebreak\smallskip
}{\par\smallskip\noindent\hfill\rule{0.3\linewidth}{0.3pt}\hfill\null\par\medskip}

\endinput
` from `exercises/`. It
% takes from germain.sty what a printed book needs — the fonts, the maths
% packages, French babel, the page geometry — and none of the transcription's
% apparatus: the book restates, it does not record, so \page, \ill, \uncertain
% and the rest have no business in it and are left undefined, so that a stray
% one is an error here as it is in `scripts/render.mjs`.
%
% The four constructs below are the book's own, fixed by the contract between
% the file and the site. Each is recognised by `scripts/render.mjs`, which
% renders the same file to `public/exercises/exercices.html`; a construct
% added here must be added there too, or rendering fails loudly.
%
%   \livretitre                         the title block
%   \chapitre{<group id>}               which mathematician the chapter is (prints nothing)
%   \begin{exercice}{<titre>}{<niveau>}{<étoiles>}
%                                       a numbered exercise; Lycée, L1, L2 or L3,
%                                       and a difficulty of 1, 2 or 3 stars
%   \manuscrit{<volume id>}{<v1>}{<v2>}  the Gallica views it comes from
%   \begin{solution}                    its solution, printed after it

% \ProvidesFile rather than \ProvidesPackage: the file is \input, not
% \usepackage'd, and the latter then warns that no package was requested.
\ProvidesFile{exercices.sty}[2026/09/22 Exercise book on mathematicians' manuscripts in Gallica]

\usepackage{amsmath,amssymb,amsthm}
\usepackage{mathrsfs}
\usepackage[french]{babel}
\usepackage{fontspec}

% Latin Modern lacks Greek and some scribal signs; `glyphs.sty` sets them in
% fallback fonts, as for the transcriptions. The path is from `exercises/`.
% Fallback fonts for the characters Latin Modern does not have.
%
% Input by `germain.sty` and `exercices.sty` right after fontspec — a
% `.sty`, not a `.tex`, because `scripts/pdf.mjs` compiles every `.tex` it
% finds under `transcripts/` as a document.
%
% The text face stays fontspec's default, Latin Modern. What it lacks, XeTeX
% drops with a line in the log (« Missing character ») and nothing on the page:
% that is how the Greek words the manuscripts quote, the barred p of per (ꝑ),
% the looped z (ʒ) and the rest vanished from the PDFs while the HTML reading
% views, set by the browser, showed them. So each fallback font gets a XeTeX
% character class, and a run of characters of that class switches the family
% to the font and back:
%
%   CMU Serif           Greek and Greek Extended, ʒ — the same Computer
%                       Modern design as Latin Modern, with polytonic Greek
%   Latin Modern Math   ∂ — the glyph the PDFs already print for $\partial$
%   Junicode            ꝑ, ꝓ, ǂ — the medievalists' font for scribal signs,
%                       scaled to Latin Modern's x-height; CMU has ǂ but not
%                       in italic, where the notes that cite it are set
%
% Only the family changes: a character in \textit or \textbf keeps its shape
% and series where the font has one. The transitions to and from every other
% class carry the tokens that class has towards an ordinary letter, so French
% spacing before « ; » or inside « » still holds next to a switched character.
% No grouping, so no brace or \endgroup can come out unbalanced; maths is left
% alone, since XeTeX inserts nothing there.
%
% To cover another character: check that one of the fonts has it (or add a
% family with \germain@fallback), then give it a \germain@fallbackchars line.

\ProvidesFile{glyphs.sty}[2026/09/23 Fallback fonts for characters Latin Modern lacks]

\makeatletter

\newfontfamily\germainCMUfont{cmun}[
  Extension=.otf, NFSSFamily=germaincmu,
  UprightFont=*rm, ItalicFont=*ti, SlantedFont=*sl,
  BoldFont=*bx, BoldItalicFont=*bi, BoldSlantedFont=*bl]
% One upright face only: named for every shape so NFSS never substitutes.
\newfontfamily\germainMathfont{latinmodern-math.otf}[
  NFSSFamily=germainmath,
  ItalicFont=latinmodern-math.otf, SlantedFont=latinmodern-math.otf,
  BoldFont=latinmodern-math.otf, BoldItalicFont=latinmodern-math.otf,
  BoldSlantedFont=latinmodern-math.otf]
\newfontfamily\germainJunicodefont{Junicode.ttf}[
  NFSSFamily=germainjunicode, Scale=MatchLowercase,
  ItalicFont=Junicode-Italic.ttf, SlantedFont=Junicode-Italic.ttf,
  BoldFont=Junicode-Bold.ttf, BoldItalicFont=Junicode-BoldItalic.ttf,
  BoldSlantedFont=Junicode-BoldItalic.ttf]

% \germain@fallback{<NFSS family>} gives the family a character class, and
% the class a macro, \germain@in<class>, that switches to it.
\def\germain@fallback#1{%
  \expandafter\newXeTeXintercharclass\csname germain@class@#1\endcsname
  \expandafter\edef\csname germain@in\number\csname germain@class@#1\endcsname\endcsname{%
    \noexpand\edef\noexpand\germain@prev{\noexpand\f@family}%
    \noexpand\fontfamily{#1}\noexpand\selectfont}}
% \germain@fallbackchars{<NFSS family>}{<first>}{<last>}: a range of code points.
\def\germain@fallbackchars#1#2#3{%
  \@tempcnta=#2\relax
  \loop
    \XeTeXcharclass\@tempcnta=\csname germain@class@#1\endcsname
    \ifnum\@tempcnta<#3\advance\@tempcnta\@ne\repeat}

\germain@fallback{germaincmu}
\germain@fallbackchars{germaincmu}{"0370}{"03FF}  % Greek and Coptic
\germain@fallbackchars{germaincmu}{"1F00}{"1FFF}  % Greek Extended
\germain@fallbackchars{germaincmu}{"0292}{"0292}  % ʒ
\germain@fallback{germainmath}
\germain@fallbackchars{germainmath}{"2202}{"2202} % ∂
\germain@fallback{germainjunicode}
\germain@fallbackchars{germainjunicode}{"A751}{"A751} % ꝑ
\germain@fallbackchars{germainjunicode}{"A753}{"A753} % ꝓ
\germain@fallbackchars{germainjunicode}{"01C2}{"01C2} % ǂ

\def\germain@prev{\rmdefault}
\def\germain@out{\fontfamily{\germain@prev}\selectfont}

% For every fallback class f and every other class k: f→k leaves the font and
% then does what an ordinary letter followed by k would; k→f does what k
% followed by an ordinary letter would, then enters it. Between two fallback
% classes, leave one and enter the other.
\def\germain@join{%
  \XeTeXinterchartokenstate=\@ne
  \@tempcntb=\z@
  \loop
    \ifcsname germain@in\number\@tempcntb\endcsname
      \germain@joinclass
    \fi
    \ifnum\@tempcntb<4095\advance\@tempcntb\@ne\repeat}
\def\germain@joinclass{%
  \@tempcnta=\z@
  \germain@joinloop}
\def\germain@joinloop{%
  \ifnum\@tempcnta=\@tempcntb\else
    \ifcsname germain@in\number\@tempcnta\endcsname
      \edef\germain@tmp{\noexpand\germain@out
        \expandafter\noexpand\csname germain@in\number\@tempcnta\endcsname}%
      \XeTeXinterchartoks\@tempcntb\@tempcnta=\expandafter{\germain@tmp}%
    \else
      \toks@=\expandafter{\the\XeTeXinterchartoks\z@\@tempcnta}%
      \edef\germain@tmp{\noexpand\germain@out\the\toks@}%
      \XeTeXinterchartoks\@tempcntb\@tempcnta=\expandafter{\germain@tmp}%
      \toks@=\expandafter{\the\XeTeXinterchartoks\@tempcnta\z@}%
      \edef\germain@tmp{\the\toks@
        \expandafter\noexpand\csname germain@in\number\@tempcntb\endcsname}%
      \XeTeXinterchartoks\@tempcnta\@tempcntb=\expandafter{\germain@tmp}%
    \fi
  \fi
  \ifnum\@tempcnta<4095
    \advance\@tempcnta\@ne
    \expandafter\germain@joinloop
  \fi}
% babel-french sets its own classes, and saves and restores the interchar
% state, each time a language is selected: join again after it does.
\addto\extrasfrench{\germain@join}
\@ifundefined{extrasenglish}{}{\addto\extrasenglish{\germain@join}}
\AtBeginDocument{\germain@join}

\makeatother


\usepackage{geometry}
\geometry{a4paper,margin=25mm}
\usepackage{xcolor}
\setlength{\emergencystretch}{3em}
\usepackage{hyperref}
\hypersetup{colorlinks=true,linkcolor=black,urlcolor=black,pdfborder={0 0 0}}

\definecolor{exgris}{gray}{0.40}
\definecolor{exstatut}{rgb}{0.63,0.29,0.18}

% --- Title ------------------------------------------------------------------
% The status line says what the book is: a first draft by a machine, like
% every reading on the site, and it is printed on the first page rather than
% left to a colophon nobody reads.

\newcommand{\livretitre}{%
  \begin{center}
    {\LARGE\bfseries Exercices}\\[6pt]
    {\large\itshape d'après les manuscrits de mathématiciens numérisés dans Gallica}\\[10pt]
    {\small\bfseries\color{exstatut}Version de travail --- énoncés et solutions
      rédigés par une machine, non relus}\\[4pt]
    {\footnotesize\color{exgris}Les manuscrits cités sont lisibles sur
      \url{https://germain.commutator.io} et sur gallica.bnf.fr.
      Source gallica.bnf.fr / Bibliothèque nationale de France.}
  \end{center}
  \bigskip}

% --- Chapters ---------------------------------------------------------------
% The mathematician a chapter is, as the catalogue's group id (`maurolico`,
% `germain`). The site filters on it; the page needs nothing from it.

\newcommand{\chapitre}[1]{\ignorespaces}

% --- Shelfmarks -------------------------------------------------------------
% `naf-5176` → « NAF 5176 »: the catalogue's slug back to the holder's
% shelfmark, as the site prints it. An unknown prefix is an error, not a guess.

\ExplSyntaxOn
\cs_new:Npn \__exercices_shelfmark:w #1-#2 \q_stop
  {
    \str_case:nnF {#1}
      {
        {fr} {Français}
        {naf} {NAF}
        {nal} {NAL}
        {latin} {Latin}
        {arsenal} {Arsenal}
        {rothschild} {Rothschild}
        {enpc} {ENPC}
      }
      { \PackageError {exercices} {Unknown~volume~prefix~`#1'~in~`#1-#2'} {} }
    \nobreakspace #2
  }
\NewDocumentCommand \exshelfmark { m } { \__exercices_shelfmark:w #1 \q_stop }

% The four levels, and nothing else: the site's filter knows these four.
\NewDocumentCommand \exniveau { m }
  {
    \str_case:nnF {#1}
      { {Lycée} {} {L1} {} {L2} {} {L3} {} }
      { \PackageError {exercices} {Unknown~level~`#1'~---~use~Lycée,~L1,~L2~or~L3} {} }
    #1
  }

% The difficulty inside a level: 1, 2 or 3 stars.
\NewDocumentCommand \exetoiles { m }
  {
    \str_case:nnF {#1}
      {
        {1} { $\bigstar$ }
        {2} { $\bigstar\bigstar$ }
        {3} { $\bigstar\bigstar\bigstar$ }
      }
      { \PackageError {exercices} {Unknown~difficulty~`#1'~---~use~1,~2~or~3} {} }
  }
\ExplSyntaxOff

% --- Exercises and solutions ------------------------------------------------
% Numbered through the whole book, not per chapter: the site numbers them the
% same way, and « exercice 7 » must name one exercise in both.

\newcounter{exercice}

\newenvironment{exercice}[3]{%
  \par\bigskip\goodbreak
  \refstepcounter{exercice}%
  \noindent{\bfseries Exercice~\theexercice\ (\exniveau{#2}, \exetoiles{#3}) --- #1}\par
  \nopagebreak\smallskip
}{\par\medskip}

\newcommand{\manuscrit}[3]{%
  \par\noindent{\small\color{exgris}Source~: \exshelfmark{#1},
    \ifnum#2=#3\relax vue~#2\else vues~#2--#3\fi}\par
  \nopagebreak\smallskip\ignorespaces}

\newenvironment{solution}{%
  \par\smallskip\noindent{\bfseries\itshape Solution}\par\nopagebreak\smallskip
}{\par\smallskip\noindent\hfill\rule{0.3\linewidth}{0.3pt}\hfill\null\par\medskip}

\endinput

\begin{document}

\livretitre

\section*{Avant-propos}

Ce livre rassemble des exercices originaux, avec leurs solutions complètes,
tirés des manuscrits que ce site transcrit et lit. Chaque chapitre porte le nom
d'un mathématicien ; il s'ouvre sur une courte présentation des feuillets qu'il
lit, puis propose des exercices qui partent de ce que ces feuillets font
réellement : une construction, un calcul, une affirmation, parfois une erreur.
On y refait le calcul du scripteur pour voir où il se trompe, on démontre ce que
la page se contente d'affirmer, on traduit la notation du temps avant de
résoudre, on compare la méthode du feuillet à celle qu'on emploierait
aujourd'hui. La ligne « Manuscrit » de chaque exercice renvoie aux vues de
Gallica sur lesquelles il repose, et l'on peut lire ces vues, leur transcription
et leur lecture modernisée sur le site.

C'est un choix, et non un inventaire : une quarantaine d'exercices, retenus pour
que chaque mathématicien soit représenté par ce que ses feuillets ont de plus
parlant, et que chaque chapitre aille du plus abordable au plus exigeant.

Les exercices ne recopient pas les lectures modernisées : ils s'en servent comme
d'un guide vers les feuillets, et leurs solutions refont chaque calcul. Tous les
nombres des solutions ont été recalculés pour ce livre. Quand un exercice porte
sur une erreur du feuillet, la solution donne la valeur juste et dit exactement
où le feuillet s'en écarte ; quand elle s'écarte de la lecture modernisée, elle
le dit aussi.

Chaque exercice porte un niveau : \emph{Lycée} (terminale, spécialité
mathématiques), \emph{L1}, \emph{L2} ou \emph{L3} (les trois années de licence
de mathématiques). Une série, une intégrale, une équation différentielle ou une
action de groupe ne sont jamais rangées au niveau Lycée ; à l'inverse, bien des
feuillets du XVII\textsuperscript{e} siècle se lisent avec les seuls outils
d'une classe de terminale. Il porte aussi une difficulté, d'une à trois
étoiles, qui se lit à l'intérieur de son niveau : une étoile pour une
application directe, deux pour une idée à trouver ou plusieurs étapes à
enchaîner, trois pour un exercice long ou délicat.

Un avertissement, qui vaut pour tout le livre. Les transcriptions et les
lectures modernisées sur lesquelles ces exercices s'appuient sont des lectures
automatiques, faites par un modèle de langage et non encore vérifiées par une
personne sur les images. Elles peuvent se tromper sur un chiffre, sur une
lettre de figure, sur l'ordre des feuillets. Le lecteur qui trouve un désaccord
entre un exercice et le feuillet doit croire le feuillet.

Les chapitres suivent l'ordre du catalogue du site, du XVII\textsuperscript{e} au
XIX\textsuperscript{e} siècle. Un exercice est rangé sous un mathématicien quand le
feuillet le nomme comme l'auteur de ce qu'il traite (« ce beau problème de
M. Roberval »), et sinon sous le mathématicien dont le catalogue porte le nom pour
son volume : les critiques anonymes de NAF 5161 sont ainsi au chapitre de
Descartes, qu'elles attaquent, et les feuillets de jeux de NAF 5176, qui ne
nomment personne, au chapitre de Mersenne. Ce rangement n'est pas une attribution.
Les attributions sont celles des feuillets et des lectures, et rien de plus :
aucun exercice ne tranche une question de priorité, et les rapprochements avec des
ouvrages imprimés, cités dans la bibliographie, sont des repères et non des
sources des feuillets. Un volume lu en entier n'a pas d'exercice : NAF 15383, la
lettre de Blaise Pascal à sa sœur Gilberte du 31 janvier 1643, avec une quittance
de rente de 1674, ne porte aucune mathématique --- tout au plus une règle de
trois.

\section*{Francesco Maurolico (1494--1575)}
\chapitre{maurolico}

Latin 17859 est la copie, faite au XVII\textsuperscript{e} siècle d'une seule
main, du volume que Francesco Maurolico fit imprimer à Messine en 1558 sur la
sphère : un \emph{Discours sur la sphère}, ses versions d'Autolycus, de Théodose
et des \emph{Phénomènes} d'Euclide, trois tables (sinus, tangentes dites « table
féconde », sécantes dites « table bienfaisante ») pour un rayon de $100\,000$,
huit problèmes chiffrés, la fin du premier livre et le second livre entier de
ses \emph{Sphériques}, et un commentaire des règles de Regiomontanus daté d'août
1530. Dans tout le chapitre, le « sinus » d'un arc est $100\,000$ fois notre
sinus, et « rejeter cinq figures » veut dire diviser par $100\,000$.

\begin{exercice}{Les soixantièmes de la table des sinus}{Lycée}{1}
\manuscrit{latin-17859}{16}{16}
La « tabella sinus recti » donne, pour chaque degré $n$, l'entier
$E(n)$ le plus proche de $100\,000\sin n^\circ$, la différence
$D(n) = E(n) - E(n-1)$, et le « soixantième » $D(n)/60$, écrit sous la forme
$a{:}b$, qui signifie $a + b/60$.

\begin{enumerate}
\item Calculer $E(23)$, $E(24)$, $D(24)$, et écrire $D(24)/60$ sous la forme
  $a{:}b$.
\item Expliquer pourquoi $E(23) + 30 \cdot \frac{D(24)}{60}$ est une valeur
  approchée de $100\,000 \sin 23^\circ 30'$. La calculer, la comparer à la valeur
  exacte, et à la valeur $39\,875$ qu'emploient les problèmes du volume. Dans
  quel sens l'interpolation se trompe-t-elle, et pourquoi ?
\item Le feuillet porte, à la ligne 59, $E(58) = 84\,805$, $E(59) = 85\,717$ et
  le soixantième « $15{:}21$ » ; à la ligne 86, $E(85) = 99\,620$,
  $E(86) = 99\,756$, la différence « $139$ » et le soixantième « $2{:}16$ ».
  Trouver les deux fautes, et dire chaque fois laquelle des colonnes est juste.
\end{enumerate}
\end{exercice}

\begin{solution}
\begin{enumerate}
\item $\sin 23^\circ = 0{,}390731\ldots$ et $\sin 24^\circ = 0{,}406737\ldots$,
  donc $E(23) = 39\,073$, $E(24) = 40\,674$ et $D(24) = 1601$. Comme
  $1601 = 26 \times 60 + 41$, le soixantième vaut $26 + \frac{41}{60}$, soit
  $26{:}41$ : c'est l'accroissement du sinus par minute d'arc entre $23^\circ$ et
  $24^\circ$.
\item Entre $23^\circ$ et $24^\circ$ on remplace l'arc de sinusoïde par sa corde :
  chaque minute ajoute $D(24)/60$, et $30'$ ajoutent trente fois cette quantité.
  On trouve $39\,073 + 30 \times 26{,}683 = 39\,073 + 800{,}5 = 39\,873{,}5$,
  soit $39\,874$ en arrondissant. La valeur exacte est
  $100\,000 \sin 23^\circ 30' = 39\,874{,}91$, soit $39\,875$ : c'est la valeur
  des problèmes, qui est donc la bonne. L'interpolation linéaire donne une
  valeur trop faible d'environ une unité et demie, parce que le sinus est
  concave sur $[0^\circ, 90^\circ]$ : sur un intervalle, sa courbe est au-dessus
  de sa corde.
\item $85\,717 - 84\,805 = 912 = 15 \times 60 + 12$ : le soixantième est
  $15{:}12$ ; le feuillet a interverti les chiffres en écrivant $15{:}21$. Les
  deux sinus et la différence sont justes. À la ligne 86,
  $99\,756 - 99\,620 = 136$, et non $139$ ; or $136 = 2 \times 60 + 16$, de sorte
  que le soixantième $2{:}16$ est bien celui de $136$. C'est donc la colonne des
  différences qui est fautive, et le soixantième qui est juste.
\end{enumerate}
\end{solution}

\begin{exercice}{Le théorème de Ptolémée et la formule d'addition}{Lycée}{2}
\manuscrit{latin-17859}{33}{34}
Les propositions VI à VIII du second livre démontrent le théorème de Ptolémée,
puis, dans un cercle bien choisi, une identité entre sinus. On refait le
chemin.

\exfigure{ptolemee}{Le quadrilatère inscrit $ABCD$ et le point $E$ de la question 1 :
$\widehat{ABE} = \widehat{DBC}$.}

\begin{enumerate}
\item Soit $ABCD$ un quadrilatère convexe inscrit dans un cercle. On place sur
  la diagonale $AC$ le point $E$ tel que $\widehat{ABE} = \widehat{DBC}$.
  Montrer que les triangles $ABE$ et $DBC$ sont semblables, puis que $CBE$ et
  $DBA$ le sont, et en déduire le théorème de Ptolémée :
  $AC \cdot BD = AB \cdot CD + AD \cdot BC$.
\item Montrer que, dans un cercle de diamètre $1$, une corde vue d'un point du
  cercle sous l'angle $\theta$ a pour longueur $\sin\theta$.
\item Soit $BD$ un diamètre d'un cercle de diamètre $1$, et $p$, $q$ deux angles
  aigus. On place $A$ sur le cercle d'un côté de $BD$ avec
  $\widehat{ADB} = p$, et $C$ de l'autre côté avec $\widehat{CDB} = q$. Calculer
  $AB$, $AD$, $CB$, $CD$ et $AC$, et déduire du théorème de Ptolémée la formule
  $\sin(p+q) = \sin p\cos q + \cos p \sin q$.
\item Obtenir de même $\sin(p-q)$ pour $p > q$, puis l'énoncé du feuillet :
  \[
  \frac{\sin p\cos q + \cos p\sin q}{\sin p\cos q - \cos p\sin q}
  = \frac{\sin(p+q)}{\sin(p-q)} .
  \]
\end{enumerate}
\end{exercice}

\begin{solution}
\begin{enumerate}
\item Les angles inscrits $\widehat{BAC}$ et $\widehat{BDC}$ interceptent le même
  arc $BC$, donc $\widehat{BAE} = \widehat{BDC}$. Avec
  $\widehat{ABE} = \widehat{DBC}$, les triangles $ABE$ et $DBC$ sont semblables,
  d'où $\frac{AE}{DC} = \frac{AB}{DB}$, soit $AE \cdot BD = AB \cdot CD$. Ensuite
  $\widehat{EBC} = \widehat{ABC} - \widehat{ABE} = \widehat{ABC} - \widehat{DBC}
  = \widehat{ABD}$, et $\widehat{BCE} = \widehat{BCA} = \widehat{BDA}$ (même arc
  $AB$) : les triangles $CBE$ et $DBA$ sont semblables, d'où
  $\frac{CE}{DA} = \frac{BC}{BD}$, soit $CE \cdot BD = AD \cdot BC$. En
  additionnant, $(AE + EC)\,BD = AB \cdot CD + AD \cdot BC$, et $AE + EC = AC$
  puisque $E$ est sur le segment $AC$ (l'angle $\widehat{ABE}$ est plus petit que
  $\widehat{ABC}$).
\item Soit une corde $MN$ vue du point $P$ du cercle sous l'angle $\theta$. Le
  diamètre issu de $M$ recoupe le cercle en $M'$ ; l'angle inscrit
  $\widehat{MM'N}$ intercepte le même arc que $\widehat{MPN}$, il vaut donc
  $\theta$ ou $180^\circ - \theta$, et le triangle $MNM'$ est rectangle en $N$.
  D'où $MN = MM' \sin\widehat{MM'N} = 1 \cdot \sin\theta$.
\item Le triangle $BAD$ est rectangle en $A$ (angle inscrit dans un demi-cercle)
  et d'hypoténuse $BD = 1$ : $AB = \sin p$, $AD = \cos p$. De même
  $CB = \sin q$, $CD = \cos q$. La corde $AC$ est vue de $D$ sous l'angle
  $\widehat{ADC} = \widehat{ADB} + \widehat{BDC} = p + q$ (les points $A$ et $C$
  sont de part et d'autre de $DB$), donc $AC = \sin(p+q)$. Le quadrilatère
  $BADC$ est inscrit, de diagonales $BD$ et $AC$ ; Ptolémée donne
  $BD \cdot AC = BA \cdot DC + AD \cdot CB$, soit
  $\sin(p+q) = \sin p \cos q + \cos p \sin q$.
\item Plaçons cette fois $C'$ du même côté que $A$, avec $\widehat{C'DB} = q < p$.
  Sur l'arc de $B$ à $D$, l'ordre est $B$, $C'$, $A$, $D$, et le quadrilatère
  $BC'AD$ a pour diagonales $BA$ et $C'D$. La corde $C'A$ est vue de $D$ sous
  l'angle $p - q$. Ptolémée donne $BA \cdot C'D = BC' \cdot AD + C'A \cdot DB$,
  soit $\sin p \cos q = \sin q \cos p + \sin(p-q)$, c'est-à-dire
  $\sin(p-q) = \sin p\cos q - \cos p \sin q$, quantité positive puisque $p > q$.
  En divisant les deux formules, on obtient l'énoncé du feuillet. La figure donne
  donc séparément les deux formules d'addition ; le feuillet, qui n'en voulait
  que le rapport, n'en tire que la proportion.
\end{enumerate}
\end{solution}

\begin{exercice}{Le « très noble théorème »}{L1}{2}
\manuscrit{latin-17859}{35}{37}
Le second livre démontre trois fois qu'un triangle sphérique $ABC$ rectangle en
$C$, dont l'angle $A$ est aigu, vérifie
\[
\frac{\sin(c+b)}{\sin(c-b)} = \frac{1 + \cos A}{1 - \cos A},
\]
où $a = BC$, $b = CA$, $c = AB$ sont les côtés (arcs de grands cercles de la
sphère unité, tous inférieurs à $90^\circ$). On le démontre ici avec des vecteurs.
On place $C = (0,0,1)$, $A = (\sin b, 0, \cos b)$, $B = (0, \sin a, \cos a)$.

\exfigure{triangle-spherique}{Le triangle sphérique $ABC$, rectangle en $C$, placé comme dans l'énoncé : $C$ au pôle, $A$ dans le plan $y = 0$, $B$ dans le plan $x = 0$.}

\begin{enumerate}
\item Vérifier que l'angle en $C$ est droit et que $\cos c = \cos a \cos b$.
\item L'angle $A$ est l'angle entre les vecteurs tangents en $A$ aux arcs $AC$ et
  $AB$, c'est-à-dire entre $C - (A\cdot C)A$ et $B - (A\cdot B)A$. Montrer que
  $\cos A = \frac{\cos a \sin b}{\sin c}$, puis que $\tan b = \tan c \cos A$.
\item En posant $g = \sin c \cos b$ et $h = \cos c \sin b$, montrer que
  $h = g \cos A$, puis démontrer l'énoncé du feuillet. Pourquoi a-t-on $c > b$ ?
\item En déduire la forme que lui donne la proposition XVI :
  $\frac{\sin(c+b)}{\sin(c-b)} = \cot^2\frac{A}{2}$.
\end{enumerate}
\end{exercice}

\begin{solution}
\begin{enumerate}
\item Les arcs $CA$ et $CB$ sont dans les plans $y = 0$ et $x = 0$, qui sont
  perpendiculaires : l'angle en $C$ est droit. Les arcs mesurent bien $b$ et $a$
  puisque $C \cdot A = \cos b$ et $C \cdot B = \cos a$. Enfin
  $\cos c = A \cdot B = \cos a \cos b$.
\item $C - (\cos b)A = (-\sin b\cos b, 0, 1 - \cos^2 b) = \sin b\,(-\cos b, 0,
  \sin b)$, de norme $\sin b$. $B - (\cos c)A = (-\cos c \sin b, \sin a,
  \cos a - \cos c \cos b)$ ; ce vecteur est orthogonal à $A$ et sa norme vaut
  $\sqrt{1 - \cos^2 c} = \sin c$ (projection de $B$ sur le plan tangent en $A$).
  Le produit scalaire de $(-\cos b, 0, \sin b)$ et de ce vecteur vaut
  $\cos b\cos c\sin b + \sin b\cos a - \sin b\cos c\cos b = \sin b\cos a$, donc
  \[
  \cos A = \frac{\cos a \sin b}{\sin c}.
  \]
  Avec $\cos a = \cos c / \cos b$, il vient
  $\cos A = \frac{\cos c \sin b}{\cos b \sin c} = \frac{\tan b}{\tan c}$.
\item $\tan b = \tan c \cos A$ s'écrit $\sin b \cos c = \cos A \sin c \cos b$,
  c'est-à-dire $h = g \cos A$. Alors
  $\sin(c+b) = \sin c\cos b + \cos c \sin b = g(1 + \cos A)$ et
  $\sin(c-b) = g - h = g(1 - \cos A)$. Comme $0 < c < 90^\circ$ et
  $0 < b < 90^\circ$, $g > 0$ ; comme $A$ est aigu, $1 - \cos A > 0$ ; le quotient
  donne l'énoncé. Enfin $\cos c = \cos a \cos b < \cos b$ puisque $0 < a < 90^\circ$,
  donc $c > b$ : l'hypoténuse est plus grande que le côté, ce qui rend
  $\sin(c-b)$ positif.
\item $1 + \cos A = 2\cos^2\frac{A}{2}$ et $1 - \cos A = 2\sin^2\frac{A}{2}$, d'où
  $\frac{1+\cos A}{1-\cos A} = \cot^2\frac{A}{2}$. C'est exactement la voie de
  Maurolico à la proposition X : deux « rectangles » $g$ et $h$ dans le rapport
  $1 : \cos A$, une proportion par composition et division (sa proposition IX),
  et l'identification de $g \pm h$ aux sinus de la somme et de la différence (sa
  proposition VIII). Maurolico attribue l'énoncé à Ménélas, dont ce serait la
  cinquième proposition du troisième livre des \emph{Sphériques} ; il le démontre
  « autrement et encore autrement », trois fois (propositions X, XII et XIV).
\end{enumerate}
\end{solution}

\begin{exercice}{Les rencontres de deux mobiles sur un cercle}{L2}{2}
\manuscrit{latin-17859}{9}{10}
Le \emph{Discours sur la sphère} affirme que deux mobiles tournant uniformément
sur un cercle, dont les vitesses sont « comme les plus petits nombres » $p$ et
$q$ ($p > q$, premiers entre eux), se rencontrent en $p - q$ points également
espacés s'ils vont dans le même sens, en $p + q$ s'ils vont en sens contraires,
et toujours aux mêmes points ; et que, si les vitesses sont incommensurables,
les points de rencontre ne se répètent jamais. On repère les positions par des
angles modulo $2\pi$, les mobiles partant ensemble de l'angle $0$ avec les
vitesses angulaires $\omega_1 > \omega_2 > 0$, $\omega_1/\omega_2 = p/q$.
\begin{enumerate}
\item Même sens : montrer que les rencontres ont lieu aux instants
  $t_k = \frac{2\pi k}{\omega_1 - \omega_2}$ ($k \in \mathbf{N}$), aux angles
  $\theta_k = \frac{2\pi k p}{p - q}$, et que l'ensemble des points de rencontre
  est formé de $p - q$ points également espacés.
\item Sens contraires : montrer qu'il y a $p + q$ points.
\item Même sens, $\omega_1/\omega_2$ irrationnel : montrer que les $\theta_k$
  sont deux à deux distincts modulo $2\pi$, puis qu'ils sont denses sur le cercle
  (ce que le feuillet ne dit pas).
\end{enumerate}
\end{exercice}

\begin{solution}
\begin{enumerate}
\item Les positions sont $\omega_1 t$ et $\omega_2 t$ ; il y a rencontre quand
  $(\omega_1 - \omega_2)t \in 2\pi\mathbf{Z}$, d'où les $t_k$, et
  $\theta_k = \omega_1 t_k = \frac{2\pi k\,\omega_1}{\omega_1 - \omega_2}
  = \frac{2\pi k p}{p-q}$. Posons $m = p - q \ge 1$. Comme
  $\operatorname{pgcd}(p, m) = \operatorname{pgcd}(p, q) = 1$, la classe de $p$ est
  inversible modulo $m$, et $k \mapsto kp \bmod m$ parcourt tous les restes
  $0, \dots, m-1$ lorsque $k$ parcourt $0, \dots, m-1$, puis recommence. Les
  points de rencontre sont donc exactement les angles $2\pi j/m$,
  $0 \le j < m$, visités périodiquement : $p - q$ points équidistants, toujours
  les mêmes.
\item La vitesse relative est $\omega_1 + \omega_2$, les rencontres ont lieu en
  $\theta_k = \frac{2\pi k p}{p+q}$, et $\operatorname{pgcd}(p, p+q) = 1$ ; le même
  argument donne $p + q$ points.
\item Posons $\beta = \frac{\omega_1}{\omega_1-\omega_2} = \frac{r}{r-1}$ avec
  $r = \omega_1/\omega_2$ ; $\beta$ est irrationnel si et seulement si $r$ l'est.
  Si $\theta_k \equiv \theta_l$, alors $(k - l)\beta \in \mathbf{Z}$, d'où $k = l$.
  Densité : soit $N \ge 1$. Parmi les $N+1$ nombres $\{k\beta\}$ (partie
  fractionnaire), $0 \le k \le N$, deux tombent dans un même intervalle
  $[\frac{j}{N}, \frac{j+1}{N}[$ : il existe $0 \le l < k \le N$ avec
  $0 < |\{k\beta\} - \{l\beta\}| < \frac1N$ (non nul par ce qui précède). Avec
  $m = k - l \ge 1$, on a donc $\{m\beta\} = \eta$ avec $0 < \eta < \frac1N$ ou
  $1 - \frac1N < \eta < 1$. Les nombres $\{jm\beta\}$, $j = 1, 2, \dots$, avancent
  alors sur le cercle $\mathbf{R}/\mathbf{Z}$ par pas constants de longueur
  $\min(\eta, 1 - \eta) < \frac1N$, dans un sens fixé : ils rencontrent tout arc
  de longueur $\frac1N$. Comme ce sont des rencontres $\theta_{jm}$, tout arc de
  longueur $2\pi/N$ contient un point de rencontre, pour tout $N$ : l'ensemble est
  dense. La conclusion du feuillet --- si les mouvements célestes sont
  incommensurables, le ciel ne revient jamais à la même disposition --- est
  juste pour deux mobiles ; il reste que la disposition revient aussi près qu'on
  veut.
\end{enumerate}

\exfigure{mobiles}{Les points de rencontre pour $p = 5$ et $q = 2$, numérotés dans l'ordre où ils sont atteints : $p - q = 3$ dans le même sens, $p + q = 7$ en sens contraires.}
\end{solution}

\section*{François Viète (1540--1603)}
\chapitre{viete}

NAL 1644 réunit deux traités latins que leurs titres, sur les feuillets mêmes,
donnent à Viète : « Francisci Vietæ De Recognitione Æquationum tractatus »
(vues 5 à 68) et « Francisci Vietæ. De æquationum Emendatione tractatus
secundus » (vues 69 à 141), suivis de pages d'astronomie qui ne portent aucun
nom. Les mains n'y sont pas identifiées, et deux notes douteuses donnent à
penser que les traités sont une copie. On y écrit en « espèces » : l'inconnue
est une voyelle, les grandeurs données des consonnes qui portent leur
dimension, et l'égalité se dit en toutes lettres. Les exercices traduisent
cette notation, puis refont sur les exemples des feuillets ce que les traités
appellent reconnaître et corriger une équation : lire ses coefficients sur ses
racines, résoudre la cubique dans le cas irréductible, abaisser la quartique à
deux équations du second degré.

\begin{exercice}{Des racines aux coefficients : la syncrise}{Lycée}{2}
\manuscrit{nal-1644}{48}{50}
\manuscrit{nal-1644}{139}{141}
Le chapitre XVI du traité titré « De Recognitione Æquationum » appelle
\emph{syncrise} la comparaison de deux équations de même forme et de mêmes
données qui ont deux racines différentes. Le traité n'admet que des racines
positives, et il écrit chaque équation avec des termes positifs de part et
d'autre : $Bx - x^2 = Z$ n'est pas pour lui la même espèce que $x^2 - Bx = Z$.
Dans les exemples chiffrés, $N$ est l'inconnue et $Q$, $C$, $QQ$, $QC$ ses
puissances deuxième à cinquième : « $124N - 1C$ æquabitur $240$ » est
$124x - x^3 = 240$. Le dernier chapitre du second traité, « De æquationum
Emendatione », donne enfin sans démonstration des équations « de multiplicibus
radicibus mire explicabiles », admirablement résolubles par plusieurs racines.
\begin{enumerate}
\item Soient $x > y > 0$ deux racines de $Bt - t^2 = Z$. En égalant $Bx - x^2$
  et $By - y^2$, puis en divisant par $x - y$, comme le fait le feuillet, montrer
  que $B = x + y$ et $Z = xy$. Écrire l'équation de cette forme dont les racines
  sont $10$ et $2$.
\item Même question pour $Bt - t^3 = Z$ : montrer que $B = x^2 + xy + y^2$ et
  $Z = xy(x + y)$, et retrouver l'exemple du feuillet, « $124N - 1C$ æquabitur
  $240$ », de racines $10$ et $2$. Développer
  $\bigl(t^2 - (x+y)t + xy\bigr)\bigl(t + x + y\bigr)$ et en déduire la troisième
  racine de $t^3 - Bt + Z = 0$. Que vaut-elle dans l'exemple ?
\item La proposition III du dernier chapitre dit que, si $B$, $D$, $G$, $H$ sont
  quatre grandeurs, l'équation qu'elle écrit a pour racines « l'une quelconque des
  quatre ». Son exemple est « $50N - 35Q + 16C - 1QQ$ æquatur $24$, fit $1N$ $1$.
  $2$. $3$. vel $4$ ». Développer $(x-1)(x-2)(x-3)(x-4)$, mettre l'équation sous
  la forme du feuillet, et trouver la faute de l'exemple. Le nombre $1$ est-il
  racine de l'équation telle qu'elle est écrite ?
\item L'exemple de la proposition IIII est « $1QC - 15QQ + 85C - 225Q + 274N$
  æquatur $120$, fit $1N$ $1$. $2$. $3$. $4$. vel $5$ ». En posant $w = x^2 - 6x$,
  écrire $(x-1)(x-5)$ et $(x-2)(x-4)$ en fonction de $w$, puis développer
  $(x-1)(x-2)(x-3)(x-4)(x-5)$ et vérifier l'exemple. Expliquer enfin pourquoi
  $274 = 120\bigl(1 + \frac12 + \frac13 + \frac14 + \frac15\bigr)$.
\end{enumerate}
\end{exercice}

\begin{solution}
\begin{enumerate}
\item De $Bx - x^2 = Z = By - y^2$ on tire $x^2 - y^2 = B(x - y)$, et, comme
  $x - y \neq 0$, $B = x + y$. Alors $Z = Bx - x^2 = (x + y)x - x^2 = xy$ : le
  coefficient est la somme des deux racines, le terme connu leur produit. Pour les
  racines $10$ et $2$ : $12x - x^2 = 20$, la « formula æquationis » du feuillet, « $12N - 1Q$
  æquabit $20$ explicabilis de $10$ vel $2$ ».
\item De $Bx - x^3 = By - y^3$ on tire $x^3 - y^3 = B(x - y)$, et, en divisant par
  $x - y$, $B = x^2 + xy + y^2$. Alors
  $Z = Bx - x^3 = x^3 + x^2y + xy^2 - x^3 = xy(x + y)$. Pour $10$ et $2$ :
  $B = 100 + 20 + 4 = 124$ et $Z = 20 \times 12 = 240$ ; c'est l'exemple du feuillet.
  À cet endroit la page écrit « $D$ plano » pour le coefficient $B$, ce que la
  marge corrige par des « $B$ », et « omnibus per $A$ divisis » alors que le
  diviseur écrit dessous est $A - E$ : deux lapsus, sans effet sur le résultat.
  Ensuite
  \[
  \bigl(t^2 - (x+y)t + xy\bigr)(t + x + y)
  = t^3 - \bigl((x+y)^2 - xy\bigr)t + xy(x+y) = t^3 - Bt + Z ,
  \]
  les termes en $t^2$ se détruisant. La troisième racine est donc $-(x + y)$,
  toujours négative ; dans l'exemple,
  $t^3 - 124t + 240 = (t - 2)(t - 10)(t + 12)$ et elle vaut $-12$. Le traité, qui
  ne connaît que les racines positives, ne la nomme pas.
\item $(x-1)(x-2) = x^2 - 3x + 2$ et $(x-3)(x-4) = x^2 - 7x + 12$, d'où
  \[
  (x-1)(x-2)(x-3)(x-4) = x^4 - 10x^3 + 35x^2 - 50x + 24 .
  \]
  Les racines sont $1$, $2$, $3$, $4$ si et seulement si
  $50x - 35x^2 + 10x^3 - x^4 = 24$ : $10$ est la somme des racines, $35$ la somme de
  leurs produits deux à deux, $50 = 6 + 8 + 12 + 24$ celle de leurs produits trois à
  trois et $24$ leur produit, comme le veut la proposition. Le feuillet écrit $16$
  au lieu de $10$ pour le coefficient de $x^3$ ; la marge le corrige en « $10$ ».
  Telle qu'elle est écrite, l'équation donne en $x = 1$ la valeur
  $50 - 35 + 16 - 1 = 30 \neq 24$ : $1$ n'est pas racine (ni $2$, $3$, $4$, qui
  donnent $72$, $186$, $408$). L'énoncé de la proposition, sur la même page, omet
  aussi le terme $-x^4$, que la marge ajoute (« $-$ Aqq. »).
\item $(x-1)(x-5) = x^2 - 6x + 5 = w + 5$ et $(x-2)(x-4) = x^2 - 6x + 8 = w + 8$,
  d'où
  \[
  (w+5)(w+8) = w^2 + 13w + 40 = x^4 - 12x^3 + 49x^2 - 78x + 40 ,
  \]
  puis, en multipliant par $x - 3$,
  \[
  (x-1)(x-2)(x-3)(x-4)(x-5) = x^5 - 15x^4 + 85x^3 - 225x^2 + 274x - 120 .
  \]
  L'équation $x^5 - 15x^4 + 85x^3 - 225x^2 + 274x = 120$ a donc bien pour racines
  $1$, $2$, $3$, $4$, $5$ : l'exemple est juste. Le coefficient de $x$ est la somme
  des produits de quatre des cinq racines ; chacun est le produit des cinq,
  $120$, divisé par la racine qui manque, d'où
  $274 = \frac{120}{1} + \frac{120}{2} + \frac{120}{3} + \frac{120}{4} + \frac{120}{5}
  = 120 + 60 + 40 + 30 + 24$.
\end{enumerate}
\end{solution}

\begin{exercice}{Le cas irréductible par deux triangles rectangles}{L1}{2}
\manuscrit{nal-1644}{16}{17}
Le chapitre VI du « De Recognitione Æquationum » traite la cubique sans terme en
$x^2$. Son théorème III (vue 16) suppose $x^3 - 3Bx = D$ avec $D^2 < 4B^3$, et
appelle « inverse » l'équation $3By - y^3 = D$. Vient ensuite, sous le titre
« Aliter, tertium theorema » (vue 17), une constitution « plus élégante », tirée
des sections angulaires : si $x^3 - 3B^2x = B^2D$ avec $B > D/2$, on prend deux
triangles rectangles de même hypoténuse $B$, tels que l'angle aigu opposé à la
hauteur (« perpendiculum ») du premier soit triple de l'angle correspondant du
second, et que le double de la base du premier soit $D$ ; alors $x$ est le double
de la base du second, et les racines de l'équation inverse
$3B^2y - y^3 = B^2D$ sont la base du second augmentée ou diminuée de la droite
dont le carré est le triple du carré de sa hauteur. On note $\theta$ l'angle aigu
du second triangle opposé à sa hauteur.

\exfigure{viete-triangles}{Les deux triangles de l'« Aliter », de même hypoténuse $B$ ; l'angle opposé à la hauteur vaut $3\theta$ dans le premier, $\theta$ dans le second. Tracé pour l'exemple de la page, $B = 10$ et $\cos 3\theta = 0{,}216$.}

\begin{enumerate}
\item Exprimer les bases et les hauteurs des deux triangles à l'aide de $B$ et de
  $\theta$. Montrer que, si $x = 2B\cos\theta$, alors
  $x^3 - 3B^2x = 2B^3\cos 3\theta$, et en déduire la première affirmation du
  feuillet. À quoi sert l'hypothèse $B > D/2$ ?
\item L'exemple de la page est $x^3 - 300x = 432$. Trouver $B$, $D$ et la base
  du premier triangle ; la page écrit « basis autem primi dupla est
  $\frac{432}{100}$ », puis « fit basis $2\frac{16}{100}$ ». Vérifier que le second
  triangle a pour base $9$ et pour hauteur $\sqrt{19}$, comme le dit la page, et
  conclure que $x = 18$.
\item Montrer que les deux autres racines de $x^3 - 3B^2x = B^2D$ sont
  $2B\cos(\theta \pm 120^\circ)$, et qu'elles sont négatives. En déduire la seconde
  affirmation du feuillet, et donner les racines de $300y - y^3 = 432$.
\item Montrer que la condition $D^2 < 4B^3$ du théorème III, écrite pour
  $x^3 - 3Bx = D$, et la condition $B > D/2$ de l'« Aliter », écrite pour
  $x^3 - 3B^2x = B^2D$, sont la même.
\item La formule qu'on appelle aujourd'hui de Cardan donne, pour $x^3 - px = q$,
  $x = \sqrt[3]{q/2 + \sqrt{\Delta}} + \sqrt[3]{q/2 - \sqrt{\Delta}}$ avec
  $\Delta = \frac{q^2}{4} - \frac{p^3}{27}$. Calculer $\Delta$ pour
  $x^3 - 300x = 432$. Calculer $(9 + i\sqrt{19})^3$, et expliquer comment la
  formule, lue dans $\mathbf{C}$, redonne $x = 18$. Que sont, sur la figure du
  feuillet, la partie réelle et la partie imaginaire de $9 + i\sqrt{19}$ ?
\end{enumerate}
\end{exercice}

\begin{solution}
\begin{enumerate}
\item L'angle opposé à la hauteur est adjacent à la base : le second triangle a
  pour base $B\cos\theta$ et pour hauteur $B\sin\theta$, le premier pour base
  $B\cos 3\theta$ et pour hauteur $B\sin 3\theta$ (avec $3\theta < 90^\circ$, puisque
  $3\theta$ est un angle aigu). De $\cos 3\theta = \cos(2\theta + \theta)
  = 4\cos^3\theta - 3\cos\theta$ on tire, pour $x = 2B\cos\theta$,
  \[
  x^3 - 3B^2x = 8B^3\cos^3\theta - 6B^3\cos\theta = 2B^3\cos 3\theta .
  \]
  Si le double de la base du premier triangle vaut $D$, c'est-à-dire
  $B\cos 3\theta = D/2$, il vient $x^3 - 3B^2x = 2B^2 \cdot \frac{D}{2} = B^2D$ :
  le double de la base du second est racine. L'hypothèse $B > D/2$ dit que la base
  $D/2$ est plus petite que l'hypoténuse $B$ : le premier triangle existe, avec
  $3\theta = \arccos\frac{D}{2B} \in {]0^\circ, 90^\circ[}$, et le second
  s'obtient en prenant le tiers de son angle, ce qui est la trisection que le
  feuillet renvoie aux « sections angulaires ».
\item $3B^2 = 300$ donne $B = 10$, et $B^2D = 432$ donne $D = 4{,}32$ : le double
  de la base du premier triangle est $\frac{432}{100}$, sa base est $2{,}16$, comme
  l'écrit la page, et $\cos 3\theta = 0{,}216$. Avec $\cos\theta = 0{,}9$ on trouve
  $4 \times 0{,}729 - 3 \times 0{,}9 = 2{,}916 - 2{,}7 = 0{,}216$ ; comme
  $\theta \mapsto \cos 3\theta$ est injective sur ${]0^\circ, 30^\circ[}$,
  c'est bien l'angle cherché ($\theta \approx 25{,}84^\circ$,
  $3\theta \approx 77{,}53^\circ$). Le second triangle a pour base $10 \times 0{,}9 = 9$
  et pour hauteur $\sqrt{100 - 81} = \sqrt{19}$. D'où $x = 18$, et en effet
  $18^3 - 300 \times 18 = 5832 - 5400 = 432$.
\item $\cos\bigl(3(\theta \pm 120^\circ)\bigr) = \cos(3\theta \pm 360^\circ)
  = \cos 3\theta$ : les nombres $2B\cos(\theta \pm 120^\circ)$ sont racines, et,
  comme $\theta$, $\theta + 120^\circ$ et $\theta - 120^\circ$ ont des cosinus
  distincts pour $0^\circ < \theta < 30^\circ$, ce sont avec $2B\cos\theta$ les trois
  racines. Or
  \[
  2B\cos(\theta \pm 120^\circ) = -B\cos\theta \mp \sqrt3\,B\sin\theta ,
  \]
  et $\theta < 30^\circ$ donne $\tan\theta < \frac{1}{\sqrt3}$, soit
  $B\cos\theta > \sqrt3\,B\sin\theta$ : les deux racines sont négatives. Si $x$ est
  racine de $x^3 - 3B^2x = B^2D$, alors $y = -x$ vérifie
  $3B^2y - y^3 = x^3 - 3B^2x = B^2D$. Les racines de l'équation inverse sont donc
  $B\cos\theta \pm \sqrt3\,B\sin\theta$, toutes deux positives : la base du second
  triangle plus ou moins $\sqrt{3\,(B\sin\theta)^2}$, la droite dont le carré est
  le triple du carré de la hauteur. Dans l'exemple, $9 \pm \sqrt{57}$, soit environ
  $16{,}550$ et $1{,}450$, comme l'écrit la page (« $9 + L\,57$ vel $9 - L\,57$ ») ;
  et $x^3 - 300x - 432 = (x - 18)(x^2 + 18x + 24)$.
\item L'équation $x^3 - 3B^2x = B^2D$ est l'équation $x^3 - 3B'x = D'$ du théorème
  III avec $B' = B^2$ et $D' = B^2D$. La condition $D'^2 < 4B'^3$ s'écrit
  $B^4D^2 < 4B^6$, soit $D^2 < 4B^2$, soit $D < 2B$ puisque $B$ et $D$ sont
  positifs. Dans l'exemple, $432^2 = 186\,624 < 4 \times 100^3$.
\item Avec $p = 300$ et $q = 432$ : $\frac{q^2}{4} = 46\,656$,
  $\frac{p^3}{27} = 1\,000\,000$, $\Delta = -953\,344 < 0$, et, comme
  $224^2 \times 19 = 953\,344$, $\sqrt{\Delta} = 224\,i\sqrt{19}$ : la formule
  demande la racine carrée d'un nombre négatif, bien que les trois racines soient
  réelles. Or
  \[
  (9 + i\sqrt{19})^3 = 729 + 243\,i\sqrt{19} - 513 - 19\,i\sqrt{19}
  = 216 + 224\,i\sqrt{19} .
  \]
  Les racines cubiques $u = 9 + i\sqrt{19}$ de $216 + 224\,i\sqrt{19}$ et
  $v = 9 - i\sqrt{19}$ de son conjugué vérifient $uv = 81 + 19 = 100 = p/3$, la
  condition qui les apparie dans la formule, et $u + v = 18$. Le nombre
  $9 + i\sqrt{19}$ est $10(\cos\theta + i\sin\theta)$ : sa partie réelle est la base
  du second triangle, sa partie imaginaire sa hauteur, et son cube,
  $1000(\cos 3\theta + i\sin 3\theta)$, a pour partie réelle $1000 \times 0{,}216 =
  216 = q/2$. Les deux triangles du feuillet sont exactement ce que la formule
  demande en nombres complexes ; la trisection de l'angle remplace l'extraction de
  la racine cubique d'un nombre complexe. Les deux autres appariements,
  $ju + \overline{j}v$ et $\overline{j}u + jv$ avec $j = \cos 120^\circ + i\sin 120^\circ$,
  donnent les racines négatives $-9 \mp \sqrt{57}$.
\end{enumerate}
\end{solution}

\begin{exercice}{La paraplérose climactique d'une quartique, et ses deux formules}{L2}{3}
\manuscrit{nal-1644}{109}{117}
Le chapitre VI du traité titré « De æquationum Emendatione tractatus secundus »
s'intitule « de climactica paraplerosi » : il abaisse une équation du quatrième
degré sans terme en $x^3$ à des équations du second degré « par le moyen d'une
cubique », en ajoutant aux deux membres de quoi en faire deux carrés. Le problème
III (vues 109 à 111) traite $x^4 + 2Gx^2 + Bx = Z$ : la « première formule »
forme le carré de $x^2 + G + \frac{y^2}{2}$, où $y$ est une grandeur auxiliaire, et
veut que l'autre membre soit le carré de $\frac{B}{2y} - yx$. La seconde
(« Aliter », vues 111 et 112), écrite pour $x^4 + Gx^2 + Bx = Z$, ajoute aux deux
membres de $x^4 = Z - Gx^2 - Bx$ la quantité $Fx^2 + \frac{F^2}{4}$. Chaque formule
est suivie de sept théorèmes, un par combinaison de signes, chacun avec un exemple
en signes cossiques ($N$, $Q$, $C$, $QQ$ pour $x$, $x^2$, $x^3$, $x^4$). Les
grandeurs $B$, $G$, $Z$ sont positives.
\begin{enumerate}
\item Montrer qu'en ajoutant aux deux membres de $x^4 + 2Gx^2 = Z - Bx$ ce qui
  manque au premier pour être le carré de $x^2 + G + \frac{y^2}{2}$, on obtient
  \[
  \Bigl(x^2 + G + \frac{y^2}{2}\Bigr)^2 = y^2x^2 - Bx + Z + G^2 + Gy^2 + \frac{y^4}{4} ,
  \]
  et que le second membre est, pour tout $x$, le carré de $\frac{B}{2y} - yx$ si et
  seulement si $u = y^2$ vérifie
  $u^3 + 4Gu^2 + (4Z + 4G^2)u = B^2$ (la « résolvante »). Montrer que cette
  équation a une racine positive et une seule.
\item Soit $D > 0$ tel que $D^2$ soit cette racine. En déduire deux équations du
  second degré dont les racines sont celles de la quartique ; l'une est celle du
  feuillet, $x^2 + Dx = \frac{B}{2D} - \frac{D^2}{2} - G$. Appliquer à l'exemple du
  théorème I (vue 113), « $1QQ + 6Q + 880N$ æquatur $1800$ » : écrire la
  résolvante, vérifier la valeur de la page, « $1C + 12Q + 7236N$ æquatur $774400$
  fit $1N$ $64$ », et trouver les quatre racines complexes de la quartique.
  Lesquelles la page donne-t-elle ?
\item On cherche aujourd'hui une factorisation
  $x^4 + 2Gx^2 + Bx - Z = (x^2 + Dx + p)(x^2 - Dx + q)$. Montrer que $D^2$ est
  alors racine de la même résolvante, et que $p = G + \frac{D^2}{2} - \frac{B}{2D}$.
  Si $r_1$, $r_2$, $r_3$, $r_4$ sont les racines de la quartique, montrer que les
  racines de la résolvante sont $(r_1 + r_2)^2$, $(r_1 + r_3)^2$ et $(r_1 + r_4)^2$.
  Les calculer pour l'exemple.
\item Montrer que la seconde formule conduit, pour $x^4 + Gx^2 + Bx = Z$, à la
  résolvante $F^3 - GF^2 + 4ZF = B^2 + 4ZG$ et à
  $x^2 + \sqrt{F - G}\,x = \frac{B}{2\sqrt{F - G}} - \frac{F}{2}$, et que $F - G$ est
  racine de la résolvante de la première formule, écrite pour la quartique
  $x^4 + 2G'x^2 + Bx = Z$ avec $G' = G/2$. Vérifier
  sur le théorème I « secundum posteriorem formulam » (vue 116), qui reprend la
  même quartique : la page trouve « $1C - 6Q + 7200N$ æquatur $817600$ fit $1N$
  $70$ » et « $70 - 6$ quadratum abs $8$ ».
\item Le théorème IIII de la seconde formule (vue 117), pour
  $x^4 - Gx^2 - Bx = Z$, écrit la résolvante $F^3 + GF^2 + 4ZF = B^2 + 4ZG$ et
  donne l'exemple « $1QQ - 4Q - 800N$ æquatur $1600$ fit $1N$ $10$ », « $1C + 4Q +
  64N$ æquatur $614400$ fit $1N$ $60$ », « $60 + 4$ est quadratum abs $8$ »,
  « $1Q - 8N$ æquatur $20$ fit $1N$ $10$ ». Refaire la construction et le calcul,
  et dire exactement où la page se trompe.
\end{enumerate}
\end{exercice}

\begin{solution}
\begin{enumerate}
\item $\bigl(x^2 + G + \frac{y^2}{2}\bigr)^2 = x^4 + 2Gx^2 + y^2x^2 + G^2 + Gy^2 +
  \frac{y^4}{4}$ : il faut ajouter $y^2x^2 + G^2 + Gy^2 + \frac{y^4}{4}$, ce qui donne
  l'égalité annoncée. Le second membre est un trinôme en $x$ de coefficient
  dominant $y^2$ et de coefficient de $x$ égal à $-B$ ; le carré
  $\bigl(\frac{B}{2y} - yx\bigr)^2 = y^2x^2 - Bx + \frac{B^2}{4y^2}$ a les mêmes, et
  les deux coïncident si et seulement si les termes constants sont égaux :
  $Z + G^2 + Gy^2 + \frac{y^4}{4} = \frac{B^2}{4y^2}$. En multipliant par $4y^2$ et
  en posant $u = y^2$ : $u^3 + 4Gu^2 + (4Z + 4G^2)u = B^2$. C'est l'équation en $E$
  de la page (vues 110 et 113), une fois reçu le « $+$ » que la marge rétablit
  devant « $G$ plano-plano $4$ ». Le polynôme
  $R(u) = u^3 + 4Gu^2 + (4Z + 4G^2)u - B^2$ vaut $-B^2 < 0$ en $0$, tend vers
  $+\infty$, et il est strictement croissant sur $[0, +\infty[$ puisque tous ses
  coefficients non constants sont positifs : il a une racine positive et une seule.
\item Avec $u = D^2$, $\bigl(x^2 + G + \frac{D^2}{2}\bigr)^2 =
  \bigl(Dx - \frac{B}{2D}\bigr)^2$, donc
  $x^2 + G + \frac{D^2}{2} = \pm\bigl(Dx - \frac{B}{2D}\bigr)$, soit
  \[
  x^2 + Dx = \frac{B}{2D} - \frac{D^2}{2} - G
  \qquad\text{ou}\qquad
  x^2 - Dx + \frac{B}{2D} + \frac{D^2}{2} + G = 0 .
  \]
  La première est celle du feuillet ; à la vue 113, le signe de « $D$ quadrato
  $\frac12$ » y est doublé de traits, mais l'exemple demande bien $-\frac{D^2}{2}$.
  Dans l'exemple, $2G = 6$, $G = 3$, $B = 880$, $Z = 1800$ ; la résolvante est
  $u^3 + 12u^2 + 7236u = 774\,400$, puisque $4 \times 1800 + 4 \times 9 = 7236$ et
  $880^2 = 774\,400$. Et $u = 64$ convient :
  $262\,144 + 49\,152 + 463\,104 = 774\,400$. Donc $D = 8$, puis
  $x^2 + 8x = 55 - 32 - 3 = 20$, de racines $2$ et $-10$ ; et
  $x^2 - 8x + 90 = 0$, de racines $4 \pm i\sqrt{74}$. En effet
  $x^4 + 6x^2 + 880x - 1800 = (x^2 + 8x - 20)(x^2 - 8x + 90)$. La page ne donne que
  la racine positive, « fit $1N$ $2$ » ; $-10$ et les deux racines non réelles ne
  sont pas de son domaine.
\item En développant,
  $(x^2 + Dx + p)(x^2 - Dx + q) = x^4 + (p + q - D^2)x^2 + D(q - p)x + pq$, d'où
  $p + q = 2G + D^2$, $D(q - p) = B$ (ce qui impose $D \neq 0$) et $pq = -Z$. Comme
  $(p + q)^2 - (q - p)^2 = 4pq$, il vient
  $(2G + D^2)^2 - \frac{B^2}{D^2} = -4Z$, et, en multipliant par $u = D^2$,
  $u(u + 2G)^2 + 4Zu = B^2$, qui est la résolvante. Enfin
  $p = \frac{(p + q) - (q - p)}{2} = G + \frac{D^2}{2} - \frac{B}{2D}$ : le facteur
  $x^2 + Dx + p$ est exactement l'équation du feuillet. La quartique n'a pas de terme
  en $x^3$, donc $r_1 + r_2 + r_3 + r_4 = 0$. En groupant $r_1, r_2$ d'un côté et
  $r_3, r_4$ de l'autre, $(x - r_1)(x - r_2)(x - r_3)(x - r_4)$ est une factorisation
  de la forme cherchée avec $D = -(r_1 + r_2) = r_3 + r_4$ : $(r_1 + r_2)^2$ est
  racine de la résolvante, et de même $(r_1 + r_3)^2$ et $(r_1 + r_4)^2$ pour les
  deux autres groupements. Dans l'exemple, avec $r_1 = 2$ :
  $(2 - 10)^2 = 64$ et $(2 + 4 \pm i\sqrt{74})^2 = (6 \pm i\sqrt{74})^2
  = -38 \pm 12\,i\sqrt{74}$. Ces trois nombres sont distincts, et une équation du
  troisième degré n'a pas plus de trois racines : ce sont les racines de la
  résolvante, $u^3 + 12u^2 + 7236u - 774\,400 = (u - 64)(u^2 + 76u + 12\,100)$. La
  racine positive que prend le feuillet est la seule réelle.
\item De $x^4 = Z - Gx^2 - Bx$ on tire
  $\bigl(x^2 + \frac{F}{2}\bigr)^2 = (F - G)x^2 - Bx + Z + \frac{F^2}{4}$. Pour
  $F > G$, le second membre est le carré de
  $\sqrt{F - G}\,x - \frac{B}{2\sqrt{F - G}}$ si et seulement si
  $Z + \frac{F^2}{4} = \frac{B^2}{4(F - G)}$, soit $(F - G)(F^2 + 4Z) = B^2$,
  c'est-à-dire $F^3 - GF^2 + 4ZF = B^2 + 4ZG$ ; et l'égalité des racines donne
  $x^2 + \frac{F}{2} = \frac{B}{2\sqrt{F - G}} - \sqrt{F - G}\,x$, l'équation
  annoncée. En posant $F = u + G$, la résolvante devient
  $u\bigl((u + G)^2 + 4Z\bigr) = B^2$, soit $u^3 + 2Gu^2 + (G^2 + 4Z)u = B^2$ : c'est
  celle de la première formule pour $G' = G/2$, puisque $4G' = 2G$ et
  $4G'^2 = G^2$. Les
  deux formules donnent donc le même $D = \sqrt{F - G}$ et la même équation du
  second degré, puisque $\frac{F}{2} = \frac{D^2}{2} + G'$. Dans l'exemple, $G = 6$,
  et $B^2 + 4ZG = 774\,400 + 43\,200 = 817\,600$ ; pour $F = 70$,
  $343\,000 - 29\,400 + 504\,000 = 817\,600$, et $70 - 6 = 64 = 8^2$, la racine $u$ de
  la question 2 ; enfin $x^2 + 8x = 55 - 35 = 20$ et $x = 2$.
\item De $x^4 = Z + Gx^2 + Bx$ on tire
  $\bigl(x^2 + \frac{F}{2}\bigr)^2 = (F + G)x^2 + Bx + Z + \frac{F^2}{4}$, carré de
  $\sqrt{F + G}\,x + \frac{B}{2\sqrt{F + G}}$ si et seulement si
  $(F + G)(F^2 + 4Z) = B^2$, soit
  \[
  F^3 + GF^2 + 4ZF = B^2 - 4ZG ,
  \]
  et non $B^2 + 4ZG$ ; puis $x^2 - \sqrt{F + G}\,x = \frac{B}{2\sqrt{F + G}} -
  \frac{F}{2}$. Pour $G = 4$, $B = 800$, $Z = 1600$ : $4Z = 6400$ et
  $B^2 - 4ZG = 640\,000 - 25\,600 = 614\,400$, qui est le second membre de la page
  (avec le signe de l'énoncé, on aurait $665\,600$). L'équation juste est
  $F^3 + 4F^2 + 6400F = 614\,400$, et $F = 60$ la vérifie :
  $216\,000 + 14\,400 + 384\,000 = 614\,400$. Avec le coefficient « $64$ » de la
  page, $F = 60$ donnerait $216\,000 + 14\,400 + 3840 = 234\,240$. Ensuite
  $\sqrt{60 + 4} = 8$, $x^2 - 8x = \frac{800}{16} - 30 = 20$, et $x = 10$ ; en effet
  $10^4 - 400 - 8000 = 1600$, et
  $x^4 - 4x^2 - 800x - 1600 = (x - 10)(x + 2)(x^2 + 8x + 80)$. La page se trompe
  donc deux fois, et seulement dans l'écriture : le signe du terme $4ZG$ dans
  l'énoncé de la résolvante, que la marge a corrigé au théorème III mais pas ici,
  et « $64N$ » pour « $6400N$ » dans l'exemple ; le nombre $614\,400$, la racine $60$
  et la racine $10$ sont justes.
\end{enumerate}
\end{solution}

\section*{Marin Mersenne (1588--1648)}
\chapitre{mersenne}

Quatre volumes nourrissent ce chapitre. Français 12357 est la copie, par une main
anonyme du XVII\textsuperscript{e} siècle, des notes que le P. Mersenne avait
écrites dans les marges et sur les feuillets blancs de son exemplaire de
l'\emph{Harmonie universelle} ; la copie le dit dans son titre, et c'est sur ce
titre que repose l'attribution des remarques. Chaque remarque renvoie à une page
du livre : la nature des sons, le mouvement des corps, les chants, les
consonances, puis huit pages de nombres. Beaucoup de ses calculs sont justes,
certains remarquablement ; d'autres sont faux, et les exercices les prennent des
deux côtés. NAF 9544, un album d'autographes, conserve une feuille d'observations
astronomiques de 1620--1621 adressée au P. Petau ; elle n'est pas signée, et
c'est la coupure du catalogue de vente et la lettre signée « Mar. Mers. M. » qui
l'accompagnait qui la donnent à Mersenne. NAF 5176, enfin, est une reliure de
papiers épars, en quatre mains au moins, que le catalogue range sous le nom de
Mersenne : une question latine sur les triangles rectangles en nombres, puis,
après un feuillet moderne de classement, deux mains rondes qui traitent des
« partis » (le partage d'une mise interrompue) et des dés. Ces feuillets de jeux
ne nomment personne : les rapprochements avec la lettre de Pascal à Fermat du
29 juillet 1654 et avec le traité de Huygens (\emph{De ratiociniis in ludo
aleae}, 1657) sont des coïncidences d'énoncés, non des attributions.
NAL 2338, enfin, les « Papiers divers de Mersenne », est le dossier de la
querelle de la cycloïde entre Torricelli et Roberval, passée par Mersenne de
1643 à 1647 ; ce sont presque toutes des copies. Les deux exercices qui en
viennent ici reposent sur des lettres que leur intitulé ou leur souscription
donnent à Torricelli ; ceux qui reposent sur les pièces de Roberval sont au
chapitre suivant. Aucune question de priorité n'y est tranchée.

\begin{exercice}{Trois valeurs de $\pi$ et un polygone de 97 côtés}{Lycée}{1}
\manuscrit{fr-12357}{8}{9}
Les feuillets blancs de devant discutent la quadrature du cercle d'après un livre
de Fabre, et opposent trois règles.

\exfigure{fabre}{La règle de Fabre (question 1) : la droite $AM$, qui joint un sommet du carré au milieu d'un côté opposé, serait égale au quart de cercle $AB$.}

\begin{enumerate}
\item Fabre fait égale au quart de la circonférence circonscrite à un carré la
  droite qui joint un sommet du carré au milieu d'un côté opposé. Quelle valeur
  de $\pi$ cette règle suppose-t-elle ?
\item Le sieur Cornu fait que, le côté du carré inscrit étant $9$, le quart de la
  circonférence est $10$. Quelle valeur de $\pi$ ? La remarque ajoute que Cornu
  fait le quart de la circonférence « trop grand d'un peu, car le polygone
  circonscrit de 97 côtés est plus petit que son cercle supposé ». Montrer que le
  périmètre du polygone régulier à $n$ côtés circonscrit à un cercle de diamètre
  $1$ vaut $n\tan\frac{\pi}{n}$ et qu'il dépasse $\pi$ ; puis, en calculant
  $97\tan\frac{\pi}{97}$, justifier l'argument de la remarque.
\item Un troisième auteur donne à la circonférence « trois fois le diamètre et
  $\frac{112}{113}$ d'une septième partie ». Quelle fraction est-ce, et quelle
  est son erreur ?
\item « Par le moyen de 70 et 99 multipliés en soi-mêmes l'on a deux nombres en
  raison double, à $\frac{1}{4900}$ près. » Expliquer.
\item « Pour faire un rond égal à un carré », on donne $8$ au diamètre du cercle et
  $10$ à la diagonale du carré. Avec $\pi = \frac{22}{7}$, la copie trouve pour le
  cercle $50\frac47$ et un excès d'$\frac{1}{88}$ sur le carré. Refaire le calcul.
\end{enumerate}
\end{exercice}

\begin{solution}
\begin{enumerate}
\item Pour un carré de côté $a$, la droite du sommet au milieu du côté opposé vaut
  $\sqrt{a^2 + a^2/4} = \frac{\sqrt5}{2}a$ ; le cercle circonscrit a pour rayon
  $\frac{a}{\sqrt2}$ et son quart de circonférence vaut $\frac{\pi a}{2\sqrt2}$.
  L'égalité donne $\pi = \sqrt2\cdot\sqrt5 = \sqrt{10} \approx 3{,}1623$, trop
  grand de $0{,}02$.
\item Le côté du carré inscrit dans un cercle de rayon $r$ est $r\sqrt2$ et le
  quart de circonférence $\frac{\pi r}{2}$. La règle $\frac{\pi r/2}{r\sqrt2} =
  \frac{10}{9}$ donne $\pi = \frac{20\sqrt2}{9} \approx 3{,}142697$, trop grand de
  $1{,}1\times 10^{-4}$. Chaque côté du polygone circonscrit est tangent au cercle
  de rayon $\frac12$ en son milieu et vu du centre sous l'angle $\frac{2\pi}{n}$ :
  sa longueur est $2\cdot\frac12\tan\frac{\pi}{n}$, et le périmètre
  $n\tan\frac{\pi}{n}$. Comme $\tan x > x$ sur $]0, \frac{\pi}{2}[$, le périmètre
  dépasse $\pi$. Or $97\tan\frac{\pi}{97} \approx 3{,}142692 < 3{,}142697$ : la
  valeur de Cornu dépasse le périmètre d'un polygone circonscrit, donc elle
  dépasse $\pi$. La remarque est exacte, et de justesse : l'écart est de cinq
  millionièmes.
\item $3 + \frac{112}{113}\cdot\frac17 = 3 + \frac{16}{113} = \frac{355}{113}
  \approx 3{,}1415929$, trop grand d'environ $2{,}7\times 10^{-7}$ seulement.
\item $99^2 = 9801$ et $2\times 70^2 = 9800$, donc $\bigl(\frac{99}{70}\bigr)^2 =
  2 + \frac{1}{4900}$ : $\frac{99}{70}$ approche $\sqrt2$ par excès. C'est une
  solution de $x^2 - 2y^2 = 1$.
\item L'aire du carré de diagonale $10$ est $\frac{10^2}{2} = 50$. Celle du cercle
  de diamètre $8$ est $\frac{22}{7}\times 16 = \frac{352}{7} = 50\frac27$, et non
  $50\frac47$. L'excès est $\frac27$, soit $\frac{2/7}{352/7} = \frac{1}{176}$ de
  l'aire du cercle, non $\frac{1}{88}$ : les deux nombres de la copie doublent
  l'excès vrai, et s'accordent entre eux, non avec le calcul. Avec la vraie
  valeur de $\pi$, le cercle vaut $50{,}27$, et l'excès est d'environ
  $\frac{1}{189}$.
\end{enumerate}
\end{solution}

\begin{exercice}{Les chants de huit notes sans septième ni sixte}{Lycée}{2}
\manuscrit{fr-12357}{21}{21}
La remarque sur la page 113 veut retirer des chants de l'octave (les $8!$ façons
d'ordonner huit notes différentes) ceux où la note \emph{ut} est voisine de
\emph{si} ou de \emph{la}. Elle compte : ceux où \emph{ut} et \emph{si} sont
voisins, $2\cdot 7! = 10\,080$ ; de même pour \emph{ut} et \emph{la} ; en tout
$20\,160$ ; mais on a compté deux fois ceux où \emph{ut} est voisin des deux,
$2\cdot 6! = 1440$. Elle conclut à $21\,600$ chants permis. La copie écrit « des
4320 chants de l'octave ».
\begin{enumerate}
\item Justifier chaque étape par le principe d'inclusion--exclusion.
\item Pourquoi le nombre des chants où \emph{ut} est voisin de \emph{si} et de
  \emph{la} est-il $2 \cdot 6!$ ?
\item Quel est le nombre total de chants, et que penser du « 4320 » ?
\end{enumerate}
\end{exercice}

\begin{solution}
\begin{enumerate}
\item Soit $X$ l'ensemble des chants où \emph{ut} et \emph{si} sont voisins, $Y$
  celui où \emph{ut} et \emph{la} le sont. En collant \emph{ut} et \emph{si} en un
  bloc, on a $7!$ ordres des sept objets et $2$ ordres dans le bloc :
  $|X| = |Y| = 2\cdot7! = 10\,080$. Les chants interdits forment $X \cup Y$, et
  $|X\cup Y| = |X| + |Y| - |X\cap Y| = 20\,160 - 1440 = 18\,720$. Les chants permis
  sont $8! - 18\,720 = 40\,320 - 18\,720 = 21\,600$.
\item \emph{ut} ne peut avoir que deux voisins ; s'il est voisin des deux, le bloc
  est « \emph{si ut la} » ou « \emph{la ut si} » : $2$ blocs, et $6!$ ordres des six
  objets restants.
\item $8! = 40\,320$ ; toute la suite du calcul de la copie suppose ce nombre, et
  « 4320 » est $40\,320$ privé de son zéro, une faute de plume ou de copie. Un
  dénombrement exhaustif des $40\,320$ permutations confirme $21\,600$.
\end{enumerate}
\end{solution}

\begin{exercice}{Cinquante-trois commas, douze demi-tons et les fractions continues}{L2}{2}
\manuscrit{fr-12357}{5}{5}
\manuscrit{fr-12357}{43}{43}
Un feuillet de musique (vue 5) attribue à chaque intervalle un nombre entier de
« commas », l'octave en comptant $53$ : dièse $3$, demi-ton $5$, ton mineur $8$, ton
majeur $9$, tierce mineure $14$, tierce majeure $17$, quarte $22$, quinte $31$,
sixte mineure $36$, sixte majeure $39$. La neuvième « observation » (vue 43)
compare l'intonation juste au tempérament égal à douze demi-tons.
\begin{enumerate}
\item Un intervalle de rapport $r$ vaut $53\log_2 r$ « commas » dans la division
  de l'octave en $53$ parties égales. Calculer ce nombre pour $\frac{128}{125}$,
  $\frac{16}{15}$, $\frac{10}{9}$, $\frac98$, $\frac65$, $\frac54$, $\frac43$, $\frac32$,
  $\frac85$, $\frac53$ et comparer au feuillet.
\item Le développement en fraction continue de $x = \log_2\frac32$ commence par
  $[0; 1, 1, 2, 2, 3, 1, 5, \dots]$. Calculer les réduites jusqu'à $\frac{31}{53}$ et
  expliquer pourquoi $12$ et $53$ sont de bonnes divisions de l'octave pour la
  quinte (on rappelle que $\bigl|x - \frac{p_k}{q_k}\bigr| < \frac{1}{q_kq_{k+1}}$).
\item Dans le tempérament égal, la tierce majeure vaut $4$ demi-tons et la tierce
  mineure $3$. Exprimer leurs écarts à la juste en cents ($1200\log_2$) et, comme la
  copie, en « commas » ($53$ à l'octave, la tierce juste comptée $17$ et $14$).
  Laquelle est la plus proche ?
\item La copie note que, sur un instrument tempéré, « trois quintes font
  exactement une treizième majeure et trois tierces majeures l'octave, ce qui ne
  peut arriver en la vraie théorie ». Calculer les écarts justes.
\end{enumerate}
\end{exercice}

\begin{solution}
\begin{enumerate}
\item On trouve respectivement $1{,}81$ ; $4{,}93$ ; $8{,}06$ ; $9{,}01$ ; $13{,}94$ ;
  $17{,}06$ ; $22{,}00$ ; $31{,}00$ ; $35{,}94$ ; $39{,}06$. Tous les nombres du
  feuillet sont les entiers les plus proches, sauf le dièse : $\frac{128}{125}$ vaut
  $1{,}81$ comma, et le « dièse » de $3$ commas du feuillet est en fait
  $\frac{25}{24}$ ($53\log_2\frac{25}{24} = 3{,}12$), le demi-ton chromatique ; le
  feuillet l'emploie bien ainsi, puisqu'il fait le ton mineur de « dièse plus
  demi-ton » ($\frac{25}{24}\cdot\frac{16}{15} = \frac{10}{9}$, $3 + 5 = 8$).
\item Les réduites sont $\frac01, \frac11, \frac12, \frac35, \frac{7}{12},
  \frac{24}{41}, \frac{31}{53}$ (par $p_k = a_kp_{k-1} + p_{k-2}$, de même pour
  $q_k$). Une division de l'octave en $q$ parts approche la quinte par $p$ parts
  avec l'erreur $|x - p/q|$ ; les réduites sont les meilleures approximations. Pour
  $\frac{7}{12}$, l'erreur est $0{,}00163$ octave, soit $1{,}96$ cent ; pour
  $\frac{31}{53}$, $5{,}7\times10^{-5}$ octave, soit $0{,}07$ cent, conformément à
  $|x - \frac{31}{53}| < \frac{1}{53\times 306}$.
\item Tierce majeure : $400 - 1200\log_2\frac54 = 400 - 386{,}31 = +13{,}69$ cents ;
  tierce mineure : $300 - 315{,}64 = -15{,}64$ cents. En commas : $4\cdot\frac{53}{12}
  = 17\frac23$ contre $17$, trop haute de $\frac23$ ; $3\cdot\frac{53}{12} =
  13\frac14$ contre $14$, trop basse de $\frac34$. Dans les deux mesures, la tierce
  majeure tempérée est plus proche de la juste que la mineure, comme le dit la
  copie.
\item $\bigl(\frac32\bigr)^3 = \frac{27}{8}$ et la treizième majeure juste vaut
  $2\cdot\frac53 = \frac{10}{3}$ : $\frac{27/8}{10/3} = \frac{81}{80}$, le comma
  syntonique. $\bigl(\frac54\bigr)^3 = \frac{125}{64}$ et $\frac{2}{125/64} =
  \frac{128}{125}$, le dièse. En demi-tons égaux, $3\times7 = 21 = 12 + 9$ et
  $3\times4 = 12$ : les deux égalités deviennent exactes.
\end{enumerate}
\end{solution}

\begin{exercice}{La latitude de Paris par l'étoile polaire et par le Soleil}{Lycée}{1}
\manuscrit{naf-9544}{38}{39}
La feuille d'observations tire d'abord la hauteur du pôle d'une étoile de la Petite
Ourse vue au méridien sous le pôle le 3 janvier 1621 ; refaite sans la faute de
soustraction qu'elle y commet, elle donne $48^\circ 47' 40''$. Au verso, « Autres
observations de la même année ». La hauteur minimale de
l'étoile polaire au méridien est $46^\circ 2' 30''$ et sa hauteur maximale
$51^\circ 28' 45''$. Le 23 décembre 1620, la hauteur méridienne du Soleil est
$17^\circ 45'$ ; on y ajoute la parallaxe, $2' 58''$, puis la déclinaison australe du
Soleil, $23^\circ 30' 22''$, et l'on obtient la hauteur de l'équateur
$41^\circ 18' 20''$.

\exfigure{latitude}{Le méridien du lieu. La distance polaire $p$ de l'étoile est exagérée pour la lisibilité ; la déclinaison du Soleil est à l'échelle.}

\begin{enumerate}
\item Une étoile circumpolaire à la distance polaire $p$ passe au méridien à la
  hauteur $\varphi + p$ au-dessus du pôle et $\varphi - p$ au-dessous. En déduire
  $\varphi$ et $p$ pour l'étoile polaire. Pourquoi cette méthode ne demande-t-elle
  ni catalogue ni obliquité ?
\item Au solstice d'hiver, montrer que la hauteur méridienne du Soleil vaut
  $90^\circ - \varphi - |\delta|$. Vérifier l'addition de la feuille, en déduire
  $\varphi$, et vérifier la déclinaison du Soleil par $\sin\delta =
  \sin\varepsilon\sin\lambda$ avec $\varepsilon = 23^\circ 31' 30''$ et la longitude
  du Soleil $\lambda = 272^\circ 13'$.
\item La réfraction atmosphérique relève les astres (d'environ $3'$ à $18^\circ$ de
  hauteur). Dans quel sens fausse-t-elle chacune des deux méthodes ? La latitude de
  Paris est aujourd'hui voisine de $48^\circ 50'$ : commenter les trois valeurs
  $48^\circ 47' 40''$, $48^\circ 45' 37\frac12''$ et $48^\circ 41' 40''$.
\end{enumerate}
\end{exercice}

\begin{solution}
\begin{enumerate}
\item La demi-somme des deux hauteurs donne $\varphi$, la demi-différence $p$ :
  $\varphi = \frac12(51^\circ 28' 45'' + 46^\circ 2' 30'') = \frac12\cdot 97^\circ 31' 15''
  = 48^\circ 45' 37\frac12''$ et $p = \frac12\cdot 5^\circ 26' 15'' = 2^\circ 43'
  7\frac12''$, les nombres de la feuille. On n'a besoin que de deux hauteurs
  mesurées, ni de la position de l'étoile ni de l'obliquité.
\item La hauteur de l'équateur au méridien est $90^\circ - \varphi$ ; un astre de
  déclinaison australe $|\delta|$ passe plus bas de $|\delta|$. D'où
  $90^\circ - \varphi = h + |\delta|$. Addition : $17^\circ 45' + 2' 58'' =
  17^\circ 47' 58''$ et $17^\circ 47' 58'' + 23^\circ 30' 22'' = 41^\circ 18' 20''$ ;
  $\varphi = 48^\circ 41' 40''$, soustraction que la feuille ne fait pas. Et
  $\sin 23^\circ 31' 30''\,\sin 272^\circ 13' = 0{,}39911\times(-0{,}99925) =
  -0{,}39881$, soit $\delta = -23^\circ 30' 23''$, à une seconde de la feuille.
\item La réfraction augmente les hauteurs observées. Pour l'étoile polaire, elle
  relève les deux passages presque également ; la demi-somme est trop forte d'environ
  une minute, et $\varphi$ surestimée. Pour le Soleil, $h$ trop fort donne une hauteur
  de l'équateur trop forte et donc $\varphi$ trop faible, de quelque $3'$ ; la parallaxe
  de $2' 58''$ ajoutée par la feuille (la vraie est d'environ $9''$) et l'obliquité de
  Tycho, trop forte d'environ $2'$, vont dans le même sens. D'où une valeur solaire
  plus basse que les valeurs stellaires, et toutes un peu au-dessous de $48^\circ 50'$ ;
  la feuille ne dit pas où les observations furent faites.
\end{enumerate}
\end{solution}

\begin{exercice}{Le partage des enjeux}{Lycée}{1}
\manuscrit{naf-5176}{40}{41}
Deux joueurs de même force jouent une partie en plusieurs jeux ; on l'interrompt,
et l'on veut partager la mise selon les chances de chacun. On note $P(i,j)$ la
part (en fraction de la mise totale) du joueur à qui il manque $i$ jeux, quand il
en manque $j$ à l'autre.
\begin{enumerate}
\item Justifier que $P(0,j) = 1$, $P(i,0) = 0$ et
  $P(i,j) = \frac12\bigl(P(i-1,j) + P(i,j-1)\bigr)$.
\item Calculer $P(1,2)$ : c'est la « seconde règle » du feuillet, qui donne
  $\frac34$ à qui il ne manque qu'un jeu contre deux. Dans une partie en deux jeux
  où chacun a misé $6$ écus, que reçoit celui qui a gagné le premier jeu ?
\item Partie en trois jeux, $8$ écus contre $8$ : calculer la part de $A$ après
  un jeu à rien, deux à rien, deux contre un. En déduire ce que « vaut » chacun
  des trois jeux gagnés par $A$ (l'accroissement de sa part), et retrouver les
  $\frac{3}{16}$, $\frac{3}{16}$, $\frac{2}{16}$ du feuillet.
\item Partie en quatre jeux : vérifier les valeurs $\frac{15}{16}$, $\frac{14}{16}$,
  $\frac{13}{16}$, $\frac{11}{16}$, $\frac{21}{32}$ du feuillet pour les positions
  trois à rien, trois contre un, deux à rien, deux contre un, un à rien, et la
  note selon laquelle « il vaut mieux avoir 2 jeux à 0 de 4 qu'un à 0 de deux ».
\item La remarque « Autre » dit que, dans une partie en deux, « le premier vaudra
  $\frac34$ ». Qu'est-ce qui est juste, qu'est-ce qui glisse ?
\end{enumerate}
\end{exercice}

\begin{solution}
\begin{enumerate}
\item S'il ne manque rien au premier, il a gagné : $P(0,j) = 1$ ; s'il ne manque
  rien à l'autre, il a perdu : $P(i,0) = 0$. Sinon, le jeu suivant est gagné par
  l'un ou l'autre avec probabilité $\frac12$, et mène en $(i-1, j)$ ou $(i, j-1)$ ;
  la part équitable est la moyenne des deux parts, ce qui est la « première
  règle » du feuillet (avoir des chances égales d'obtenir $10$ ou $12$ vaut $11$).
\item $P(1,1) = \frac12(1 + 0) = \frac12$ et $P(1,2) = \frac12\bigl(P(0,2) + P(1,1)\bigr)
  = \frac12\bigl(1 + \frac12\bigr) = \frac34$. Sur $12$ écus, le gagnant du premier
  jeu en reçoit $9$ et l'autre $3$.
\item $P(1,1) = \frac12$, $P(1,2) = \frac34$, $P(2,1) = \frac14$, $P(2,2) = \frac12$,
  $P(1,3) = \frac12(1 + \frac34) = \frac78$, $P(2,3) = \frac12(\frac78 + \frac12)
  = \frac{11}{16}$. Un jeu à rien : $A$ doit en gagner deux, $B$ trois,
  $P(2,3) = \frac{11}{16}$ ; deux à rien : $P(1,3) = \frac{14}{16}$ ; deux contre
  un : $P(1,2) = \frac{12}{16}$. Partant de $\frac{8}{16}$, le premier jeu vaut
  $\frac{11}{16} - \frac{8}{16} = \frac{3}{16}$, le second
  $\frac{14}{16} - \frac{11}{16} = \frac{3}{16}$, et, si $B$ a gagné entre-temps
  (deux contre un), le troisième vaut $1 - \frac{14}{16} = \frac{2}{16}$ à partir
  de la position deux à rien ; les trois jeux de $A$ font bien
  $\frac{3+3+2}{16} = \frac12$ de la mise totale, soit $8$ écus sur $16$. Ce sont
  les nombres que donne la lettre de Pascal à Fermat du 29 juillet 1654 ($12$, $12$
  et $8$ pistoles sur $64$) : la récurrence impose ces valeurs à quiconque la
  résout juste, et le feuillet ne nomme personne.
\item $P(1,4) = \frac12(1 + P(1,3)) = \frac{15}{16}$ (trois à rien) ;
  $P(1,3) = \frac{14}{16}$ (trois contre un) ; $P(2,4) = \frac12\bigl(P(1,4) +
  P(2,3)\bigr) = \frac12\bigl(\frac{15}{16} + \frac{11}{16}\bigr) = \frac{13}{16}$
  (deux à rien) ; $P(2,3) = \frac{11}{16}$ (deux contre un) ;
  $P(3,4) = \frac12\bigl(P(2,4) + P(3,3)\bigr) = \frac12\bigl(\frac{13}{16} +
  \frac12\bigr) = \frac{21}{32}$ (un à rien). Deux jeux à zéro dans une partie en
  quatre valent $\frac{13}{16}$, un jeu à zéro dans une partie en deux vaut
  $\frac34 = \frac{12}{16}$ : la note a raison, l'écart est $\frac{1}{16}$.
\item Il est juste qu'à égalité la partie recommence avec le nombre de jeux
  restants. Mais « le premier vaudra $\frac34$ » confond la part du gagnant
  ($\frac34$ de la mise) avec ce que vaut le jeu au sens employé partout ailleurs
  sur le feuillet, l'accroissement de la part : $\frac34 - \frac12 = \frac14$.
\end{enumerate}
\end{solution}

\begin{exercice}{La formule générale des partis, et trois joueurs}{L1}{2}
\manuscrit{naf-5176}{40}{41}
On garde les notations de l'exercice précédent.
\begin{enumerate}
\item Montrer que la partie est décidée au plus tard après $N = i + j - 1$ jeux,
  et que l'on peut supposer, sans changer les chances, que l'on joue toujours ces
  $N$ jeux. En déduire
  \[
  P(i,j) = \frac{1}{2^{N}}\sum_{k=i}^{N}\binom{N}{k}.
  \]
  Vérifier que cette formule satisfait la récurrence et les conditions au bord.
\item Trois joueurs $A$, $B$, $C$ de même force jouent une partie en deux jeux.
  Le feuillet considère deux positions : $A$ et $B$ ont chacun un jeu, $C$ aucun ;
  puis $A$ a un jeu, $B$ et $C$ aucun. Calculer les parts, et retrouver
  $\frac{12}{27}$, $\frac{12}{27}$, $\frac{3}{27}$ puis $\frac{17}{27}$,
  $\frac{5}{27}$, $\frac{5}{27}$.
\item La « règle générale » du feuillet : une égale facilité à obtenir $a$, $b$,
  $c$ ou $d$ vaut $\frac{a+b+c+d}{4}$. Le feuillet la justifie par un jeu
  équitable. Pour deux valeurs : montrer qu'avec $11$ en main, en jouant $11$
  contre $11$ à condition que le gagnant donne $9$ au perdant, on se procure
  exactement la même chance d'avoir $9$ ou $13$ ; généraliser à deux valeurs
  $a < b$.
\end{enumerate}
\end{exercice}

\begin{solution}
\begin{enumerate}
\item En $N$ jeux, l'un des deux atteint son but : si le premier gagne moins de
  $i$ jeux, l'autre en gagne au moins $N - (i-1) = j$. Deux joueurs ne peuvent pas
  atteindre tous deux leur but dans ces $N$ jeux ($i + j > N$). Continuer à jouer
  après la décision ne change pas le vainqueur ; le premier gagne donc si et
  seulement si, parmi $N$ jeux indépendants équiprobables, il en gagne au moins
  $i$, d'où la formule. Conditions au bord : pour $i = 0$ (et $j \ge 1$),
  $N = j - 1$ et la somme porte sur tous les $k$ de $0$ à $N$ ; elle vaut $2^N$,
  et $P = 1$. Pour $j = 0$ (et $i \ge 1$), $N = i - 1 < i$, la somme est vide et
  $P = 0$. Récurrence : les positions $(i-1, j)$ et $(i, j-1)$ ont chacune
  $N - 1$ jeux au plus à jouer, donc
  $P(i-1,j) = 2^{-(N-1)}\sum_{k \ge i-1}\binom{N-1}{k}$ et
  $P(i,j-1) = 2^{-(N-1)}\sum_{k \ge i}\binom{N-1}{k}$. Avec la formule de Pascal
  $\binom{N}{k} = \binom{N-1}{k-1} + \binom{N-1}{k}$,
  \[
  \sum_{k \ge i}\binom{N}{k} = \sum_{k\ge i-1}\binom{N-1}{k} + \sum_{k \ge i}\binom{N-1}{k},
  \]
  et en divisant par $2^N$ on obtient $\frac12\bigl(P(i-1,j) + P(i,j-1)\bigr)$. Par
  exemple $P(2,3) = \frac{1}{16}\bigl(\binom42 + \binom43 + \binom44\bigr) =
  \frac{11}{16}$.
\item Première position : au jeu suivant, $A$ gagne tout avec probabilité
  $\frac13$, perd tout si $B$ gagne, et si $C$ gagne, chacun a un jeu et la part
  de chacun est $\frac13$ par symétrie. Donc $A$ a
  $\frac13\bigl(1 + 0 + \frac13\bigr) = \frac49 = \frac{12}{27}$, de même $B$, et
  $C$ a $1 - \frac{24}{27} = \frac{3}{27}$. Seconde position : si $A$ gagne, il a
  tout ; si $B$ (ou $C$) gagne, on est dans la première position avec $A$ parmi
  les deux qui ont un jeu. Donc $A$ a
  $\frac13\bigl(1 + \frac{12}{27} + \frac{12}{27}\bigr) = \frac{17}{27}$, et $B$
  et $C$ ont chacun $\frac12\bigl(1 - \frac{17}{27}\bigr) = \frac{5}{27}$. Ce
  sont les positions des propositions VIII et IX du traité de Huygens, et la
  seconde est celle dont discutent Pascal et Fermat en 1654 ; rapprochements
  seulement.
\item En jouant $11$ contre $11$, le gagnant ramasse $22$ et en donne $9$ : il lui
  reste $13$ ; le perdant reçoit $9$. Le jeu étant équitable, avoir $11$ en main
  vaut avoir une chance égale d'obtenir $9$ ou $13$. En général, avec
  $x = \frac{a+b}{2}$ en main, on joue $x$ contre $x$, le gagnant donnant $a$ au
  perdant : le gagnant a $2x - a = b$, le perdant $a$. C'est la définition de
  l'espérance d'une variable à deux valeurs équiprobables, justifiée par
  l'équivalence avec un jeu équitable ; c'est aussi, presque mot pour mot, la
  démonstration de la première proposition du traité de Huygens de 1657.
\end{enumerate}
\end{solution}

\begin{exercice}{L'aire et la tangente des trochoïdes}{L1}{2}
\manuscrit{nal-2338}{3}{4}
La lettre à Mersenne datée « Dat: Flor: Kal: Maij: anno 1644 » et signée
« Euang\textsuperscript{a} Torricellius » définit les trochoïdes : un cercle de
rayon $r$ roule sur une droite, et chaque point de son rayon prolongé, à la
distance $d$ du centre, décrit une « cycloïde secondaire » ; la cycloïde
« primaire » est celle de $d = r$. Le « cercle générateur propre » d'une
trochoïde est le cercle dont le diamètre est égal à sa hauteur. On prend
\[
x(t) = rt - d\sin t, \qquad y(t) = r - d\cos t, \qquad 0 \le t \le 2\pi,
\]
de sorte que la base de l'arche est le segment de la droite $y = r - d$ qui va de
$(0, r-d)$ à $(2\pi r, r-d)$. La lettre énonce que l'espace compris entre la
trochoïde et sa base est au triangle de même base et de même hauteur « comme la
circonférence du cercle générateur propre avec le double de la base, au double de
la base », et qu'il est au cercle générateur propre dans le rapport qu'on en tire
« per conuersionem rationis ».

\exfigure{trochoides}{Une arche de trochoïde pour $d = \frac r2$, $d = r$ (la cycloïde primaire) et $d = \frac{3r}{2}$, décrites par des points du même cercle qui roule. La base de chaque arche est la droite $y = r - d$.}

\begin{enumerate}
\item On suppose $d \le r$. Montrer que $x$ est croissante sur $[0, 2\pi]$, et que
  l'aire de l'espace vaut
  $\mathcal A = \int_0^{2\pi}\bigl(y(t) - (r - d)\bigr)\,x'(t)\,dt = 2\pi r d + \pi d^2$.
\item Vérifier le premier énoncé, puis sa seconde forme « en lignes droites » :
  l'espace est au triangle comme la hauteur de la trochoïde avec le double de
  l'axe de la cycloïde primaire, au double de cet axe. Montrer que l'espace
  diminué du triangle est égal au cercle générateur propre, et en déduire le
  rapport de l'espace à ce cercle (la « conversion de rapport » passe de $a : b$ à
  $a : (a - b)$). Que donne $d = r$ ?
\item La tangente selon la lettre : on trace le cercle générateur propre tangent à
  la base qui passe par le point $P$ de la trochoïde (on prend celui qui a pour
  centre le centre du cercle qui roule, dans la position où il porte $P$) ; la
  tangente en $P$ à ce cercle coupe la base en $T$ ; on prend sur la base, à partir
  de $T$ et du côté convenable, le point $S$ tel que la hauteur de la trochoïde
  soit à celle de la primaire comme $PT$ à $TS$ ; la droite $SP$ est la tangente.
  En écrivant le vecteur vitesse $(x'(t), y'(t))$ comme la somme de la vitesse du
  centre et de celle de la rotation, démontrer la règle. Pourquoi peut-on, comme
  le dit la lettre, remplacer la base par n'importe quelle parallèle ?
\item Montrer que, pour la cycloïde primaire, la tangente en $P$ passe par le point
  le plus haut du cercle qui roule dans la position où il porte $P$.
\item On suppose $d > r$. Montrer que $x$ n'est plus monotone et que la courbe se
  recoupe au-dessus de chacun de ses points les plus bas. Que reste-t-il de la
  question 1 ?
\end{enumerate}
\end{exercice}

\begin{solution}
\begin{enumerate}
\item $x'(t) = r - d\cos t \ge r - d \ge 0$, et $x'$ ne s'annule qu'en des points
  isolés : $x$ est strictement croissante de $[0, 2\pi]$ sur $[0, 2\pi r]$, et
  l'arche est le graphe d'une fonction de l'abscisse au-dessus de la base.
  L'aire de l'espace est l'intégrale en $X$ de la hauteur au-dessus de la base, et
  le changement de variable $X = x(t)$ donne l'intégrale de l'énoncé. Comme
  $y - (r - d) = d(1 - \cos t)$,
  \[
  \mathcal A = d\int_0^{2\pi}(1 - \cos t)(r - d\cos t)\,dt
  = d\int_0^{2\pi}\bigl(r - (r + d)\cos t + d\cos^2 t\bigr)\,dt
  = d\,(2\pi r + \pi d).
  \]
\item La base vaut $B = 2\pi r$, la circonférence génératrice $C = 2\pi d$, la
  hauteur $2d$, et le triangle $\tau = \frac12\cdot 2\pi r\cdot 2d = 2\pi r d$.
  Alors
  \[
  \frac{\mathcal A}{\tau} = \frac{2r + d}{2r} = \frac{2\pi d + 2\cdot 2\pi r}{2\cdot 2\pi r}
  = \frac{C + 2B}{2B},
  \]
  et, en divisant les deux termes par $\pi$, c'est aussi $\frac{2d + 2\cdot 2r}{2\cdot 2r}$,
  où $2d$ est la hauteur de la trochoïde et $2r$ l'axe de la primaire : les deux
  formes sont le même nombre. Ensuite $\mathcal A - \tau = \pi d^2$, le cercle
  générateur propre. La conversion de $\mathcal A : \tau = (C + 2B) : 2B$ donne
  $\mathcal A : (\mathcal A - \tau) = (C + 2B) : C$, c'est-à-dire que l'espace est
  au cercle générateur propre comme $(C + 2B) : C = (2r + d) : d$. Pour $d = r$,
  $\mathcal A = 3\pi r^2$ : l'espace de la cycloïde vaut une fois et demie son
  triangle et trois fois le cercle qui roule.
\item Le centre du cercle qui roule est $\Omega(t) = (rt, r)$, de vitesse $r(1, 0)$,
  et $\overrightarrow{\Omega P} = (-d\sin t, -d\cos t)$ a pour dérivée
  $d\,u$ avec $u = (-\cos t, \sin t)$ ; ainsi
  \[
  (x'(t), y'(t)) = r(1, 0) + d\,u .
  \]
  Le cercle de centre $\Omega(t)$ et de rayon $d$ passe par $P$ et a son point le
  plus bas sur la droite $y = r - d$ : c'est le cercle générateur propre tangent à
  la base de la règle, et $u$ dirige sa tangente en $P$. Si $\sin t \ne 0$, cette
  tangente coupe la base en $T = P + \lambda u$ pour un réel $\lambda \ne 0$.
  Posons $S = T + \lambda\frac{r}{d}(1, 0)$ : alors $TS = \frac rd\,PT$,
  c'est-à-dire $PT : TS = d : r = 2d : 2r$, le rapport des hauteurs, et le signe
  de $\lambda$ fixe le « côté convenable ». On a
  \[
  \overrightarrow{PS} = \lambda u + \lambda\frac rd(1, 0) = \frac{\lambda}{d}\bigl(d\,u + r(1, 0)\bigr)
  = \frac{\lambda}{d}\,(x'(t), y'(t)),
  \]
  donc $SP$ porte la vitesse : c'est la tangente. (Si $\sin t = 0$, $P$ est au plus
  haut ou au plus bas, la vitesse est horizontale et la règle dégénère.) Une
  parallèle à la base se déduit de la base par une homothétie de centre $P$ ; elle
  envoie $T$ et $S$ sur deux points $T'$, $S'$ de la parallèle, conserve le rapport
  $PT' : T'S'$ et la direction de $PS$ : la construction faite sur la parallèle
  donne la même droite.
\item Pour $d = r$, le cercle générateur propre est le cercle qui roule ; son point
  le plus haut est $H = (rt, 2r)$, et $\overrightarrow{PH} = r(\sin t, 1 + \cos t)$.
  La vitesse vaut $r(1 - \cos t, \sin t)$, et
  \[
  \det\bigl(\overrightarrow{PH}, (x', y')\bigr) = r^2\bigl(\sin^2 t - (1 + \cos t)(1 - \cos t)\bigr) = 0 :
  \]
  la tangente passe par $H$. Autrement dit : le point de contact $K = (rt, 0)$ est
  immobile à l'instant considéré, la vitesse de $P$ est perpendiculaire à $KP$, et,
  $KH$ étant un diamètre, l'angle $\widehat{KPH}$ est droit.
\item $x'(0) = r - d < 0$ : $x$ décroît d'abord, et $x(t) < 0$ pour $t > 0$ petit.
  La fonction $\varphi(t) = d\sin t - rt$ vérifie $\varphi(0) = 0$,
  $\varphi'(0) = d - r > 0$ et $\varphi(\pi) = -\pi r < 0$ ; elle s'annule donc en
  un $t_0 \in\, ]0, \pi[$, où $x(t_0) = 0$. Comme $x(-t) = -x(t)$ et $y(-t) = y(t)$,
  l'arche précédente passe au même point $(0, r - d\cos t_0)$ pour $t = -t_0$ : la
  courbe se recoupe au-dessus du point le plus bas $(0, r - d)$ et fait une boucle
  dont ce point est le plus bas ; de même en $2\pi r$. Le calcul de la question 1 reste un calcul
  juste, et l'intégrale vaut toujours $2\pi r d + \pi d^2$ ; mais le changement de
  variable qui en faisait l'aire d'une région sous un graphe demandait $x$
  monotone. Ce n'est plus qu'une aire algébrique, où les portions parcourues à
  rebours ($x' < 0$) comptent négativement, et « l'espace compris entre la courbe
  et sa base » n'est plus une région bornée par un arc simple. La lettre dit
  « quodlibet spatium » sans distinguer : son énoncé vaut tel quel pour la
  cycloïde primaire et les trochoïdes raccourcies ($d < r$).
\end{enumerate}

\exfigure{trochoide-tangente}{La règle de la question 3 pour $d = 0{,}6\,r$ : $PT$ est tangente au cercle générateur propre de centre $\Omega$, $TS = \frac rd\,PT$, et $SP$ est la tangente à la trochoïde.}
\end{solution}

\begin{exercice}{Le solide autour de l'axe : 11 à 18, 16 à 11, une seule erreur}{L2}{3}
\manuscrit{nal-2338}{4}{5}
\manuscrit{nal-2338}{16}{19}
La lettre du 1\textsuperscript{er} mai 1644, signée « Euang\textsuperscript{a}
Torricellius », donne sans démonstration une suite de rapports pour la cycloïde
primaire engendrée par un cercle de rayon $r$ : l'arche
$x = r(t - \sin t)$, $y = r(1 - \cos t)$, $0 \le t \le 2\pi$, et l'axe est la droite
$x = \pi r$. Le centre de gravité de la cycloïde divise l'axe, du sommet à la base,
comme $7$ à $5$ ; le solide engendré autour de la base est au cylindre de même axe
et de même diamètre comme $5$ à $8$ ; le solide engendré autour de l'axe est au
cylindre « de même axe et de même base » comme $11$ à $18$ ; le rapport de ces deux
solides est « ineffable », composé de $44 : 45$ et du rapport d'un cercle au carré
circonscrit ; enfin le centre de gravité de la demi-cycloïde est sur la parallèle
à l'axe qui coupe la base de la cycloïde entière en parties $19$ et $8$, ou celle de
la demi-cycloïde en parties $16$ et $11$, la plus grande du côté de la courbe. La
pièce des vues 16--19, écrite à Mersenne à la première personne et close par
« Torricellium », rapporte un soupçon de Roberval (« quæ Geometricè falsa
suspicatur »), rappelle qu'Archimède donne le cercle au carré du diamètre comme
$11$ à $14$, et demande d'où l'on croit qu'il a tiré « le rapport qu'il réduisait
aux nombres 11 et 18 ». On admet que $\pi^2$ est irrationnel, et la règle de
Pappus--Guldin : une figure plane d'aire $A$, tournant autour d'une droite de son
plan qui ne la traverse pas, engendre un volume $2\pi\delta A$, où $\delta$ est la
distance de son centre de gravité à la droite.

\exfigure{cycloide-centre}{La cycloïde, son centre de gravité $G_0$ sur l'axe, et la demi-cycloïde (ombrée) avec son centre $G$ à la distance $\delta$ de l'axe. La parallèle en pointillé coupe la base de la demi-cycloïde en parties $16$ et $11$, comme le veut la lettre.}

\begin{enumerate}
\item Calculer l'aire de la cycloïde et la hauteur $\bar y$ de son centre de
  gravité. Vérifier les rapports $7 : 5$ et $5 : 8$.
\item Calculer le volume $V$ du solide engendré autour de l'axe (on pourra écrire
  $V = 2\pi\int_0^{\pi r}(\pi r - x)\,y\,dx$), puis son rapport au cylindre de rayon
  $\pi r$ et de hauteur $2r$. Le comparer à $\frac{11}{18}$. Ce rapport peut-il
  être rationnel ?
\item Soit $\delta$ la distance à l'axe du centre de gravité de la demi-cycloïde.
  Montrer que les énoncés $19 : 8$ et $16 : 11$ donnent le même $\delta$, et que ce
  $\delta$, porté dans la règle de Pappus--Guldin, redonne exactement
  $\frac{11}{18}$. Calculer le vrai $\delta$, et les vrais rapports des parties des
  deux bases.
\item Montrer que le « $44 : 45$ » se déduit de $11 : 18$ et de $5 : 8$. Donner la
  vraie valeur du rapport des deux solides, et montrer qu'elle n'est pas de la
  forme $q\cdot\frac{\pi}{4}$ avec $q$ rationnel.
\item Que vaut le rapport de la question 2 si l'on prend $\pi = \frac{22}{7}$ ? Pour
  quelle valeur de $\pi$ vaudrait-il $\frac{11}{18}$ ? Cette valeur est-elle
  compatible avec les bornes d'Archimède $3\frac{10}{71} < \pi < 3\frac17$ ? Que
  répondre à la question de la pièce ?
\end{enumerate}
\end{exercice}

\begin{solution}
\begin{enumerate}
\item Avec $dx = r(1 - \cos t)\,dt$ et $\int_0^{2\pi}(1 - \cos t)^2\,dt = 3\pi$,
  l'aire vaut $A = r^2\cdot 3\pi = 3\pi r^2$. Le moment par rapport à la base est
  $\int \frac{y^2}{2}\,dx = \frac{r^3}{2}\int_0^{2\pi}(1 - \cos t)^3\,dt$ ; en
  développant, les termes en $\cos t$ et $\cos^3 t$ ont une intégrale nulle, et
  $\int_0^{2\pi}(1 + 3\cos^2 t)\,dt = 2\pi + 3\pi = 5\pi$. D'où
  $\bar y = \frac{5\pi r^3/2}{3\pi r^2} = \frac{5r}{6}$. L'axe mesure $2r$ ; le
  centre est à $\frac{7r}{6}$ du sommet et à $\frac{5r}{6}$ de la base : $7 : 5$,
  exact. Autour de la base, Pappus--Guldin donne
  $2\pi\cdot\frac{5r}{6}\cdot 3\pi r^2 = 5\pi^2r^3$, et le cylindre de rayon $2r$ et
  de longueur $2\pi r$ vaut $8\pi^2r^3$ : $5 : 8$, exact.
\item Le solide est engendré par la moitié de l'espace située d'un côté de l'axe ;
  par tranches cylindriques,
  \[
  V = 2\pi\int_0^{\pi r}(\pi r - x)\,y\,dx
  = 2\pi r^3\int_0^{\pi}(\pi - t + \sin t)(1 - \cos t)^2\,dt .
  \]
  Avec $(1 - \cos t)^2 = \frac32 - 2\cos t + \frac12\cos 2t$, et, par parties,
  $\int_0^\pi(\pi - t)\cos t\,dt = 2$, $\int_0^\pi(\pi - t)\cos 2t\,dt = 0$, on
  trouve $\int_0^\pi(\pi - t)(1 - \cos t)^2\,dt = \frac{3\pi^2}{4} - 4$ ; et
  $\int_0^\pi \sin t\,(1 - \cos t)^2\,dt = \bigl[\frac{(1 - \cos t)^3}{3}\bigr]_0^\pi
  = \frac83$. D'où
  \[
  V = 2\pi r^3\Bigl(\frac{3\pi^2}{4} - \frac43\Bigr) = \frac{\pi r^3}{6}\,(9\pi^2 - 16).
  \]
  Le cylindre vaut $\pi(\pi r)^2\cdot 2r = 2\pi^3r^3$, et le rapport est
  \[
  \frac{9\pi^2 - 16}{12\pi^2} = \frac34 - \frac{4}{3\pi^2} \approx 0{,}61491,
  \]
  contre $\frac{11}{18} \approx 0{,}61111$ : la lettre donne un rapport trop petit
  de $0{,}0038$, soit $0{,}6$ pour cent de la valeur vraie. Si ce rapport était un
  rationnel $q$, on aurait $\pi^2 = \frac{4}{3(3/4 - q)}$ rationnel : il est
  irrationnel, et $11 : 18$ ne peut pas être exact.
\item La demi-cycloïde a pour base le segment de longueur $\pi r$ qui va du point
  de rebroussement au pied de l'axe. L'énoncé $16 : 11$, la plus grande partie
  du côté de la courbe, met le centre à $\delta = \frac{11}{27}\pi r$ de l'axe.
  L'énoncé $19 : 8$ coupe la base entière, de longueur $2\pi r$, en
  $\frac{38}{27}\pi r$ et $\frac{16}{27}\pi r$ : la parallèle est à
  $\frac{16}{27}\pi r$ d'un point de rebroussement, donc à
  $\pi r - \frac{16}{27}\pi r = \frac{11}{27}\pi r$ de l'axe. C'est le même énoncé.
  Le solide autour de l'axe est celui qu'engendre la demi-cycloïde, d'aire
  $\frac32\pi r^2$ ; la règle donne $V = 2\pi\delta\cdot\frac32\pi r^2 =
  3\pi^2r^2\delta$, et
  \[
  \frac{V}{2\pi^3r^3} = \frac{3\delta}{2\pi r} = \frac{3}{2}\cdot\frac{11}{27} = \frac{11}{18}
  \]
  pour $\delta = \frac{11}{27}\pi r$. Les deux énoncés faux de la lettre sont donc une seule donnée lue dans les deux
  sens ; la lettre de Roberval à Torricelli le dit à sa manière, qui écrit que ce
  rapport a été « déduit du centre trouvé auparavant » (vue 79). Le vrai $\delta$
  est $\frac{V}{3\pi^2r^2} = \frac{9\pi^2 - 16}{18\pi}\,r = \Bigl(\frac{\pi}{2} -
  \frac{8}{9\pi}\Bigr)r \approx 1{,}2879\,r$, contre $\frac{11}{27}\pi r \approx
  1{,}2799\,r$. Les vrais rapports sont, pour la base de la demi-cycloïde,
  $(\pi r - \delta) : \delta = (9\pi^2 + 16) : (9\pi^2 - 16) \approx 1{,}4394$, et non
  $\frac{16}{11} \approx 1{,}4545$ ; pour la base entière,
  $(\pi r + \delta) : (\pi r - \delta) = (27\pi^2 - 16) : (9\pi^2 + 16) \approx
  2{,}3895$, et non $\frac{19}{8} = 2{,}375$. Le soupçon que rapporte la pièce des
  vues 16--19 est fondé pour la demi-cycloïde ; pour la cycloïde entière, $7 : 5$
  est exact.
\item Les deux cylindres sont entre eux comme $2\pi^3r^3 : 8\pi^2r^3 = \pi : 4$,
  le rapport d'un cercle au carré circonscrit. Avec les nombres de la lettre,
  \[
  \frac{V_{\text{axe}}}{V_{\text{base}}} = \frac{11/18}{5/8}\cdot\frac{\pi}{4}
  = \frac{44}{45}\cdot\frac{\pi}{4} = \frac{11\pi}{45} \approx 0{,}7679 :
  \]
  le $44 : 45$ est la conséquence exacte de $11 : 18$ et de $5 : 8$. La vraie valeur
  est $\frac{\pi r^3(9\pi^2 - 16)/6}{5\pi^2r^3} = \frac{9\pi^2 - 16}{30\pi} \approx
  0{,}7727$. Si elle valait $q\frac{\pi}{4}$, on aurait
  $\pi^2\bigl(9 - \frac{15q}{2}\bigr) = 16$ ; $q = \frac65$ donne $0 = 16$, et toute
  autre valeur rend $\pi^2$ rationnel. Le mot « ineffable » (irrationnel) est juste
  pour la valeur de la lettre, et le serait pour la vraie, mais cela demande de
  savoir que $\pi$ n'est racine d'aucune équation du second degré à coefficients
  rationnels ; la composition, elle, est fausse.
\item Avec $\pi = \frac{22}{7}$, $\frac34 - \frac{4}{3\pi^2} = \frac34 -
  \frac{49}{363} = \frac{893}{1452} \approx 0{,}61501$, loin de $0{,}61111$. Le
  rapport vaut $\frac{11}{18}$ quand $\frac{4}{3\pi^2} = \frac34 - \frac{11}{18} =
  \frac{5}{36}$, c'est-à-dire $\pi^2 = \frac{48}{5}$, $\pi = \frac{4\sqrt{15}}{5}
  \approx 3{,}0984$, bien au-dessous de $3\frac{10}{71} \approx 3{,}1408$. Comme
  $\frac34 - \frac{4}{3\pi^2}$ croît avec $\pi$, les bornes d'Archimède
  l'enferment entre $0{,}61484$ et $0{,}61501$, et $\frac{11}{18}$ en sort. La
  réponse à la question de la pièce est donc négative : $11 : 18$ n'est pas la
  valeur vraie « réduite » par le $11 : 14$ d'Archimède. C'est l'argument de la
  lettre de Roberval (vue 79), selon qui ce rapport « s'égare loin hors de ces
  bornes ». La pièce ne dit pas comment il a été obtenu, et le calcul de la
  question 3 montre seulement qu'il est lié au centre $16 : 11$.
\end{enumerate}
\end{solution}

\begin{exercice}{Aucun triangle rectangle en nombres n'a pour aire un carré}{L3}{3}
\manuscrit{naf-5176}{32}{32}
La note de la vue 32 pose, parmi quatre questions latines : « trouver des
triangles rectangles en nombres dont l'aire soit un carré, ou démontrer
l'impossibilité de la question ». On démontre l'impossibilité (c'est ce qu'on
appelle aujourd'hui le théorème du triangle rectangle de Fermat).
\begin{enumerate}
\item Montrer qu'on peut se ramener à un triangle primitif, de côtés
  $m^2 - n^2$, $2mn$, $m^2 + n^2$ avec $m > n \ge 1$ premiers entre eux et de
  parités différentes (on admet cette paramétrisation des triplets primitifs).
\item Montrer que $m$, $n$, $m+n$, $m-n$ sont deux à deux premiers entre eux, et
  que si l'aire est un carré, chacun d'eux est un carré :
  $m = x^2$, $n = y^2$, $m + n = u^2$, $m - n = v^2$.
\item Poser $r = \frac{u+v}{2}$, $s = \frac{u-v}{2}$. Montrer que $r$, $s$ sont des
  entiers positifs, que $r^2 + s^2 = x^2$ et $2rs = y^2$, et en déduire un
  triangle rectangle en nombres, d'aire carrée, d'hypoténuse strictement plus
  petite. Conclure par descente infinie.
\item En considérant le triangle $(X^4 - Y^4, 2X^2Y^2, X^4 + Y^4)$, en déduire
  que $X^4 - Y^4 = Z^2$ n'a pas de solution en entiers avec $XYZ \neq 0$, puis que
  $X^4 + Y^4 = W^4$ n'en a pas non plus.
\end{enumerate}
\end{exercice}

\begin{solution}
\begin{enumerate}
\item Si un triangle en entiers a pour côtés $d\,(a', b', c')$ avec $(a',b',c')$
  primitif, son aire est $d^2$ fois celle du triangle primitif ; l'une est un
  carré si et seulement si l'autre l'est (un rationnel positif $q$ tel que $d^2q$
  soit le carré d'un entier est le carré d'un rationnel, et l'aire primitive est
  entière, donc carré d'un entier). On suppose donc le triangle primitif.
\item $m$ et $n$ sont premiers entre eux et de parités différentes, donc
  $m \pm n$ sont impairs. Un diviseur commun à $m+n$ et $m-n$ divise $2m$ et $2n$ ;
  étant impair, il divise $m$ et $n$, donc vaut $1$. Un diviseur commun de $m$ et
  $m \pm n$ divise $n$, donc vaut $1$ ; de même pour $n$. L'aire vaut
  $mn(m-n)(m+n)$, produit de quatre entiers positifs deux à deux premiers entre
  eux ; si c'est un carré, l'exposant de chaque nombre premier dans le produit est
  pair, et chaque nombre premier ne divise qu'un facteur : chaque facteur est un
  carré.
\item $u$ et $v$ sont impairs ($u^2 = m + n$ et $v^2 = m - n$ sont impairs) et
  $u > v > 0$ : $r$ et $s$ sont des entiers positifs. Puis
  $r^2 + s^2 = \frac{u^2 + v^2}{2} = m = x^2$ et $2rs = \frac{u^2 - v^2}{2} = n
  = y^2$. Le triangle de côtés $r$, $s$, $x$ est rectangle, d'aire
  $\frac{rs}{2} = \frac{y^2}{4}$ ; $y^2 = 2rs$ est pair, donc $y$ est pair et
  l'aire est le carré de l'entier $y/2$. Son hypoténuse vérifie
  $x \le x^2 = m < m^2 + n^2$. On aurait ainsi une suite infinie strictement
  décroissante d'entiers positifs (les hypoténuses), ce qui est impossible.
\item Ce triangle est rectangle :
  $(X^4 - Y^4)^2 + 4X^4Y^4 = (X^4 + Y^4)^2$, et son aire vaut
  $X^2Y^2(X^4 - Y^4)$. Si $X^4 - Y^4 = Z^2$ avec $XYZ \neq 0$ (on peut supposer
  $X > Y > 0$), l'aire vaut $(XYZ)^2$, carré non nul : contradiction. Enfin, si
  $X^4 + Y^4 = W^4$ avec $XYW \ne 0$, alors $W^4 - Y^4 = (X^2)^2$, ce qui vient
  d'être exclu. La note ne dit pas qui propose la question, ni si elle en avait
  la réponse.
\end{enumerate}
\end{solution}

\section*{René Descartes (1596--1650)}
\chapitre{descartes}

Aucun feuillet de ce chapitre n'est de Descartes. NAF 5161, que le catalogue
range sous son nom, réunit vingt feuillets d'écrits dirigés contre lui ou
échangés à son sujet, en cinq mains au moins, sans date et presque tous sans
signature : deux lettres-traités « à un cher ami » contre sa théorie des racines
(vues 3 à 20), le début d'un mémoire contre sa solution du lieu à quatre droites,
une note de travail sur les triangles rectangles en nombres, la copie d'une
lettre signée du nom de Roberval, un mémoire recopié au net sur le « centre
d'agitation » d'un corps qui oscille (vues 31 à 39), et une lettre à un
« Révérend Père » sur les thèses de géométrie du P. Fabri. Aucune de ces pièces
n'est attribuée ici au-delà de ce que portent les feuillets. Les pages citées par
les deux lettres-traités s'accordent, selon la transcription, avec \emph{La
Géométrie} imprimée en 1637 à la suite du \emph{Discours de la méthode}.

\begin{exercice}{Une racine réelle et deux « rêveries »}{Lycée}{1}
\manuscrit{naf-5161}{7}{8}
Dans la \emph{Géométrie}, Descartes dit que l'équation
$x^3 - 6x^2 + 13x - 10 = 0$ a « une racine réelle et deux imaginaires ». L'auteur
de la première lettre-traité répond qu'elle n'en a qu'une, $+2$, et que les deux
autres ne sont « que des rêveries et des extravagances de sa belle méthode ».
\begin{enumerate}
\item Vérifier que $2$ est racine et factoriser le polynôme.
\item Montrer que $2$ est la seule racine réelle.
\item Résoudre l'équation dans $\mathbf{C}$. Qui a raison, et en quel sens ?
\item L'auteur propose, si l'on tient à des racines « réelles » et « imaginaires »,
  d'appeler ainsi les positives et les négatives. Que penser de cette proposition ?
\end{enumerate}
\end{exercice}

\begin{solution}
\begin{enumerate}
\item $8 - 24 + 26 - 10 = 0$. La division par $x - 2$ donne
  $x^3 - 6x^2 + 13x - 10 = (x-2)(x^2 - 4x + 5)$.
\item $x^2 - 4x + 5 = (x-2)^2 + 1 > 0$ pour tout réel $x$ (discriminant
  $16 - 20 = -4$) : le seul zéro réel du produit est $2$.
\item $x^2 - 4x + 5 = 0$ donne $x = 2 \pm i$. L'équation a trois racines complexes,
  $2$, $2 + i$, $2 - i$, dont une réelle : c'est exactement le compte de Descartes,
  et il est devenu un théorème (tout polynôme de degré $n$ a $n$ racines complexes
  comptées avec multiplicité). L'auteur a raison sur les nombres réels --- aucun
  binôme réel $x \pm q$ autre que $x - 2$ ne divise le polynôme --- et tort d'en
  conclure que les deux autres racines n'existent pas.
\item Elle confond deux distinctions : les racines négatives sont des nombres
  réels, et les racines « imaginaires » de Descartes ne sont ni positives ni
  négatives. On ne peut pas renommer l'une par l'autre.
\end{enumerate}
\end{solution}

\begin{exercice}{Le centre d'agitation d'un secteur de cylindre}{L2}{2}
\manuscrit{naf-5161}{31}{39}
Le mémoire recopié au net oppose deux règles pour la longueur du pendule simple
équivalent à un corps qui oscille autour d'un axe horizontal. Il prend un secteur
de cylindre droit homogène tournant autour de son propre axe ; sa section est un
secteur de cercle de centre $I$, de rayon $r$ et d'angle au centre $2\alpha$, dont
la bissectrice $IN$ est verticale. On note $\rho$ la distance d'un élément de masse
à l'axe et $\theta$ l'angle de $\vec{I\,\cdot}$ avec la verticale. On admet que la
longueur du pendule équivalent (centre d'oscillation) est
$\ell = \frac{\int\rho^2\,dm}{\int\rho\cos\theta\,dm}$ (exercice sur les pendules
composés de Français 12357).

\exfigure{secteur}{La section du secteur, tournant autour de l'axe perpendiculaire en $I$ ; tracé pour $2\alpha = 110^\circ$, avec $O$ et $Q$ aux distances des questions 1 et 2.}

\begin{enumerate}
\item Montrer que le centre de gravité est à la distance
  $IO = \frac23 r\frac{\sin\alpha}{\alpha}$ de l'axe, comme le dit le mémoire
  (« comme l'arc est à sa corde, ainsi les deux tiers du demi-diamètre à $IO$ »).
\item Montrer que $\ell = IQ = \frac34 r\frac{\alpha}{\sin\alpha}$ : c'est le
  « centre de percussion » du mémoire.
\item La règle attribuée à Descartes revient à
  $\ell_D = \frac{\int\rho^2\,dm}{\int\rho\,dm}$. Calculer $\ell_D$ et montrer qu'elle
  place le centre trop près de l'axe. Pour quels corps les deux règles
  coïncident-elles ?
\item Pour quel angle au centre $Q$ tombe-t-il sur l'arc ($IQ = r$) ? Que vaut $IQ$
  pour le demi-cylindre ?
\end{enumerate}
\end{exercice}

\begin{solution}
En coordonnées polaires, $dm = \rho\,d\rho\,d\theta$ (densité $1$, épaisseur $1$),
$0 \le \rho \le r$, $-\alpha \le \theta \le \alpha$.
\begin{enumerate}
\item $M = \int dm = \alpha r^2$ et $\int\rho\cos\theta\,dm = \int_0^r\rho^2d\rho
  \int_{-\alpha}^{\alpha}\cos\theta\,d\theta = \frac{2r^3\sin\alpha}{3}$, d'où
  $IO = \frac{2r^3\sin\alpha}{3\alpha r^2} = \frac23r\frac{\sin\alpha}{\alpha}$.
  C'est la proportion du mémoire, puisque $\frac{\text{corde}}{\text{arc}} =
  \frac{2r\sin\alpha}{2r\alpha}$.
\item $\int\rho^2\,dm = 2\alpha\int_0^r\rho^3d\rho = \frac{\alpha r^4}{2}$, donc
  $\ell = \frac{\alpha r^4/2}{2r^3\sin\alpha/3} = \frac34 r\frac{\alpha}{\sin\alpha}$.
\item $\int\rho\,dm = 2\alpha\frac{r^3}{3}$, donc $\ell_D = \frac{\alpha r^4/2}{2\alpha
  r^3/3} = \frac34r$, quel que soit $\alpha$ --- le point $P$ du mémoire, « même au
  demi-cylindre ». Comme $\sin\alpha < \alpha$, $\ell_D < \ell$ : la règle place le
  centre trop près de l'axe. En général $\int\rho\cos\theta\,dm \le \int\rho\,dm$,
  avec égalité si et seulement si $\cos\theta = 1$ pour presque toute la masse,
  c'est-à-dire pour un corps situé tout entier dans le demi-plan vertical issu de
  l'axe vers le bas (une figure plane pendant dans le plan de l'axe) : c'est le
  défaut que le mémoire nomme en disant que la règle oublie « la direction de
  l'agitation ».
\item $IQ = r$ équivaut à $\frac{\alpha}{\sin\alpha} = \frac43$ ; la fonction
  $\alpha \mapsto \alpha/\sin\alpha$ croît de $1$ à $\frac{\pi}{2}$ sur
  $]0, \frac{\pi}{2}]$, et l'équation a une seule solution, $\alpha \approx
  1{,}2756$ rad, soit un angle au centre $2\alpha \approx 146^\circ$ ; au-delà, $Q$
  est hors du secteur. Pour le demi-cylindre, $\alpha = \frac{\pi}{2}$ et
  $IQ = \frac{3\pi}{8}r \approx 1{,}178\,r$.
\end{enumerate}
\end{solution}

\begin{exercice}{La règle des signes de Descartes, bien énoncée}{L2}{3}
\manuscrit{naf-5161}{19}{20}
L'auteur attaque la règle de Descartes « pour connaître combien en chaque équation
il y a de racines positives et combien de négatives », qu'il lit comme un compte
exact : autant de racines positives que de « changements de signes », autant de
négatives que de « signes semblables qui s'entresuivent ». On note $V(P)$ le nombre
de changements de signe dans la suite des coefficients non nuls d'un polynôme réel
$P$, et $Z(P)$ le nombre de ses racines strictement positives, comptées avec
multiplicité.
\begin{enumerate}
\item Montrer que $V(P) \equiv Z(P) \pmod 2$ (comparer les signes du coefficient
  dominant et du coefficient non nul de plus bas degré, et le signe de $P$ près de
  $0^+$ et en $+\infty$).
\item On suppose $P(0) \neq 0$. Montrer que $V(P') \le V(P)$ et, par le théorème de
  Rolle, que $Z(P') \ge Z(P) - 1$. En déduire par récurrence sur le degré que
  $Z(P) \le V(P)$, puis que $V(P) - Z(P)$ est un entier pair et positif. Que faire
  si $P(0) = 0$ ?
\item Vérifier la règle sur l'exemple de Descartes,
  $x^4 - 4x^3 - 19x^2 + 106x - 120 = (x-2)(x-3)(x-4)(x+5)$, puis sur l'objection
  de l'auteur, $x^3 - Bx^2 + B^2x - B^3$. L'auteur a-t-il raison ?
\item L'auteur prend ensuite $x^3 - 13x + 12$ et $x^3 - x^2 + 12$, où « manque » un
  terme, et observe qu'elles ont une racine négative bien qu'il n'y ait aucun signe
  semblable qui s'entresuive. Vérifier, et dire comment compter correctement les
  racines négatives.
\item Il affirme enfin qu'en ces deux formes, $x^3 - sx + t$ et $x^3 - Bx^2 + t$
  (lettres positives), « il peut y avoir deux fausses racines ». Montrer que c'est
  faux.
\end{enumerate}
\end{exercice}

\begin{solution}
\begin{enumerate}
\item Écrivons $P = a_nx^n + \cdots + a_mx^m$, $a_n a_m \ne 0$. Le nombre de
  changements de signe entre $a_n$ et $a_m$ est pair si et seulement si $a_n$ et
  $a_m$ ont le même signe. Par ailleurs $P(x) \sim a_mx^m$ en $0^+$ et
  $P(x) \sim a_nx^n$ en $+\infty$ ; $P$ ne change de signe sur $]0, +\infty[$
  qu'en ses racines positives de multiplicité impaire, donc le nombre de racines
  positives comptées avec multiplicité est pair si et seulement si $P$ a le même
  signe près de $0$ et à l'infini, c'est-à-dire si $a_na_m > 0$. D'où
  $V \equiv Z \pmod 2$.
\item Si $P = \sum_{k=0}^n a_kx^k$ avec $a_0 \ne 0$, les coefficients de
  $P' = \sum k a_k x^{k-1}$ sont $ka_k$, de même signe que $a_k$ pour $k \ge 1$ :
  la suite des signes de $P'$ est celle de $P$ privée de son dernier terme, et
  $V(P') \le V(P)$. Soient $r_1 < \cdots < r_s$ les racines positives de $P$, de
  multiplicités $m_1, \dots, m_s$ ; $r_i$ est racine de $P'$ de multiplicité
  $m_i - 1$, et Rolle donne une racine de $P'$ dans chaque $]r_i, r_{i+1}[$ : donc
  $Z(P') \ge \sum(m_i - 1) + (s - 1) = Z(P) - 1$. Récurrence : pour le degré $0$,
  $Z = V = 0$. Si la propriété vaut pour $P'$ (en divisant au besoin $P'$ par une
  puissance de $x$, ce qui ne change ni $Z$ ni $V$), alors
  $Z(P) \le Z(P') + 1 \le V(P') + 1 \le V(P) + 1$ ; avec $Z(P) \equiv V(P)
  \pmod 2$, on obtient $Z(P) \le V(P)$. Si $P(0) = 0$, on écrit $P = x^mQ$ avec
  $Q(0) \ne 0$ : $Z$ et $V$ sont les mêmes pour $P$ et $Q$.
\item Signes $+,-,-,+,-$ : trois changements et trois racines positives $2, 3, 4$ ;
  pour les négatives on applique la règle à $P(-x) = x^4 + 4x^3 - 19x^2 - 106x - 120$,
  signes $+,+,-,-,-$, un changement et une racine négative $-5$. Pour
  $x^3 - Bx^2 + B^2x - B^3 = (x - B)(x^2 + B^2)$ : trois changements et une seule
  racine positive. L'auteur a raison contre une règle lue comme un compte exact ;
  il n'a pas raison contre la règle bien énoncée, puisque l'écart $3 - 1 = 2$ est
  pair. Les racines manquantes sont $\pm iB$.
\item $x^3 - 13x + 12 = (x-1)(x-3)(x+4)$ et $x^3 - x^2 + 12 = (x+2)(x^2 - 3x + 6)$ :
  racines négatives $-4$ et $-2$. Les suites des coefficients non nuls, $+,-,+$,
  n'ont aucune permanence. La règle des permanences ne vaut que pour un polynôme
  complet ; la règle correcte applique le compte des changements à $P(-x)$ :
  $-x^3 + 13x + 12$ et $-x^3 - x^2 + 12$ ont chacun un changement, donc au plus une
  racine négative, et exactement une par parité.
\item Pour $P = x^3 - sx + t$ ou $x^3 - Bx^2 + t$ ($s, B, t > 0$), $P(-x)$ vaut
  $-x^3 + sx + t$ ou $-x^3 - Bx^2 + t$ : un seul changement de signe, donc au plus
  une racine négative. Et il y en a une : $P(0) = t > 0$ et
  $P(x) \to -\infty$ en $-\infty$. Chacune a exactement une racine négative ;
  l'affirmation de l'auteur est fausse.
\end{enumerate}
\end{solution}

\begin{exercice}{Un second triangle rectangle d'aire 6}{L3}{2}
\manuscrit{naf-5161}{27}{28}
Une note de travail pose : « Il faut trouver un triangle primitif dont l'aire soit
un sextuple quarré et qui ne soit pas semblable au triangle de 3. 4. 5. » ; en
divisant ses côtés par la racine du carré, on aura un triangle rectangle à côtés
rationnels d'aire $6$. Le procédé : du triangle donné, former un triangle avec les
deux plus petits côtés, puis un troisième avec les deux plus grands côtés du
second. « Former un triangle » de deux nombres $p > q$, c'est prendre
$p^2 - q^2$, $2pq$, $p^2 + q^2$.
\begin{enumerate}
\item Exécuter le procédé à partir de $(3,4,5)$. Vérifier que l'aire du troisième
  triangle est $6\cdot 70^2$ et en déduire un triangle rationnel d'aire $6$.
\item Généraliser : si $(a, b, c)$ est un triangle rectangle rationnel d'aire $n$
  avec $a \ne b$, le triangle formé de $c^2$ et $2ab$ a pour côtés
  $(a^2 - b^2)^2$, $4abc^2$, $c^4 + 4a^2b^2$ et pour aire
  $n\bigl(2c(a^2-b^2)\bigr)^2$. En déduire un nouveau triangle rationnel d'aire $n$.
\item Au triangle $(a, b, c)$ on associe le point $M = \bigl(\frac{c^2}{4},
  \frac{c(b^2 - a^2)}{8}\bigr)$. Montrer que $M$ est sur la courbe
  $y^2 = x^3 - n^2x$.
\item Pour $(3,4,5)$, $M = \bigl(\frac{25}{4}, \frac{35}{8}\bigr)$. Calculer
  l'abscisse du point d'intersection (autre que $M$) de la tangente en $M$ avec la
  courbe, et la comparer au nouveau triangle.
\item La même page refait le chemin à partir de $(5, 12, 13)$. Donner le triangle
  rationnel d'aire $30$ que l'on obtient.
\end{enumerate}
\end{exercice}

\begin{solution}
\begin{enumerate}
\item $(3,4) \mapsto (16 - 9, 24, 25) = (7, 24, 25)$ ; $(24, 25) \mapsto
  (625 - 576, 1200, 1201) = (49, 1200, 1201)$. Aire $\frac12\cdot49\cdot1200 =
  29\,400 = 6\cdot 4900 = 6\cdot 70^2$. En divisant par $70$ :
  $\bigl(\frac{7}{10}, \frac{120}{7}, \frac{1201}{70}\bigr)$, d'aire
  $\frac12\cdot\frac{7}{10}\cdot\frac{120}{7} = 6$, et
  $\frac{49}{100} + \frac{14400}{49} = \frac{2401 + 1\,440\,000}{4900} =
  \frac{1\,442\,401}{4900} = \bigl(\frac{1201}{70}\bigr)^2$.
\item Former un triangle de $p = c^2$ et $q = 2ab$ ($p > q$ car
  $c^2 - 2ab = (a-b)^2 > 0$) : $p^2 - q^2 = c^4 - 4a^2b^2 = (a^2+b^2)^2 - 4a^2b^2 =
  (a^2 - b^2)^2$, $2pq = 4abc^2$, $p^2 + q^2 = c^4 + 4a^2b^2$. Aire :
  $\frac12(a^2-b^2)^2\cdot 4abc^2 = 2abc^2(a^2-b^2)^2 = \frac{ab}{2}\,4c^2(a^2 -
  b^2)^2 = n\bigl(2c(a^2-b^2)\bigr)^2$. En divisant les trois côtés par
  $2c|a^2 - b^2|$, on obtient un triangle rectangle rationnel d'aire $n$. Pour
  $(3,4,5)$, $2c|a^2-b^2| = 70$ : c'est le diviseur de la note.
\item Avec $n = \frac{ab}{2}$ : $x^3 - n^2x = x(x^2 - n^2) = \frac{c^2}{4}\cdot
  \frac{c^4 - 4a^2b^2}{16} = \frac{c^2(a^2-b^2)^2}{64} = y^2$.
\item $n = 6$, $y^2 = x^3 - 36x$. Pente de la tangente :
  $\lambda = \frac{3x_0^2 - 36}{2y_0} = \frac{3\cdot\frac{625}{16} - 36}{\frac{35}{4}}
  = \frac{1299/16}{35/4} = \frac{1299}{140}$. La droite coupe la cubique en trois
  points dont la somme des abscisses est $\lambda^2$ (coefficient de $x^2$ dans
  $x^3 - (\lambda x + \mu)^2 - 36x$), $M$ comptant deux fois :
  $x_2 = \lambda^2 - 2x_0 = \frac{1\,687\,401}{19\,600} - \frac{25}{2} =
  \frac{1\,442\,401}{19\,600} = \bigl(\frac{1201}{140}\bigr)^2$. C'est
  $\frac{c'^2}{4}$ pour $c' = \frac{1201}{70}$ : le point associé au nouveau
  triangle. Le procédé de la note est la « duplication » d'un point sur cette
  courbe ; la note ne nomme aucun auteur, et le lien avec les courbes elliptiques
  est bien postérieur.
\item $(5,12) \mapsto (119, 120, 169)$ ; $(120, 169) \mapsto (14\,161, 40\,560,
  42\,961)$, d'aire $30\cdot(2\cdot13\cdot119)^2 = 30\cdot 3094^2$. En divisant
  par $3094$ : $\bigl(\frac{119}{26}, \frac{1560}{119}, \frac{42\,961}{3094}\bigr)$,
  d'aire $\frac12\cdot\frac{119}{26}\cdot\frac{1560}{119} = \frac{780}{26} = 30$.
\end{enumerate}
\end{solution}

\section*{Pierre de Fermat (1601--1665)}
\chapitre{fermat}

NAL 2339, que le catalogue et le feuillet de titre de 1888 intitulent
« Mémoires de Fermat », relie une introduction aux lieux plans et solides, une
méthode « ad disquirendam maximam et minimam » avec une méthode des tangentes,
copiées pour être envoyées par le P. Mersenne à Descartes, un appendice qui
résout les cubiques et les quartiques par l'intersection de deux coniques, et
une méthode d'élimination envoyée à Carcavi avec une lettre du 20 avril 1650.
Les feuillets ne nomment Fermat que dans trois intitulés (vues 37, 41 et 49) ;
les autres pièces ne portent aucun nom, et l'on y dit « l'auteur ». La méthode
des maxima « adégale » deux valeurs voisines d'une expression, divise par
l'accroissement et le supprime : on la refait ici, puis on dit ce qui en fait
une démonstration.

\begin{exercice}{Le plus grand rectangle, par adæqualité puis par la dérivée}{Lycée}{1}
\manuscrit{nal-2339}{27}{28}
Sous le titre « Copie d'un escrit enuoyé par le R. Pere Marçenne a des Cartes »,
qui nomme l'envoyeur et le destinataire mais non l'auteur, le feuillet donne en
une page une méthode pour trouver le maximum et le minimum. On écrit la
grandeur à rendre maximale au moyen de l'inconnue $A$, puis on remplace $A$ par
$A + E$ ; on « adégale » les deux expressions (« Adæquentur, vt Loquitur
Diophantus »), on ôte les termes communs, on divise par $E$, on efface les
termes qui contiennent encore $E$, et l'on égale le reste. L'exemple est de
couper une droite $AC = B$ en un point $E$ de sorte que le rectangle
$AE \cdot EC$ soit le plus grand. On écrit $x$ pour l'inconnue $A = AE$, et $e$
pour l'accroissement, que le feuillet appelle $E$ comme le point.
\begin{enumerate}
\item Exprimer le rectangle $f(x) = AE \cdot EC$, puis $f(x + e)$. Refaire le
  calcul du feuillet jusqu'à sa conclusion, « B æquabitur A bis ».
\item Montrer que, pour $e \neq 0$, $\dfrac{f(x+e) - f(x)}{e} = B - 2x - e$. À
  quoi revient l'opération qui « efface » $e$ ? Retrouver le résultat avec la
  dérivée de $f$.
\item Démontrer que le milieu de $AC$ donne bien le plus grand rectangle, et
  donner sa valeur. Pour $B = 10$, comparer ce rectangle à celui qu'on obtient
  pour $AE = 4$.
\item Appliquer la même règle à la somme des carrés $AE^2 + EC^2$. Quel point
  trouve-t-on ? Est-ce un maximum ? Qu'en conclure sur ce que la règle
  démontre ?
\end{enumerate}
\end{exercice}

\begin{solution}
\begin{enumerate}
\item $EC = B - x$, donc $f(x) = x(B - x) = Bx - x^2$, que le feuillet écrit
  « B in A $-$ Aq ». Puis
  $f(x+e) = (x+e)(B - x - e) = Bx - x^2 + Be - 2xe - e^2$. On adégale
  $f(x+e)$ et $f(x)$ et l'on ôte les termes communs $Bx - x^2$ : il reste
  $Be \approx 2xe + e^2$. On divise par $e$ : $B \approx 2x + e$. On efface $e$ :
  $B = 2x$, et $x = \frac B2$. Il faut couper $AC$ en son milieu.
\item $f(x+e) - f(x) = Be - 2xe - e^2 = e\,(B - 2x - e)$, d'où le quotient
  $B - 2x - e$ pour $e \neq 0$. C'est une fonction affine de $e$ ; « effacer »
  $e$, c'est prendre sa valeur en $e = 0$, qui est sa limite quand $e$ tend vers
  $0$. Cette limite est le nombre dérivé $f'(x) = B - 2x$, et la règle écrit
  $f'(x) = 0$. On retrouve directement $f'(x) = B - 2x$, qui s'annule en
  $x = \frac B2$.
\item $f'(x) = B - 2x$ est positive sur $\left[0, \frac B2\right[$ et négative
  sur $\left]\frac B2, B\right]$ : $f$ croît puis décroît, et son maximum sur
  $[0, B]$ est atteint en $\frac B2$. On le voit aussi sans dérivée :
  $f\bigl(\frac B2 + h\bigr) = \bigl(\frac B2 + h\bigr)\bigl(\frac B2 - h\bigr)
  = \frac{B^2}{4} - h^2 \le \frac{B^2}{4}$, avec égalité seulement pour $h = 0$.
  Le plus grand rectangle est le carré de côté $\frac B2$, d'aire
  $\frac{B^2}{4}$. Pour $B = 10$ : $25$ au milieu, et $4 \times 6 = 24$ pour
  $AE = 4$, c'est-à-dire $25 - 1^2$.
\item $g(x) = x^2 + (B - x)^2 = 2x^2 - 2Bx + B^2$, et
  $g(x+e) - g(x) = 4xe - 2Be + 2e^2$. On divise par $e$ : $4x - 2B + 2e$ ; on
  efface $e$ : $4x = 2B$, soit encore $x = \frac B2$. Mais
  $g\bigl(\frac B2 + h\bigr) = \frac{B^2}{2} + 2h^2 \ge \frac{B^2}{2}$ : le milieu
  donne le \emph{minimum} de la somme des carrés, ce qu'on pouvait prévoir
  puisque $AE^2 + EC^2 = B^2 - 2\,AE \cdot EC$. La règle a donné le même point
  pour un maximum et pour un minimum : elle trouve les points où la dérivée
  s'annule, ce qui est une condition nécessaire pour un extremum intérieur, mais
  elle ne dit pas lequel des deux on a, ni même qu'on en a un. Pour
  $k(x) = x^3$, elle donne $3x^2 = 0$, soit $x = 0$, alors que $x^3$ est
  croissante et n'a en $0$ ni maximum ni minimum. Le « Nec potest generalior
  dari methodus » du feuillet parle de la généralité de la règle, non de ce
  qu'elle démontre : il faut chaque fois une vérification, comme à la question
  précédente.
\end{enumerate}
\end{solution}

\begin{exercice}{La tangente à la parabole : l'inégalité des deux côtés}{L1}{2}
\manuscrit{nal-2339}{29}{30}
Sur le feuillet suivant, de la même main, « De tangentibus Linearum curuarum »
ramène à la méthode des maxima la recherche des tangentes. Soit une parabole de
sommet $d$ et de diamètre $dc$, un point $B$ de la courbe, d'ordonnée $BC$, et
la tangente cherchée $BE$, qui coupe le diamètre en $E$, au-delà du sommet.
L'auteur prend sur $BE$ un point $O$ entre $B$ et $E$, mène l'ordonnée $OI$, et
écrit que, $O$ étant hors de la parabole, le rapport de $cd$ à $di$ est plus
grand que celui de $BC^2$ à $OI^2$, donc que celui de $CE^2$ à $IE^2$. Il pose
$cd = d$, $CE = a$, $CI = e$\footnote{Le feuillet écrit les deux inconnues $A$ et
$E$. Il appelle d'abord $cd$ « B data », puis $d$ ; une note marginale française
le relève : « Il a icy nommé B ce qu'il nomme d par apres ».}, « adégale », et
trouve $a = 2d$. On rapporte la parabole à son sommet : $u$ est l'abscisse
comptée de $d$ vers $c$ le long du diamètre, $v$ l'ordonnée, et la courbe a pour
équation $v^2 = pu$, les coordonnées étant obliques si les ordonnées ne sont pas
perpendiculaires au diamètre ; un point est dit hors de la parabole si
$v^2 > pu$. Pour tout réel $e$, on note $I$ le point du diamètre tel que
$CI = e$, compté positivement de $C$ vers $E$, et $O$ le point de la droite $BE$
situé sur l'ordonnée menée par $I$.

\exfigure{fermat-tangente}{La figure du feuillet, ordonnées perpendiculaires au diamètre : la tangente $BE$, avec $dE = cd$ puisque $CE = 2\,cd$, et un point $O$ de $BE$ entre $B$ et $E$.}

\begin{enumerate}
\item Montrer que $\dfrac{BC^2}{OI^2} = \dfrac{CE^2}{IE^2}$, et que, pour
  $0 < e < d$, le point $O$ est hors de la parabole si et seulement si
  $\dfrac{cd}{di} > \dfrac{BC^2}{OI^2}$.
\item Soit $F(e) = d\,(a - e)^2 - (d - e)\,a^2$. Montrer que, pour tout réel
  $e$, le point $O$ est hors de la parabole si et seulement si $F(e) > 0$, et
  que $F(e) = e\,(a^2 - 2da + de)$. Refaire l'adæqualité du feuillet, et dire
  quelle quantité elle calcule.
\item Le feuillet ne prend que des points $O$ entre $B$ et $E$, c'est-à-dire
  $0 < e < a$. Montrer que $F(e) > 0$ pour tout $e \in \left]0, a\right[$ si et
  seulement si $a \ge 2d$. Pour $a = 3d$, trouver le second point commun de la
  droite $BE$ et de la parabole.
\item Montrer qu'en demandant aussi $F(e) > 0$ pour $e < 0$ voisin de $0$ (des
  points $O$ au-delà de $B$), on obtient $a = 2d$, et qu'alors la droite ne
  rencontre la parabole qu'en $B$. Expliquer pourquoi, l'inégalité étant prise
  des deux côtés, le calcul du feuillet devient une démonstration.
\item Retrouver $CE = 2\,cd$ à l'aide de la dérivée de $u \mapsto \sqrt{pu}$.
\end{enumerate}
\end{exercice}

\begin{solution}
\begin{enumerate}
\item Les ordonnées $OI$ et $BC$ sont parallèles et $O$ est sur $BE$ : les
  triangles $EIO$ et $ECB$ sont semblables (Thalès), d'où
  $\frac{OI}{BC} = \frac{IE}{CE}$, et l'égalité des carrés. Soit $w$ l'ordonnée
  de la parabole au point $I$ : $w^2 = p \cdot di$, et $BC^2 = p \cdot cd$. Le
  point $O$ est hors de la parabole si et seulement si $OI^2 > p \cdot di$,
  c'est-à-dire $OI^2 > BC^2 \cdot \frac{di}{cd}$, soit, puisque $di > 0$,
  $\frac{cd}{di} > \frac{BC^2}{OI^2}$.
\item Le point $O$ a pour abscisse $u = cd - CI = d - e$ et, par Thalès, pour
  ordonnée $v = BC \cdot \frac{a - e}{a}$ (avec son signe). Comme
  $BC^2 = pd$, la condition $v^2 > pu$ s'écrit
  $pd\,\frac{(a-e)^2}{a^2} > p\,(d - e)$, soit, en multipliant par
  $\frac{a^2}{p} > 0$, $d\,(a - e)^2 > a^2(d - e)$ : c'est $F(e) > 0$, sans
  qu'il faille supposer $d - e > 0$. C'est aussi l'inégalité du feuillet après
  le produit des extrêmes et des moyens,
  $da^2 + de^2 - 2dae > da^2 - a^2e$ (« d in Aq $+$ d in Eq $-$ d in A in
  E$''$ maius erit quàm d in A quadr. $-$ Aq in E »). En développant,
  $F(e) = -2dae + de^2 + a^2e = e\,(a^2 - 2da + de)$. L'adæqualité : on ôte
  $da^2$ commun, $de^2 + a^2e \approx 2dae$ ; on divise par $e$,
  $de + a^2 \approx 2da$ ; on efface $de$, $a^2 = 2da$, d'où $a = 2d$. Diviser
  $F(e)$ par $e$ puis faire $e = 0$, c'est calculer
  $F'(0) = a^2 - 2da$ : la règle écrit $F'(0) = 0$.
\item Pour $e > 0$, $F(e) > 0$ équivaut à $a^2 - 2da + de > 0$. Si $a \ge 2d$,
  $a^2 - 2da = a(a - 2d) \ge 0$ et $de > 0$ : l'inégalité vaut pour tout
  $e > 0$. Si $a < 2d$, $a^2 - 2da < 0$, et $a^2 - 2da + de < 0$ pour
  $0 < e < \frac{a(2d - a)}{d}$ : les points $O$ proches de $B$ sont dans la
  parabole. L'inégalité que le feuillet écrit, pour les seuls points entre $B$
  et $E$, ne donne donc que $a \ge 2d$, et non $a = 2d$. Les droites avec
  $a > 2d$, qui font avec le diamètre un angle plus petit que la tangente, restent hors de la
  parabole entre $B$ et $E$ mais la recoupent au-delà de $B$, en
  $e^* = \frac{a(2d - a)}{d} < 0$. Pour $a = 3d$,
  $F(e) = de\,(3d + e)$ et $e^* = -3d$ : le point commun a pour abscisse
  $u = d + 3d = 4d$ et pour ordonnée $v = BC \cdot \frac{3d + 3d}{3d} = 2\,BC$ ;
  on vérifie $v^2 = 4pd = p \cdot 4d$.
\item Pour $e < 0$, $F(e) > 0$ équivaut à $a^2 - 2da + de < 0$. Si
  $a^2 - 2da > 0$, cette quantité est positive pour $e < 0$ assez proche de
  $0$ : il faut donc $a^2 - 2da \le 0$, soit $a \le 2d$. Avec la question
  précédente, $a = 2d$, et alors $F(e) = de^2 > 0$ pour tout $e \neq 0$ : tout
  point de la droite autre que $B$ est hors de la parabole, qui est touchée sans
  être coupée. Pris des deux côtés, l'énoncé « $O$ est hors de la parabole »
  dit que $F$, nulle en $0$, a en $0$ un minimum ; le quotient $F(e)/e$ est
  alors positif pour $e > 0$ et négatif pour $e < 0$, et sa limite $F'(0)$ est
  à la fois positive ou nulle et négative ou nulle : $F'(0) = 0$. C'est exactement le calcul du
  feuillet, qui est donc une condition nécessaire démontrée, et la vérification
  $F(e) = de^2$ en donne la réciproque. D'un seul côté, l'inégalité n'impose que
  $F'(0) \ge 0$, et l'égalité que le feuillet en tire n'est pas justifiée par
  elle.
\item Au point $B$, $u = d$ et $v = \sqrt{pd}$ ; la pente de la tangente est
  $v'(d) = \frac{\sqrt p}{2\sqrt d}$. La tangente rencontre le diamètre
  ($v = 0$) à l'abscisse $u = d - \frac{v(d)}{v'(d)} = d - 2d = -d$ : le point
  $E$ est à la distance $d$ du sommet, de l'autre côté, et $CE = 2d = 2\,cd$. La
  sous-tangente est double de l'abscisse, « quod quidem ita se habet ».
\end{enumerate}
\end{solution}

\begin{exercice}{Éliminer une racine seconde : l'analogie redoublée et le résultant}{L2}{3}
\manuscrit{nal-2339}{41}{43}
La pièce latine que son titre dit envoyée par M. de Fermat à M. de Carcavi le
20 avril 1650 (« A Domino de Fermat ad Dominum De Carcaui Die 20a. Aprilis Anno
1650 Missus ») enseigne à « ramener les racines secondes aux premières »,
c'est-à-dire à éliminer une inconnue entre deux équations. Son exemple, en
écrivant $x$ pour la racine première $A$ et $y$ pour la racine seconde $E$, est
\[ x^3 + y^3 = Z, \qquad Bx + y^2 + Dy = N^2 . \]
On pose $P = Z - x^3$ et $Q = N^2 - Bx$, de sorte que les deux équations
s'écrivent $P = y^3$ et $Q = y^2 + Dy$. L'« analogie de l'égalité double » écrit
la proportion $P : y^3 = Q : (y^2 + Dy)$, égale le produit des extrêmes à celui
des moyens et divise par $y$ ; on recommence jusqu'à ce que $y$ ne paraisse plus
qu'au premier degré, et une division le donne ; on porte enfin cette valeur dans
l'une des équations. Les grandeurs $B$, $D$, $N$, $Z$ sont des réels, avec
$D \neq 0$.
\begin{enumerate}
\item Le feuillet écrit le produit des extrêmes et des moyens
  « $Z^sE^2 - A^3E^2 + Z^sDE - A^3DE$ Æle $N^2E^2 - BAE^3$ ». Trouver la faute.
  Montrer que l'égalité juste, $P(y^2 + Dy) = Qy^3$, est une combinaison des
  deux équations à coefficients polynomiaux, et qu'après division par $y$ elle
  n'est plus que du second degré en $y$.
\item Refaire la seconde analogie, entre $PD = Qy^2 - Py$ et $Q = y^2 + Dy$, et
  en déduire $y\,(Q^2 - DP) = P\,(Q + D^2)$. Vérifier que le quotient du
  feuillet,
  \[ \frac{ZD^2 - x^3D^2 + N^2Z - N^2x^3 - BZx + Bx^4}
          {N^4 - N^2Bx - BxN^2 + B^2x^2 - ZD + x^3D} , \]
  est bien $\dfrac{P(Q + D^2)}{Q^2 - DP}$, et qu'en une solution commune réelle
  avec $y \neq 0$ son dénominateur ne s'annule pas.
\item La division euclidienne de $y^3 - P$ par $y^2 + Dy - Q$, en $y$, donne
  une autre expression de $y$. La trouver, et montrer que
  \[ P(Q + D^2)^2 - (P + DQ)(Q^2 - DP) = D\,R, \qquad
     R = P^2 + D^3P + 3DPQ - Q^3 . \]
  Que peut-on en conclure sur les deux expressions de $y$ ?
\item Montrer que $R$ est le résultant en $y$ des polynômes $y^3 - P$ et
  $y^2 + Dy - Q$ (on pourra utiliser les racines complexes $y_1$, $y_2$ du
  second). Quel est son degré en $x$ ? Que signifie $R(x) = 0$ ?
\item Le dernier précepte porte la valeur de $y$ du feuillet dans
  $Q = y^2 + Dy$. Montrer qu'après avoir chassé le dénominateur on obtient
  $-Q^2R = 0$, équation de degré 8 en $x$. Pour $B = D = 1$, $N^2 = 3$, $Z = 2$,
  vérifier que $(1, 1)$ est solution du système, et que $x = 3$ annule
  l'équation finale sans être l'abscisse d'aucune solution.
\end{enumerate}
\end{exercice}

\begin{solution}
\begin{enumerate}
\item Le produit des extrêmes est $P(y^2 + Dy) = Zy^2 - x^3y^2 + ZDy - x^3Dy$,
  comme au feuillet ; celui des moyens est $Qy^3 = N^2y^3 - Bxy^3$. Le feuillet
  écrit $N^2E^2$ au lieu de $N^2E^3$, alors que le terme voisin, $BAE^3$, a bien
  l'exposant 3 ; la ligne suivante, divisée par $E$, porte $N^2E^2$, qui est
  juste : la faute est restée sans suite dans le calcul. L'égalité juste s'écrit
  \[ P(y^2 + Dy) - Qy^3 = (y^2 + Dy)\,(P - y^3) - y^3\,(Q - y^2 - Dy) , \]
  combinaison des deux équations $P - y^3 = 0$ et $Q - y^2 - Dy = 0$ à
  coefficients polynomiaux : toute solution commune la vérifie. Divisée par $y$,
  elle devient $P(y + D) = Qy^2$, du second degré en $y$ au lieu du troisième ;
  c'est le « vno ad minus gradu depressior » du feuillet. La division perd les
  solutions où $y = 0$, qui demandent $P = 0$ et $Q = 0$ à la fois.
\item La proportion $PD : (Qy^2 - Py) = Q : (y^2 + Dy)$ donne
  $PD\,(y^2 + Dy) = Q\,(Qy^2 - Py)$, soit, développé comme au feuillet,
  $ZDy^2 + ZD^2y - x^3Dy^2 - x^3D^2y = N^4y^2 - N^2Bxy^2 - N^2Zy + N^2x^3y
  - BxN^2y^2 + B^2x^2y^2 + BZxy - Bx^4y$ ; tous les termes du feuillet sont
  justes. Divisée par $y$ : $PD\,(y + D) = Q\,(Qy - P)$, d'où
  $y\,(Q^2 - DP) = PD^2 + PQ = P\,(Q + D^2)$. Au numérateur du feuillet,
  $ZD^2 - x^3D^2 = PD^2$ et $N^2Z - N^2x^3 - BZx + Bx^4 = (N^2 - Bx)(Z - x^3)
  = PQ$ ; au dénominateur, $N^4 - 2N^2Bx + B^2x^2 = Q^2$ (le feuillet laisse
  séparés les deux termes $-N^2BA$ et $-BAN^2$) et $-ZD + x^3D = -DP$. En une
  solution commune,
  $Q^2 - DP = (y^2 + Dy)^2 - Dy^3 = y^2(y^2 + Dy + D^2)
  = y^2\bigl[\bigl(y + \frac D2\bigr)^2 + \frac{3D^2}{4}\bigr]$, strictement
  positif pour $y$ réel non nul.
\item $(y - D)(y^2 + Dy - Q) = y^3 - (Q + D^2)y + DQ$, donc
  \[ y^3 - P = (y - D)(y^2 + Dy - Q) + (Q + D^2)\,y - (P + DQ) . \]
  En une solution commune le reste est nul, et $y = \dfrac{P + DQ}{Q + D^2}$ si
  $Q + D^2 \neq 0$ : une seule division, au lieu de deux analogies. En
  développant, $P(Q + D^2)^2 = PQ^2 + 2D^2PQ + D^4P$ et
  $(P + DQ)(Q^2 - DP) = PQ^2 - DP^2 + DQ^3 - D^2PQ$, dont la différence est
  $DP^2 + D^4P + 3D^2PQ - DQ^3 = D\,R$. Là où les dénominateurs ne s'annulent pas, les deux quotients sont donc
  égaux exactement quand $R = 0$ (on a supposé $D \neq 0$), en particulier en toute
  solution commune : le feuillet et la division euclidienne donnent le même $y$.
\item Le polynôme $g = y^2 + Dy - Q$ est unitaire, de racines $y_1$, $y_2$ avec
  $y_1 + y_2 = -D$ et $y_1y_2 = -Q$. Le résultant de $f = y^3 - P$ et de $g$
  est $\prod_i f(y_i)$ (au signe $(-1)^{3 \cdot 2} = 1$ près), soit
  \[ (y_1^3 - P)(y_2^3 - P) = (y_1y_2)^3 - P\,(y_1^3 + y_2^3) + P^2 . \]
  Or $(y_1y_2)^3 = -Q^3$ et
  $y_1^3 + y_2^3 = (y_1 + y_2)^3 - 3y_1y_2(y_1 + y_2) = -D^3 - 3DQ$, d'où
  $P^2 + D^3P + 3DPQ - Q^3 = R$. Comme $P$ est de degré 3 en $x$ et $Q$ de degré
  1, $R$ est de degré 6, de terme dominant $x^6$ venu de $P^2$ : c'est le
  produit $3 \times 2$ des degrés en $y$. $R(x) = 0$ signifie que, pour cette
  valeur de $x$, les deux équations en $y$ ont une racine complexe commune : les
  racines de $R$ sont exactement les abscisses des solutions complexes du
  système.
\item Avec $y = \frac{P(Q + D^2)}{Q^2 - DP}$, l'équation $Q - y^2 - Dy = 0$,
  multipliée par $(Q^2 - DP)^2$, devient
  \[ Q(Q^2 - DP)^2 - P^2(Q + D^2)^2 - DP(Q + D^2)(Q^2 - DP) = 0 . \]
  En développant, le premier membre vaut
  $Q^5 - 2DPQ^3 + D^2P^2Q - P^2Q^2 - 2D^2P^2Q - D^4P^2
  - DPQ^3 + D^2P^2Q - D^3PQ^2 + D^4P^2$, soit
  $Q^5 - 3DPQ^3 - P^2Q^2 - D^3PQ^2 = -Q^2R$, de degré $2 + 6 = 8$. L'équation
  finale contient donc, outre le résultant, le facteur parasite
  $(N^2 - Bx)^2$. Pour les valeurs proposées, $1 + 1 = 2$ et $1 + 1 + 1 = 3$ :
  $(1, 1)$ est solution ; $P = 2 - x^3$, $Q = 3 - x$, et
  $R = x^6 + 3x^4 - 13x^3 - 9x^2 + 21x - 3
  = (x - 1)(x^5 + x^4 + 4x^3 - 9x^2 - 18x + 3)$, tandis que l'équation finale
  est $(x - 3)^2\,R = 0$. En $x = 3$, $Q = 0$ et $P = -25$ : la seconde équation
  donne $y^2 + y = 0$, soit $y = 0$ ou $y = -1$, et aucun ne vérifie
  $y^3 = -25$ ; d'ailleurs $R(3) = 625 - 25 = 600 \neq 0$. La racine double
  $x = 3$ est étrangère. Le feuillet n'écrit pas cette équation finale et
  affirme seulement que la méthode est générale ; elle élimine toujours, mais
  pour qu'elle soit exacte il faut écarter ce facteur, ce que fait le
  résultant.
\end{enumerate}
\end{solution}

\section*{Gilles Personne de Roberval (1602--1675)}
\chapitre{roberval}

Quatre volumes nomment Roberval, et aucun des feuillets lus ici n'est donné pour
être de sa main. Latin 11195, que le catalogue range sous son nom, relie un
traité latin anonyme, « De Recognitione Æquationum », écrit dans la notation de
l'école de Viète : il fabrique les équations de degré deux, trois et quatre à
partir de leurs racines, et emprunte trois lemmes de maximum à « D. Rob. »,
qu'une note marginale renvoie à « son traité de mécanique des plans inclinés » ;
c'est une citation de Roberval, non une attribution du traité. Guglielmo Libri
décrit, dans un cahier monté dans l'album NAF 9544, cinq manuscrits de Roberval,
dont il rapporte des énoncés. Et une lettre anonyme de NAF 5161, sur les thèses
du P. Fabri, propose « ce beau problème de M. Roberval ». Et NAL 2338 conserve,
en copie, deux lettres dont l'intitulé le nomme, l'une à Mersenne, l'autre à
Torricelli, et la « De Vacuo Narratio » adressée à des Noyers, qui rapporte les
expériences de Rouen de « D. de Paschal ». Les exercices partent de
ces énoncés tels que les feuillets les rapportent ; on écrit les équations sous
forme unitaire, avec $=$ pour la double barre du feuillet.

\begin{exercice}{Mercure, eau et vin dans les tubes de Rouen}{Lycée}{1}
\manuscrit{nal-2338}{97}{103}
La « De Vacuo Narratio », dont le titre nomme Roberval
(« Æ.\textsuperscript{i} P.\textsuperscript{i} de Roberual ») et qui finit sur
« Parisijs 12 Calend. Octob. 1647 », rapporte les expériences faites à Rouen par
« D. de Paschal ». Le mercure s'arrête toujours à $2\frac{7}{24}$ pieds au-dessus
du mercure de la cuvette ; « D. de Paschal » calcule d'avance qu'il faut à l'eau
« environ » $31\frac19$ pieds, et au vin « à peu près » $31\frac23$ pieds, pour lui
faire équilibre, et deux tubes de 40 pieds attachés à un mât de navire le
confirment. La pièce donne enfin le poids de l'eau à celui du mercure, sous le
même volume, « comme 1 à $13\frac47$ ». On admet la loi de la statique des
fluides : deux colonnes de liquides, surmontées de vide, font équilibre au même
air quand les produits de leurs hauteurs par leurs densités sont égaux. On prend
le pied du roi égal à $32{,}5$ cm, et l'on compte 12 pouces au pied.

\exfigure{tubes-rouen}{Les hauteurs de la pièce, à la même échelle, au-dessus de la surface de chaque cuvette ; les deux tubes de 40 pieds étaient attachés à un mât.}

\begin{enumerate}
\item Calculer la hauteur d'eau qui équilibre $2\frac{7}{24}$ pieds de mercure
  avec le rapport $1 : 13\frac47$, et la comparer à $31\frac19$. Quel rapport des
  densités les deux hauteurs de « D. de Paschal » donnent-elles exactement ?
  Convertir $2\frac{7}{24}$ et $31\frac19$ pieds en unités actuelles.
\item Quelle densité du vin, rapportée à celle de l'eau, les hauteurs $31\frac19$
  et $31\frac23$ supposent-elles ? Les adversaires avaient accordé que le vin,
  plus riche en « esprits », laisserait plus de vide que l'eau au sommet de
  tubes de même hauteur. Pourquoi l'expérience les réfute-t-elle ?
\item La pièce affirme que, dans un tube qui contient du mercure surmonté d'eau
  et qui plonge dans une cuvette de mercure pur, la hauteur de l'eau vaut
  $13\frac47$ fois ce qui manque au mercure pour atteindre $2\frac{7}{24}$ pieds.
  Le démontrer. Un tube droit de section constante, dont le sommet est à 4 pieds
  au-dessus de la cuvette, contient $2\frac12$ pieds de mercure et $1\frac12$ pied
  d'eau ; on l'ouvre, et dans l'état final le mercure occupe le bas du tube et
  l'eau le surmonte. Calculer la hauteur finale du mercure et celle du vide laissé
  au sommet (on néglige la variation du niveau de la cuvette).
\item Si l'on verse sur le mercure de la cuvette une hauteur $e$ d'eau, les
  colonnes montent, dit la pièce, de la $13\frac47$\textsuperscript{e} partie de
  $e$. Le justifier, et calculer la montée en pouces pour $e = 1$ pied.
\item Pourquoi aucun siphon ne fait-il passer l'eau par-dessus une montagne de
  plus de 31 pieds, ni aucune pompe qui ne fait qu'aspirer ne l'élève-t-elle plus
  haut, alors qu'une pompe qui refoule le peut ?
\end{enumerate}
\end{exercice}

\begin{solution}
\begin{enumerate}
\item L'équilibre s'écrit $\rho_e h_e = \rho_m h_m$, soit $h_e = 13\frac47\,h_m$.
  Avec $h_m = \frac{55}{24}$ et $13\frac47 = \frac{95}{7}$,
  \[
  h_e = \frac{55}{24}\cdot\frac{95}{7} = \frac{5225}{168} = 31\tfrac{17}{168}
  \approx 31{,}101
  \]
  pieds, contre $31\frac19 \approx 31{,}111$ : un centième de pied d'écart, un huitième de
  pouce environ, et la pièce dit « circiter ». Les deux hauteurs de
  « D. de Paschal » donnent exactement le rapport
  $\frac{280/9}{55/24} = \frac{448}{33} = 13\frac{19}{33} \approx 13{,}576$, contre
  $13\frac47 \approx 13{,}571$ : trois dix-millièmes d'écart relatif, les deux
  données s'accordent. En unités actuelles, $2\frac{7}{24}\times 32{,}5 \approx
  74{,}5$ cm et $31\frac19\times 0{,}325 \approx 10{,}1$ m. La pression moyenne au
  niveau de la mer tient aujourd'hui 76 cm de mercure, ou $10{,}3$ m d'eau :
  l'écart, de 2 pour cent environ, est de l'ordre des variations de la pression d'un
  jour à l'autre. Le rapport actuel des densités du mercure et de l'eau, à la
  température ordinaire, est voisin de $13{,}55$.
\item Les hauteurs sont en raison inverse des densités :
  $\frac{\rho_v}{\rho_e} = \frac{31\frac19}{31\frac23} = \frac{280/9}{95/3} =
  \frac{56}{57} \approx 0{,}982$. Le vin est un peu plus léger que l'eau, donc il
  monte plus haut, et dans deux tubes de même hauteur le vide au-dessus du vin est
  plus court que le vide au-dessus de l'eau, de $31\frac23 - 31\frac19 = \frac59$
  pied, soit $6\frac23$ pouces. C'est exactement le contraire de ce que prédisait
  la doctrine des esprits, et exactement ce qu'annonçait le calcul par les poids,
  fait avant l'expérience : seul le poids du liquide compte.
\item À la hauteur de la surface de la cuvette, la colonne du tube (mercure de
  hauteur $h$, eau de hauteur $h_e$, vide au-dessus) doit peser autant que la
  colonne de mercure seul qui fait équilibre à l'air :
  $\rho_m h + \rho_e h_e = \rho_m H$, avec $H = 2\frac{7}{24}$ pieds. D'où
  \[
  h_e = \frac{\rho_m}{\rho_e}\,(H - h) = 13\tfrac47\,(H - h),
  \]
  ce qu'affirme la pièce. Dans l'application, l'eau garde sa hauteur de $1\frac12$
  pied : la section est constante, et l'eau, plus légère, ne peut sortir par le bas
  à travers le mercure. Le mercure s'écoule en partie dans la cuvette jusqu'à
  \[
  h = \frac{55}{24} - \frac{7}{95}\cdot\frac32 = \frac{55}{24} - \frac{21}{190}
  = \frac{4973}{2280} \approx 2{,}181
  \]
  pieds. Le haut de l'eau est à $h + \frac32 \approx 3{,}681$ pieds, et le vide mesure
  $4 - \frac{8393}{2280} = \frac{727}{2280} \approx 0{,}319$ pied, soit un peu moins
  de quatre pouces. On vérifie que $\frac32 = \frac{95}{7}\cdot\frac{21}{190}$.
\item La surface du mercure de la cuvette porte maintenant l'air et l'eau versée.
  Si $h'$ est la nouvelle hauteur du mercure au-dessus de cette surface,
  $\rho_m h' = \rho_m H + \rho_e e$, d'où $h' - H = \frac{\rho_e}{\rho_m}\,e =
  \frac{7}{95}\,e$ : la $13\frac47$\textsuperscript{e} partie de $e$. Pour
  $e = 1$ pied, la montée est $\frac{7}{95} \approx 0{,}074$ pied, soit
  $\frac{84}{95} \approx 0{,}88$ pouce. Le raisonnement ne change pas si le tube
  contient aussi de l'eau : seul compte le poids posé sur la surface de la cuvette.
\item Dans une colonne d'eau immobile, la pression à la hauteur $z$ au-dessus
  d'une surface libre exposée à l'air vaut $p_0 - \rho_e g z$, et elle ne peut
  devenir négative : l'eau ne peut être tenue qu'à $z \le \frac{p_0}{\rho_e g}$,
  soit $31\frac19$ pieds environ. Si le haut d'un siphon est plus haut, le liquide
  se sépare au sommet, chaque branche retombe à cette hauteur au-dessus de sa
  cuvette et le sommet paraît vide : c'est ce que la pièce rapporte pour les
  siphons de mercure au-delà de $2\frac{7}{24}$ pieds. Une pompe aspirante ne fait
  que laisser un vide au-dessus de l'eau, que l'air pousse dans le tuyau : la même
  borne s'applique. Une pompe foulante pousse l'eau par un piston avec la pression
  que lui donne la machine, qui n'est pas bornée par celle de l'air.
\end{enumerate}
\end{solution}

\begin{exercice}{Les lemmes de maximum empruntés à « D. Rob. »}{L1}{1}
\manuscrit{latin-11195}{9}{10}
Le traité emprunte à « D. Rob. » deux lemmes de maximum. Le premier : de tous les
solides $zA - A^3$ (« appliqués » au plan $z$ et « déficients » d'un cube), le plus
grand est celui qui est « appliqué aux deux tiers », « double du défaut ». Le
second : si une droite $x$ est divisée en deux parties $u$ et $v$, le plus grand
solide $u^2v$ est celui où $u$ est double de $v$.
\begin{enumerate}
\item Trouver le maximum de $f(A) = zA - A^3$ sur $A > 0$ ($z > 0$) et vérifier les
  deux descriptions du feuillet.
\item Trouver le maximum de $u^2v$ sous $u + v = x$, $u, v \ge 0$.
\item La cubique $(A - B)(A - C)(A + B + C)$, avec $B, C > 0$ (deux côtés au-dessus
  et un côté au-dessous égal à leur somme), s'écrit $A^3 - zA + T$. Exprimer $z$ et
  $T$. Montrer qu'une équation $A^3 - zA + T = 0$ avec $z, T > 0$ a deux racines
  positives (distinctes ou non) si et seulement si $27T^2 \le 4z^3$, et qu'elle a
  toujours une racine négative.
\item De même, montrer que $A^3 - xA^2 + T = 0$ ($x, T > 0$) a deux racines
  positives si et seulement si $T \le \frac{4}{27}x^3$, la « limitation » que le
  traité écrit en clair.
\end{enumerate}
\end{exercice}

\begin{solution}
\begin{enumerate}
\item $f'(A) = z - 3A^2$ s'annule en $A_0 = \sqrt{z/3}$, $f$ croît puis décroît :
  le maximum est $f(A_0) = A_0(z - A_0^2) = \frac{2z}{3}\sqrt{\frac z3}$. Avec
  $\ell = A_0$, $z = 3\ell^2$ : le maximum vaut $2\ell^3$, double du « défaut »
  $\ell^3$, et c'est $\frac23 z\cdot\ell$ : le solide de base les deux tiers du plan
  et de hauteur $\ell$, la droite « dont le carré est le tiers du plan ».
\item $g(u) = u^2(x - u)$ a pour dérivée $2ux - 3u^2$, nulle en $u = \frac23x$ ;
  alors $v = \frac13 x$, $u = 2v$, et le maximum vaut $\frac49x^2\cdot\frac{x}{3} =
  \frac{4}{27}x^3$.
\item Le coefficient de $A^2$ est $(B + C) - B - C = 0$ ; celui de $A$ est
  $BC - (B+C)^2 = -(B^2 + BC + C^2)$ ; le terme constant $BC(B+C)$. Donc
  $z = B^2 + BC + C^2$, $T = BC(B+C)$. Soit $h(A) = A^3 - zA + T$ : $h(0) = T > 0$
  et $h \to -\infty$ en $-\infty$, d'où une racine négative. Sur $]0, +\infty[$,
  $h$ décroît puis croît, avec un minimum en $A_0 = \sqrt{z/3}$ égal à
  $T - \frac{2z}{3}\sqrt{\frac z3}$ ; $h$ a des racines positives si et seulement
  si ce minimum est $\le 0$, c'est-à-dire $T \le \frac{2z}{3}\sqrt{\frac{z}{3}}$,
  soit $27T^2 \le 4z^3$ (deux racines, confondues en cas d'égalité). C'est la
  condition de discriminant $4z^3 - 27T^2 \ge 0$, que le traité atteint par un
  maximum.
\item $k(A) = A^3 - xA^2 + T$ : $k'(A) = A(3A - 2x)$, minimum sur $]0, +\infty[$ en
  $A = \frac23x$, égal à $T - \frac{4}{27}x^3$. D'où la condition
  $T \le \frac{4}{27}x^3$, qui est encore la positivité du discriminant
  $T(4x^3 - 27T)$.
\end{enumerate}
\end{solution}

\begin{exercice}{Le centre des demi-fuseaux paraboliques}{L1}{2}
\manuscrit{nal-2338}{83}{83}
La lettre intitulée « Epistola Ægid. Personerii de Roberual ad R. P. Mersennum »
examine des propositions de Torricelli. À l'article 10, elle offre le centre de
gravité de la parabole « a priori », et celui des paraboles de tous les degrés et
de leurs solides, et donne sans démonstration cet exemple : dans le demi-fuseau
parabolique, le centre de gravité divise l'axe en deux parties, celle du côté du
sommet étant à celle du côté de la base comme $11$ à $5$ pour la parabole
« quadratiue », $13$ à $7$ pour la cubique, $15$ à $9$ pour la
« quadrato-quadratiue », $17$ à $11$ pour la « quadrato-cubique », « et ainsi à
l'infini, en ajoutant toujours 2 à chacun des termes du rapport précédent ».

Soit une parabole de degré $n$ ($n \ge 1$ entier) de sommet $V$, d'axe
perpendiculaire à sa base, de hauteur $h$, dont la base est une ordonnée de
longueur $2a$ : le point de la courbe dont l'ordonnée est $u$ ($|u| \le a$) est à
la distance $h\,(|u|/a)^n$ de $V$ le long de l'axe. Pour $n = 2$, les carrés des
ordonnées sont comme les portions de l'axe : c'est la parabole « quadratiue ». Le
fuseau est le solide engendré par le segment de parabole tournant autour de sa
base ; le demi-fuseau en est la moitié que coupe le plan perpendiculaire à la base
mené par l'axe de la parabole. Son sommet est une extrémité de la base, sa base un
disque de rayon $h$, son axe la demi-base, de longueur $a$. On admet que le centre
de gravité d'un solide de révolution homogène est sur son axe, à la distance
$\bar s = \int s\,\rho(s)^2\,ds \big/ \int \rho(s)^2\,ds$ de l'origine, où $\rho(s)$
est le rayon de la section à l'abscisse $s$.

\exfigure{demi-fuseau}{Le segment de parabole (ici $n = 2$) et le demi-fuseau qu'engendre sa moitié en tournant autour de la base ; son axe est la demi-base, de longueur $a$.}

\begin{enumerate}
\item Montrer que la section du demi-fuseau à la distance $s$ de sa base est un
  disque de rayon $\rho(s) = h\bigl(1 - (s/a)^n\bigr)$. Calculer le rapport du
  demi-fuseau au cylindre de rayon $h$ et de longueur $a$ qui l'enveloppe : ce sont
  les rapports que la lettre dit avoir trouvés et qu'elle « passe ».
\item Montrer que le centre de gravité est à la distance
  $\bar s = \frac{2n + 1}{4(n + 2)}\,a$ de la base.
\item En déduire que les deux parties de l'axe sont comme $(2n + 7) : (2n + 1)$, et
  vérifier les quatre rapports et la règle de la lettre. Pourquoi écrit-elle
  $15 : 9$ plutôt que $5 : 3$ ?
\item Que donne la règle pour $n = 1$ ? Reconnaître le solide et contrôler le
  résultat par un centre de gravité connu. Que devient le rapport quand $n$ tend
  vers l'infini, et pourquoi fallait-il s'y attendre ?
\end{enumerate}
\end{exercice}

\begin{solution}
\begin{enumerate}
\item Le point de la base à la distance $s$ du pied de l'axe porte, au-dessus de
  lui, un segment de hauteur $h - h(s/a)^n$ jusqu'à la courbe. En tournant autour de
  la base, ce segment décrit un disque perpendiculaire à la base, de rayon
  $\rho(s) = h\bigl(1 - (s/a)^n\bigr)$. Avec $t = s/a$,
  \[
  V = \pi\int_0^a \rho(s)^2\,ds = \pi h^2 a\int_0^1 (1 - t^n)^2\,dt
  = \pi h^2 a\Bigl(1 - \frac{2}{n+1} + \frac{1}{2n+1}\Bigr)
  = \frac{2n^2}{(n+1)(2n+1)}\,\pi h^2 a .
  \]
  Le rapport au cylindre est $\frac{2n^2}{(n+1)(2n+1)}$ : $\frac{8}{15}$,
  $\frac{9}{14}$, $\frac{32}{45}$, $\frac{25}{33}$ pour $n = 2, 3, 4, 5$. Pour
  $n = 2$, ce sont les huit quinzièmes du fuseau parabolique que Libri rapporte
  d'un manuscrit de Roberval (NAF 9544, dans ce chapitre).
\item Le numérateur vaut
  \[
  \int_0^a s\,\rho(s)^2\,ds = h^2a^2\int_0^1 t(1 - t^n)^2\,dt
  = h^2a^2\Bigl(\frac12 - \frac{2}{n+2} + \frac{1}{2n+2}\Bigr)
  = \frac{n^2}{2(n+1)(n+2)}\,h^2a^2 ,
  \]
  et le dénominateur $\frac{2n^2}{(n+1)(2n+1)}\,h^2 a$ ; leur quotient est
  $\bar s = \frac{(2n+1)}{4(n+2)}\,a$.
\item La partie du côté du sommet vaut $a - \bar s = \frac{2n + 7}{4(n+2)}\,a$ et
  celle du côté de la base $\frac{2n+1}{4(n+2)}\,a$ : elles sont comme
  $(2n + 7) : (2n + 1)$. Pour $n = 2, 3, 4, 5$, on trouve $11 : 5$, $13 : 7$,
  $15 : 9$, $17 : 11$, les quatre rapports de la lettre ; et le passage de $n$ à
  $n + 1$ ajoute $2$ à chaque terme, ce qui est la règle. Les quatre énoncés et la
  règle sont exacts. La lettre écrit les termes que donne la règle, et non le
  rapport réduit : $5 : 3$ cacherait la progression $11, 13, 15, 17$ et
  $5, 7, 9, 11$. Que les quatre rapports sortent de ce calcul confirme la lecture
  du solide adoptée ici ; la lettre ne le décrit que par ses mots (« fusi
  parabolici dimidium plano ad ipsius axem erecto resectum »).
\item Pour $n = 1$, la règle donne $9 : 3 = 3 : 1$. Le rayon $h(1 - s/a)$ décroît
  linéairement jusqu'à zéro au sommet : le demi-fuseau est un cône de base le
  disque de rayon $h$, et l'on sait que le centre de gravité d'un cône est au
  quart de son axe à partir de la base ; en effet $\bar s = \frac{3}{12}\,a =
  \frac a4$, et les parties sont comme $\frac34 : \frac14$. Quand $n \to \infty$,
  $(2n + 7)/(2n + 1) \to 1$ : pour $0 \le t < 1$, $t^n \to 0$ et $\rho(s) \to h$, le
  demi-fuseau tend vers le cylindre de rayon $h$ et de longueur $a$, dont le
  centre est au milieu de l'axe.
\end{enumerate}
\end{solution}

\begin{exercice}{Les paraboles de tous les degrés, selon la description de Libri}{L1}{2}
\manuscrit{naf-9544}{42}{49}
Libri rapporte que le premier manuscrit de Roberval contient un petit écrit sur la
quadrature de la parabole (mars 1651) où, par la méthode des indivisibles, le
conoïde parabolique est la moitié du cylindre de même base et de même hauteur, le
« fuseau parabolique » (solide engendré par un segment de parabole tournant autour
de sa base) les $\frac{8}{15}$ du cylindre qui l'enveloppe, puis la quadrature de
« toutes les paraboles à l'infini ».

\exfigure{paraboles}{Les courbes $y = x^n$ pour $n = 1, \ldots, 5$ ; l'aire ombrée, sous $y = x^2$, est le tiers du carré.}

\begin{enumerate}
\item Montrer que l'aire sous $y = x^n$ entre $0$ et $a$ est la fraction
  $\frac{1}{n+1}$ du rectangle circonscrit.
\item On fait tourner autour de l'axe des $x$ la région $0 \le y \le x^{1/n}$,
  $0 \le x \le a$ (un « conoïde » de la courbe $x = y^n$ autour de son axe). Montrer
  que son volume est au cylindre circonscrit comme $n$ à $n + 2$, et retrouver la
  moitié pour la parabole.
\item Montrer que le fuseau engendré par $y = h\bigl(1 - \frac{x^2}{b^2}\bigr)$,
  $-b \le x \le b$, tournant autour de l'axe des $x$, vaut les $\frac{8}{15}$ du
  cylindre de rayon $h$ et de longueur $2b$. Où a-t-on déjà rencontré le nombre
  $\frac{8}{15}$ dans ce livre, et pourquoi est-ce le même calcul ?
\end{enumerate}
\end{exercice}

\begin{solution}
\begin{enumerate}
\item $\int_0^a x^n\,dx = \frac{a^{n+1}}{n+1} = \frac{1}{n+1}\,(a\cdot a^n)$.
\item $V = \pi\int_0^a x^{2/n}\,dx = \pi\frac{a^{1+2/n}}{1 + 2/n} =
  \frac{n}{n+2}\,\pi a^{2/n}\,a$, et le cylindre circonscrit, de rayon $a^{1/n}$ et
  de longueur $a$, vaut $\pi a^{2/n}a$. Pour $n = 2$ : $\frac24 = \frac12$.
\item $V = \pi h^2\int_{-b}^{b}\bigl(1 - \frac{x^2}{b^2}\bigr)^2dx =
  2\pi h^2 b\int_0^1(1 - u^2)^2du = 2\pi h^2 b\cdot\frac{8}{15}$, et le cylindre vaut
  $2\pi h^2 b$. C'est l'intégrale $\int_0^1(1-u^2)^2du$ de la spirale des huit
  quinzièmes de Français 12357, où l'aire polaire $\frac12\int r^2d\theta$ fait
  apparaître, comme le volume de révolution $\pi\int y^2dx$, le carré de
  $1 - u^2$.
\end{enumerate}
\end{solution}

\begin{exercice}{Le beau problème de M. Roberval : un cylindre oblique mesuré au compas}{L3}{3}
\manuscrit{naf-5161}{42}{42}
La lettre sur les thèses du P. Fabri propose « ce beau problème de M. Roberval » : « Étant donné un
cylindre oblique, retrancher d'un seul trait de compas, sur un cylindre droit, une
superficie égale à la superficie du cylindre oblique donné. » On note
$E(k) = \int_0^{\pi/2}\sqrt{1 - k^2\sin^2u}\,du$ ($0 \le k < 1$) l'intégrale
elliptique complète de seconde espèce.

\exfigure{cylindre-oblique}{Le cylindre oblique de la question 1, et sa section droite, perpendiculaire aux génératrices.}

\begin{enumerate}
\item Le cylindre oblique a pour base un cercle de rayon $r$, ses génératrices ont
  la longueur $\ell$ et font l'angle $\beta$ avec le plan de la base. Montrer que
  sa surface latérale est $\ell$ fois le périmètre de sa section droite, que cette
  section est une ellipse de demi-axes $r$ et $r\sin\beta$, et que la surface vaut
  $4r\ell\,E(\cos\beta)$.
\item Un « trait de compas » sur un cylindre droit de rayon $R$ : la pointe fixe est
  en un point du cylindre, l'ouverture est $\rho > 2R$, et le trait est
  l'intersection du cylindre et de la sphère de centre la pointe et de rayon
  $\rho$. Montrer que la partie du cylindre intérieure à la sphère a pour aire
  $8R\rho\,E\bigl(\frac{2R}{\rho}\bigr)$.
\item En déduire qu'on résout le problème en choisissant $\rho = \sqrt{\frac{r\ell}{\cos\beta}}$
  et $R = \frac{\rho\cos\beta}{2}$, deux longueurs constructibles à la règle et au
  compas. Que se passe-t-il si le cylindre droit est imposé ?
\end{enumerate}
\end{exercice}

\begin{solution}
\begin{enumerate}
\item Découpons la surface latérale en bandes étroites le long des génératrices :
  une bande de largeur $ds$ mesurée perpendiculairement aux génératrices a l'aire
  $\ell\,ds$, et $ds$ est l'élément de longueur de la section droite ; d'où
  $\ell\times$ périmètre. Prenons la base dans le plan $z = 0$ et les génératrices
  dirigées par $w = (\cos\beta, 0, \sin\beta)$. La section droite est la projection
  orthogonale de la base sur le plan orthogonal à $w$ : le vecteur $(0, 1, 0)$ se
  projette sans changer de longueur, $(1, 0, 0)$ en un vecteur de longueur
  $\sqrt{1 - \cos^2\beta} = \sin\beta$. C'est donc une ellipse de demi-axes $r$ et
  $r\sin\beta$, d'excentricité $k = \cos\beta$. Son périmètre, avec la
  paramétrisation $(r\sin u, r\sin\beta\cos u)$, vaut
  $\int_0^{2\pi} r\sqrt{\cos^2u + \sin^2\beta\sin^2u}\,du = 4r\int_0^{\pi/2}
  \sqrt{1 - \cos^2\beta\sin^2u}\,du = 4rE(\cos\beta)$. La surface est
  $4r\ell E(\cos\beta)$.
\item Paramétrons le cylindre par $(R\cos\phi, R\sin\phi, z)$ et la pointe en
  $(R, 0, 0)$. La distance au carré vaut $2R^2(1 - \cos\phi) + z^2 =
  4R^2\sin^2\frac{\phi}{2} + z^2$. Pour $\rho > 2R$, chaque génératrice est coupée
  selon $|z| \le \sqrt{\rho^2 - 4R^2\sin^2(\phi/2)}$, et l'aire (élément
  $R\,d\phi\,dz$) vaut
  $\int_{-\pi}^{\pi} 2R\sqrt{\rho^2 - 4R^2\sin^2\frac\phi2}\,d\phi$. Avec
  $\phi = 2u$ et la parité, c'est
  $8R\rho\int_0^{\pi/2}\sqrt{1 - \frac{4R^2}{\rho^2}\sin^2u}\,du =
  8R\rho\,E\bigl(\frac{2R}{\rho}\bigr)$. Le trait est formé de deux courbes fermées
  qui bornent cette zone.
\item Les deux aires sont égales dès que $\frac{2R}{\rho} = \cos\beta$ et
  $8R\rho = 4r\ell$, soit $2R\rho = r\ell$ ; alors $\rho^2\cos\beta = r\ell$, d'où
  les valeurs annoncées, et $\rho > 2R$ puisque $\cos\beta < 1$. Une moyenne
  proportionnelle et une quatrième proportionnelle se construisent à la règle et
  au compas. Si le cylindre droit (donc $R$) est imposé, les deux conditions
  portent sur la seule inconnue $\rho$ et ne peuvent en général être remplies
  ensemble ; il faut alors comparer les deux intégrales elliptiques autrement. La
  lettre ne donne pas de solution ; celle-ci repose sur une notion bien
  postérieure.
\end{enumerate}
\end{solution}

\section*{Joseph-Louis Lagrange (1736--1813)}
\chapitre{lagrange}

Deux des exercices de ce chapitre viennent de pièces conservées dans les papiers
de Sophie Germain : un billet de Lagrange à « Mademoiselle Germaine », daté d'un
jour du calendrier républicain mais sans année (NAF 4073), et le résumé, de la
main de Sophie Germain, d'un mémoire de Lagrange sur l'équation $x^2 - ay^2 = 1$
paru dans les Mémoires de l'Académie de Turin (Français 9115). Le troisième vient
de NAF 5166, les « Papiers de Lagrange », où un mémoire intitulé « Eclaircissement
d'une difficulté singulière … Par J. L. Lagrange » tire un paradoxe de l'équation
de Laplace à la surface d'un sphéroïde.

\begin{exercice}{Dater un billet de Lagrange par le calendrier républicain}{Lycée}{1}
\manuscrit{naf-4073}{29}{29}
Un billet de Lagrange à « Mademoiselle Germaine » est daté « Paris ce 17 germinal
(mercredi) » et la prévient qu'il sera à ses ordres « le 19 et le 20. (vendredi, et
samedi) ». Il ne porte pas d'année. Dans le calendrier républicain, l'année compte
douze mois de trente jours (vendémiaire, brumaire, frimaire, nivôse, pluviôse,
ventôse, germinal, \ldots) suivis de jours complémentaires. Le 1\textsuperscript{er}
vendémiaire tombe, pour les ans I à XIV, les 22 septembre 1792, 1793, 1794, 23 septembre
1795, 22 septembre 1796, 1797, 1798, 23 septembre 1799, 1800, 1801, 1802, 24 septembre 1803,
23 septembre 1804 et 1805 ; le 22 septembre 1792 était un samedi.
\begin{enumerate}
\item La transcription note que le « 17 » s'accorde avec le reste du billet, « où le
  19 est un vendredi et le 20 un samedi ». Cet accord confirme-t-il la lecture 17 ?
\item Montrer que le 17 germinal tombe le même jour de la semaine que le
  1\textsuperscript{er} vendémiaire de la même année.
\item Déterminer le jour de la semaine du 1\textsuperscript{er} vendémiaire de chaque
  année de l'an I à l'an XIV (on comptera $365$ ou $366$ jours d'un
  1\textsuperscript{er} vendémiaire au suivant, selon la date). En déduire les dates
  possibles du billet.
\end{enumerate}
\end{exercice}

\begin{solution}
\begin{enumerate}
\item Non : si le 17 est un mercredi, le 19 est un vendredi et le 20 un samedi,
  quels que soient le mois et l'année. Les trois indications n'en font qu'une ; elles
  confirment seulement que les trois quantièmes ont été lus de façon cohérente.
\item Le 17 germinal est le jour $6\times30 + 17 = 197$ de l'année ; il suit le
  1\textsuperscript{er} vendémiaire de $196 = 28\times 7$ jours, donc tombe le même
  jour de la semaine.
\item D'un 1\textsuperscript{er} vendémiaire au suivant, il y a $365$ jours, plus un
  si l'intervalle contient un 29 février (1796 et 1804 sont bissextiles, 1800 ne l'est
  pas), plus le décalage de la date de septembre ($+1$, $0$ ou $-1$) ; le jour de la
  semaine avance d'autant de jours que ce nombre dépasse $364 = 52\times7$. On trouve : an I samedi, II dimanche, III lundi, IV mercredi (23 septembre 1795 :
  $366$ jours après le 22 septembre 1794), V jeudi, VI vendredi, VII samedi,
  VIII lundi, IX mardi, X mercredi, XI jeudi, XII samedi, XIII dimanche, XIV lundi. Les
  seuls mercredis sont les ans IV et X. Le 17 germinal an IV est le 6 avril 1796
  ($196$ jours après le 23 septembre 1795), le 17 germinal an X le 7 avril 1802 : le
  billet est de l'un de ces deux jours, et d'aucun autre.
\end{enumerate}
\end{solution}

\begin{exercice}{L'équation $x^2 - ay^2 = 1$ dans un extrait de Turin}{L1}{2}
\manuscrit{fr-9115}{749}{749}
La dernière page écrite du volume résume, d'après un mémoire de Lagrange des Mémoires
de Turin, l'étude de $x^2 - ay^2 = 1$, $a$ étant un entier positif non carré : le
lemme $(x^2 - ay^2)(x'^2 - ay'^2) = (xx' \pm ayy')^2 - a(xy' \pm yx')^2$ ; puis, si
$p^2 - aq^2 = 1$, les solutions
\[
x = \frac{(p + q\sqrt a)^m + (p - q\sqrt a)^m}{2}, \qquad
y = \frac{(p + q\sqrt a)^m - (p - q\sqrt a)^m}{2\sqrt a} ;
\]
enfin, « La Grange démontre ensuite »\footnote{« La Grange » est la graphie du
feuillet, courante à l'époque pour Lagrange ; on la garde dans la citation.} que si $(p,q)$ est la plus petite solution en
entiers positifs, toutes s'obtiennent ainsi.
\begin{enumerate}
\item Vérifier le lemme.
\item Montrer que $x$ et $y$ ainsi définis sont entiers et vérifient $x^2 - ay^2 = 1$.
  Pour $m = 2$, la copie écrit
  $q' = \frac{(p+q\sqrt a)^2 + (p - q\sqrt a)^2}{2\sqrt a}$ : calculer cette quantité,
  et la valeur juste de $q'$.
\item Pour $a = 2$ et $(p,q) = (3,2)$, donner les deux solutions suivantes. Pour
  $a = 13$, vérifier que $(649, 180)$ est solution, et calculer la suivante.
\item On démontre l'énoncé attribué à Lagrange ; on pose $\varepsilon = p + q\sqrt a$.
  Montrer que si $u^2 - av^2 = 1$ avec $u$, $v$ entiers et $\xi = u + v\sqrt a > 1$,
  alors $u > 0$ et $v > 0$. En déduire que $\varepsilon$ est le plus petit des nombres
  $\xi > 1$ de cette sorte.
\item Soit $(x,y)$ une solution en entiers positifs, $\eta = x + y\sqrt a$, et $m$
  l'entier tel que $\varepsilon^m \le \eta < \varepsilon^{m+1}$. Montrer que
  $\eta\,\varepsilon^{-m} = \eta\,(p - q\sqrt a)^m$ s'écrit $u + v\sqrt a$ avec $u$, $v$
  entiers et $u^2 - av^2 = 1$, et conclure.
\end{enumerate}
\end{exercice}

\begin{solution}
\begin{enumerate}
\item $(xx' + ayy')^2 - a(xy' + yx')^2 = x^2x'^2 + 2axx'yy' + a^2y^2y'^2 - ax^2y'^2 -
  2axx'yy' - ay^2x'^2 = x^2x'^2 - ax^2y'^2 - ay^2x'^2 + a^2y^2y'^2$, qui est
  $(x^2 - ay^2)(x'^2 - ay'^2)$ ; de même avec les signes inférieurs. C'est l'identité
  qu'on attribue à Brahmagupta.
\item Par récurrence sur $m$, $(p + q\sqrt a)^m = X_m + Y_m\sqrt a$ avec $X_m$, $Y_m$
  entiers, et $(p - q\sqrt a)^m = X_m - Y_m\sqrt a$ (on multiplie par
  $p \pm q\sqrt a$ : $X_{m+1} = pX_m + aqY_m$, $Y_{m+1} = qX_m + pY_m$). Donc
  $x = X_m$ et $y = Y_m$ sont entiers, et
  $x^2 - ay^2 = (x + y\sqrt a)(x - y\sqrt a) = \bigl((p+q\sqrt a)(p-q\sqrt a)\bigr)^m =
  (p^2 - aq^2)^m = 1$. Pour $m = 2$, $(p+q\sqrt a)^2 + (p-q\sqrt a)^2 = 2(p^2 + aq^2)$,
  et la quantité de la copie vaut $\frac{p^2 + aq^2}{\sqrt a}$, qui n'est pas même
  rationnelle ; il faut le signe moins :
  $q' = \frac{4pq\sqrt a}{2\sqrt a} = 2pq$, avec $p' = p^2 + aq^2$.
\item $a = 2$ : $(17, 12)$, puisque $17 = 9 + 8$ et $12 = 2\cdot3\cdot2$, puis
  $(3\cdot17 + 2\cdot2\cdot12,\ 2\cdot17 + 3\cdot12) = (99, 70)$ ; en effet
  $289 - 288 = 1$ et $9801 - 9800 = 1$. $a = 13$ : $649^2 = 421\,201$ et
  $13\times180^2 = 421\,200$. La suivante est $p' = 649^2 + 13\times180^2 = 842\,401$,
  $q' = 2\times649\times180 = 233\,640$.
\item $(u + v\sqrt a)(u - v\sqrt a) = 1$, donc $u - v\sqrt a = \frac1\xi \in\, ]0,1[$,
  et $u = \frac12\bigl(\xi + \frac1\xi\bigr) > 0$,
  $v\sqrt a = \frac12\bigl(\xi - \frac1\xi\bigr) > 0$. Les $\xi > 1$ sont donc
  exactement les $x + y\sqrt a$ des solutions en entiers positifs ; pour celles-ci,
  $x = \sqrt{1 + ay^2}$ croît avec $y$, et la plus petite solution donne le plus petit
  $\xi$, qui est $\varepsilon$.
\item Comme $\eta \ge \varepsilon > 1$, l'entier $m \ge 1$ existe. On a
  $\varepsilon^{-1} = p - q\sqrt a$, donc $\eta\varepsilon^{-m} = (x + y\sqrt a)(p -
  q\sqrt a)^m$, qui, développé, est de la forme $u + v\sqrt a$ avec $u$, $v$ entiers ;
  et $u^2 - av^2 = (x^2 - ay^2)(p^2 - aq^2)^m = 1$ par le lemme. Or
  $1 \le u + v\sqrt a < \varepsilon$. Si $u + v\sqrt a > 1$, le (a) en fait une
  solution positive plus petite que $\varepsilon$, ce qui est exclu ; donc
  $u + v\sqrt a = 1$ et $\eta = \varepsilon^m$ : $(x,y)$ est la solution de rang $m$
  de la formule. L'extrait ne rappelle pas que $a$ doit être positif et non carré ;
  c'est ce qui assure l'existence de $(p,q)$ et l'irrationalité de $\sqrt a$, utilisée
  en identifiant $x + y\sqrt a$.
\end{enumerate}
\end{solution}

\begin{exercice}{Un « paradoxe dans le Calcul intégral » : la couche sphérique uniforme}{L2}{3}
\manuscrit{naf-5166}{21}{33}
\manuscrit{naf-5166}{41}{41}
Le mémoire des vues 17 à 41, « Eclaircissement d'une difficulté singuliere qui se
rencontre dans le calcul de l'attraction des sphéroïdes très peu differens de la
sphere », dont le titre porte « Par J. L. Lagrange », examine l'équation que
« M. Laplace trouve » à la surface d'un sphéroïde homogène de rayon
$R = a(1+y')$, $y'$ très petit : $\frac V2 + a\frac{dV}{dr} = -\frac{2\pi a^2}{3}$.
Il écrit $V = \frac{4\pi a^3}{3r} + u$, où $r$ est la distance du point attiré au
centre, et
\[
u = a^3\int \rho\,y'\,d\pi'\,d\theta'\sin\theta', \qquad
\rho = \frac{1}{\sqrt{r^2 - 2ra\mu + a^2}},
\]
$\mu$ étant le cosinus de l'angle du rayon $r$ et du rayon de la
molécule.\footnote{Sur ces vues, $\pi'$ est la longitude de la molécule (notre
$\varpi'$), et $\pi$ désigne aussi l'angle de $180^\circ$ ; $d\pi'\,d\theta'\sin\theta'$
est l'élément d'aire $d\Omega'$ de la sphère unité, et l'intégrale porte sur toute
la sphère. Le $d$ du feuillet sert aussi pour les dérivées partielles :
$\frac{du}{dr}$ est notre $\frac{\partial u}{\partial r}$. On garde ses lettres $u$,
$\rho$, $\mu$, $y'$.} Au premier ordre en $y'$, $u$ est le potentiel d'une couche de
densité superficielle $ay'$ étalée sur la sphère de rayon $a$. En dérivant sous le
signe $\int$ et en faisant ensuite $r = a$, le facteur
$\frac\rho2 + a\frac{d\rho}{dr}$ est une « équation identique » (vue 27), d'où
$\frac u2 + a\frac{du}{dr} = 0$ ; mais pour $y'$ constant le calcul donne
$-2\pi a^2y'$ (vue 25), « ce qui est un paradoxe dans le Calcul integral »
(vue 31). On prend ici, comme aux n\textsuperscript{os} 3 et 7 du mémoire, $y'$
constant, et positif, et l'on pose $G(r) = \frac u2 + r\frac{\partial u}{\partial r}$.
\begin{enumerate}
\item (N\textsuperscript{o} 3, vue 25.) En portant le pôle des angles sur le rayon
  $r$, de sorte que $\mu = \cos\theta'$, montrer que
  \[
  \int\frac{d\pi'\,d\theta'\sin\theta'}{\sqrt{r^2 - 2ra\mu + a^2}}
  = \frac{2\pi\bigl((r+a) - |r-a|\bigr)}{ar},
  \]
  et en déduire $u$ pour $r > a$ et pour $r < a$. Le feuillet complète l'intégrale en
  $\frac{2\pi(r+a)}{ra} - \frac{2\pi(r-a)}{ra} = \frac{4\pi}{r}$ : quel cas
  traite-t-il sans le dire ? $u$ est-il continu en $r = a$ ?
\item Calculer $G(r)$ pour $r > a$ et pour $r < a$, et ses limites quand
  $r \to a^+$ et quand $r \to a^-$. Même question pour $\frac{\partial u}{\partial r}$,
  et pour la quantité $\frac u2 + a\frac{\partial u}{\partial r}$ du feuillet. Laquelle
  de ces limites le feuillet trouve-t-il aux vues 25 et 33 ?
\item La « valeur directe ». On fait $r = a$ avant d'intégrer, comme aux vues 23
  et 27. Calculer, pour $\mu < 1$, les valeurs de $\rho$ et de
  $\frac{\partial\rho}{\partial r}$ en $r = a$, et vérifier que
  $\frac\rho2 + a\frac{\partial\rho}{\partial r} = 0$. Le feuillet, tel que la
  transcription le lit, écrit en $r = a$ : $\rho = \frac{1}{a\sqrt2(1-\mu)}$ et
  $\frac{d\rho}{dr} = \frac{1}{2a\sqrt2(1-\mu)}$ ; corriger. Calculer ensuite les
  intégrales, prises en $r = a$, $u_d = a^3y'\int\rho\,d\Omega'$ et
  $u'_d = a^3y'\int\frac{\partial\rho}{\partial r}\,d\Omega'$, puis
  $G_d = \frac{u_d}{2} + a\,u'_d$, et comparer aux limites de la question 2.
\item « je ferai d'abord le calcul sans rien negliger » (vue 31). Montrer que, pour
  $r \ne a$, $r\frac{\partial\rho}{\partial r} = -\frac\rho2 - \frac{(r^2-a^2)\rho^3}{2}$,
  puis que $G(r) = -\frac{a^3(r^2-a^2)\,y'}{2}\int\rho^3\,d\Omega'$. Calculer
  $\int\rho^3\,d\Omega'$ pour $r > a$ et pour $r < a$ ; le feuillet donne
  $\frac{4\pi}{r(r^2-a^2)}$ « quelle que soit la valeur de $r$ » (vue 33) : qu'en
  penser ? Retrouver les résultats de la question 2. Enfin, $\delta \in\, ]0,2[$ étant
  fixé, montrer que la part de $(r^2-a^2)\int\rho^3\,d\Omega'$ qui vient des
  directions $\mu \le 1 - \delta$ tend vers $0$ quand $r \to a^+$, et dire où le
  raisonnement du n\textsuperscript{o} 5 (vue 29) est en défaut.
\item Laquelle des trois valeurs faut-il garder ? Le sphéroïde est ici la sphère
  homogène de rayon $a(1+y')$, et le point attiré est sur sa surface,
  $r = a(1+y)$ avec $y = y'$. Calculer exactement
  $\frac V2 + a\frac{\partial V}{\partial r}$ en ce point. Le n\textsuperscript{o} 10
  (vue 41) écrit, en négligeant les $y^2$,
  $\frac V2 + a\frac{dV}{dr} = -\frac{2\pi a^2}{3} + 2\pi a^2y + \frac u2 + a\frac{du}{dr}$ :
  vérifier cette égalité, et en déduire quelle valeur de
  $\frac u2 + a\frac{du}{dr}$ est la bonne, et pourquoi.
\end{enumerate}
\end{exercice}

\begin{solution}
\begin{enumerate}
\item Avec le pôle sur le rayon $r$, $d\Omega' = \sin\theta'\,d\theta'\,d\pi'$,
  $\mu = \cos\theta'$ et $d\mu = -\sin\theta'\,d\theta'$ ; l'intégrale vaut
  $2\pi\int_{-1}^{1}\frac{d\mu}{\sqrt{r^2 - 2ar\mu + a^2}}$. Une primitive en $\mu$ est
  $-\frac{1}{ar}\sqrt{r^2 - 2ar\mu + a^2}$, et
  \[
  2\pi\int_{-1}^{1}\frac{d\mu}{\sqrt{r^2 - 2ar\mu + a^2}}
  = \frac{2\pi}{ar}\Bigl(\sqrt{(r+a)^2} - \sqrt{(r-a)^2}\Bigr)
  = \frac{2\pi\bigl((r+a) - |r-a|\bigr)}{ar}.
  \]
  Elle vaut $\frac{4\pi}{r}$ pour $r \ge a$ et $\frac{4\pi}{a}$ pour $r \le a$. D'où
  \[
  u = \frac{4\pi a^3y'}{r}\quad(r \ge a), \qquad u = 4\pi a^2y'\quad(r \le a).
  \]
  En écrivant $r - a$ pour la racine en $\mu = 1$, le feuillet suppose $r \ge a$ :
  c'est le point extérieur à la sphère de rayon $a$. Les deux expressions valent
  $4\pi a^2y'$ en $r = a$ : $u$ est continu. À l'intérieur il est constant : une
  couche sphérique uniforme n'attire pas un point placé au dedans (théorème de
  Newton).
\item Pour $r > a$, $\frac{\partial u}{\partial r} = -\frac{4\pi a^3y'}{r^2}$ et
  $G(r) = \frac{2\pi a^3y'}{r} - \frac{4\pi a^3y'}{r} = -\frac{2\pi a^3y'}{r}$ ; pour
  $r < a$, $\frac{\partial u}{\partial r} = 0$ et $G(r) = 2\pi a^2y'$. Donc
  \[
  \lim_{r\to a^+}G = -2\pi a^2y', \qquad \lim_{r\to a^-}G = +2\pi a^2y', \qquad
  \lim_{r\to a^+}\frac{\partial u}{\partial r} = -4\pi ay', \qquad
  \lim_{r\to a^-}\frac{\partial u}{\partial r} = 0 .
  \]
  La quantité du feuillet diffère de $G$ de $(a - r)\frac{\partial u}{\partial r}$,
  qui tend vers $0$ : pour $r > a$ elle vaut
  $\frac{2\pi a^3y'}{r} - \frac{4\pi a^4y'}{r^2}$, et elle a les mêmes limites,
  $-2\pi a^2y'$ et $+2\pi a^2y'$. La dérivée radiale saute de $-4\pi ay'$, soit
  $-4\pi$ fois la densité $ay'$ de la couche : c'est la relation de saut d'un
  potentiel de simple couche. Les valeurs $-2\pi a^2y'$ de la vue 25 et
  $-2a^2\pi y'$ de la vue 33 sont la limite extérieure.
\item En $r = a$, $r^2 - 2ar\mu + a^2 = 2a^2(1-\mu)$, donc
  $\rho = \frac{1}{a\sqrt{2(1-\mu)}}$. Comme
  $\frac{\partial\rho}{\partial r} = -(r - a\mu)\rho^3$,
  \[
  \frac{\partial\rho}{\partial r}\Big|_{r=a}
  = -\frac{a(1-\mu)}{2\sqrt2\,a^3(1-\mu)^{3/2}}
  = -\frac{1}{2\sqrt2\,a^2\sqrt{1-\mu}},
  \qquad
  \frac\rho2 + a\frac{\partial\rho}{\partial r}
  = \frac{1}{2\sqrt2\,a\sqrt{1-\mu}} - \frac{1}{2\sqrt2\,a\sqrt{1-\mu}} = 0 .
  \]
  Les valeurs lues sur la vue 27 n'ont pas de signe moins (comme $\frac{du}{dr}$ à
  la vue 23), le radical n'y couvre que le $2$, et la dérivée y porte $a$ où il faut
  $a^2$ ; avec elles la somme ne s'annulerait pas. Comme
  $\int_{-1}^{1}\frac{d\mu}{\sqrt{1-\mu}} = 2\sqrt2$, les intégrales convergent et
  \[
  u_d = \frac{2\pi a^2y'}{\sqrt2}\cdot2\sqrt2 = 4\pi a^2y', \qquad
  u'_d = -\frac{2\pi ay'}{2\sqrt2}\cdot2\sqrt2 = -2\pi ay', \qquad
  G_d = 2\pi a^2y' - 2\pi a^2y' = 0 .
  \]
  $u_d$ est bien $u(a)$. Mais $u'_d = \frac12(-4\pi ay' + 0)$ et
  $G_d = \frac12(-2\pi a^2y' + 2\pi a^2y')$ : la valeur directe est la moyenne des
  deux limites latérales. Le calcul des vues 23 et 27, signes rétablis, ne se trompe
  pas sur cette intégrale ; il se trompe en la prenant pour la dérivée de $u$ en
  $r = a$, où $u$ n'est pas dérivable.
\item On a $r(r - a\mu) = \frac12(r^2 - 2ar\mu + a^2) + \frac12(r^2 - a^2)$ et
  $\rho^3(r^2 - 2ar\mu + a^2) = \rho$, d'où
  $r\frac{\partial\rho}{\partial r} = -r(r-a\mu)\rho^3
  = -\frac\rho2 - \frac{(r^2-a^2)\rho^3}{2}$. Pour $r$ dans un intervalle fermé qui
  ne contient pas $a$, $r^2 - 2ar\mu + a^2 \ge (r-a)^2$ est minoré par une constante
  positive, la fonction à intégrer et sa dérivée en $r$ sont continues, et l'on peut
  dériver sous le signe $\int$ : $G(r) = a^3y'\int\bigl(\frac\rho2 +
  r\frac{\partial\rho}{\partial r}\bigr)d\Omega' = -\frac{a^3(r^2-a^2)\,y'}{2}
  \int\rho^3\,d\Omega'$. Une primitive de $(r^2 - 2ar\mu + a^2)^{-3/2}$ en $\mu$ est
  $\frac{\rho}{ar}$, donc
  \[
  \int\rho^3\,d\Omega' = \frac{2\pi}{ar}\Bigl(\frac{1}{|r-a|} - \frac{1}{r+a}\Bigr)
  = \begin{cases}
  \frac{4\pi}{r(r^2-a^2)} & (r > a),\\
  \frac{4\pi}{a(a^2-r^2)} & (r < a).
  \end{cases}
  \]
  La valeur du feuillet ne vaut que pour $r > a$ : pour $r < a$ elle est négative,
  alors que $\rho^3 > 0$. On retrouve $G(r) = -\frac{2\pi a^3y'}{r}$ pour $r > a$ et
  $G(r) = +2\pi a^2y'$ pour $r < a$. Pour $\mu < 1$ fixé, $\rho$ reste borné quand
  $r \to a$, et $(r^2-a^2)\rho^3 \to 0$ ; pourtant $(r^2-a^2)\int\rho^3\,d\Omega'$ vaut
  $\frac{4\pi}{r}$ pour $r > a$ et tend vers $\frac{4\pi}{a}$. En
  $\mu = 1 - \delta$, $r^2 - 2ar\mu + a^2 = (r-a)^2 + 2ar\delta$, et la part des
  directions $\mu \le 1 - \delta$ est
  \[
  2\pi(r^2-a^2)\int_{-1}^{1-\delta}\rho^3\,d\mu
  = \frac{2\pi(r^2-a^2)}{ar}\Bigl(\frac{1}{\sqrt{(r-a)^2 + 2ar\delta}} - \frac{1}{r+a}\Bigr),
  \]
  dont la parenthèse tend vers $\frac{1}{a\sqrt{2\delta}} - \frac{1}{2a}$ : elle tend
  vers $0$. Toute l'intégrale se ramasse dans une calotte aussi petite qu'on veut
  autour de la direction du point attiré, là où la fonction à intégrer devient
  infinie ; la fonction $\frac{r(r^2-a^2)}{4\pi}\rho^3$, positive et d'intégrale $1$,
  est ce qu'on appelle aujourd'hui le noyau de Poisson de la boule. Le
  n\textsuperscript{o} 5 intègre la limite $\frac\rho2 + a\frac{d\rho}{dr}$, nulle
  pour $\mu < 1$, au lieu de prendre la limite des intégrales : il échange la limite
  $r \to a$ et l'intégration, ce que rien ne permet ici (la convergence n'est ni
  uniforme ni dominée près de $\mu = 1$). La somme « de toutes les valeurs
  successives de la differentielle » n'est pas en cause.
\item La sphère entière a pour potentiel extérieur
  $V = \frac{4\pi a^3(1+y')^3}{3r}$ ; en $r = a(1+y')$,
  \[
  \frac V2 + a\frac{\partial V}{\partial r}
  = \frac{2\pi a^2(1+y')^2}{3} - \frac{4\pi a^2(1+y')}{3}
  = -\frac{2\pi a^2}{3}\,(1 - y'^2),
  \]
  soit $-\frac{2\pi a^2}{3}$ au premier ordre, sans terme en $y'$. Par ailleurs, en
  $r = a(1+y)$, $\frac{2\pi a^3}{3r} - \frac{4\pi a^4}{3r^2}
  = \frac{2\pi a^2}{3}(1 - y) - \frac{4\pi a^2}{3}(1 - 2y) + O(y^2)
  = -\frac{2\pi a^2}{3} + 2\pi a^2y + O(y^2)$, ce qui est l'égalité du
  n\textsuperscript{o} 10. Il faut donc
  $\frac u2 + a\frac{du}{dr} = -2\pi a^2y'$ : c'est la limite extérieure. La limite
  intérieure donnerait $-\frac{2\pi a^2}{3} + 4\pi a^2y'$, et la valeur directe
  $-\frac{2\pi a^2}{3} + 2\pi a^2y'$, le résultat de la vue 23, qui n'est l'équation
  de Laplace que si $y = 0$, comme le feuillet le dit. La raison : le point attiré,
  à la distance $a(1+y') > a$, est hors de la sphère de rayon $a$ sur laquelle on a
  étalé la couche ; c'est le côté extérieur qui compte, et l'on vérifie que
  $V - \frac{4\pi a^3}{3r} = \frac{4\pi a^3}{3r}\bigl((1+y')^3 - 1\bigr)
  = \frac{4\pi a^3y'}{r} + O(y'^2)$ est bien la branche extérieure de $u$. C'est le
  remède du mémoire : calculer « sans rien negliger » pour $r > a$, comme aux
  n\textsuperscript{os} 6 et 7, et ne faire $r = a$ qu'à la fin.
\end{enumerate}

\exfigure{couche-spherique}{Questions 1 à 3 : $u$ est continu en $r = a$, mais $G$ y saute de $+2\pi a^2y'$ à $-2\pi a^2y'$ ; la valeur directe $G_d = 0$ est la moyenne des deux limites.}
\end{solution}

\section*{Sophie Germain (1776--1831)}
\chapitre{germain}

Deux volumes de ses papiers nourrissent ce chapitre. Français 9115, « Recueil de
dissertations et problèmes mathématiques et physiques », relie sans ordre de
sujet ni de date sept cent cinquante vues : la théorie des plaques élastiques,
avec un long mémoire « sur la question proposée par la première classe de
l'Institut » qui confronte ses sons à ceux que Chladni a notés, des pages sur la
chaleur, puis de l'arithmétique, dont les trois pages que la littérature appelle
le « manuscrit C ». Français 9118 est le dossier des lettres qu'elle reçut, de ses
brouillons à Gauss signés « Le Blanc » et de feuillets détachés de sa main. Les
exercices vont des plaques vibrantes à l'arithmétique des \emph{Disquisitiones
arithmeticae} de Gauss (1801) et au théorème de Fermat.

\begin{exercice}{Les sons de la plaque rectangulaire $9 : 8$}{Lycée}{1}
\manuscrit{fr-9115}{524}{533}
Au N\textsuperscript{o} 19 du long mémoire, elle compare aux notes de Chladni les sons
d'une plaque de verre rectangulaire dont les côtés sont comme $9$ à $8$. On note $T|L$
la figure formée de $T$ lignes nodales droites parallèles au petit côté (elles
partagent le grand) et de $L$ lignes parallèles au grand côté. Dans sa théorie, le son
de $T|L$ est proportionnel à $\frac{T^2}{9^2} + \frac{L^2}{8^2}$ ; en multipliant par
$81$, elle l'écrit $N = \frac{64T^2 + 81L^2}{64}$. On rappelle les intervalles : demi-ton
mineur $\frac{25}{24}$, demi-ton majeur $\frac{16}{15}$, ton mineur $\frac{10}{9}$, ton
majeur $\frac98$ ; un intervalle de rapport $r$ vaut $1200\log_2 r$ cents.

\exfigure{plaque}{Les figures $2|2$, $2|3$ et $3|3$, schématiques : $T$ lignes nodales partagent le grand côté, $L$ le petit.}

\begin{enumerate}
\item Calculer $N$, sous forme de fraction irréductible, pour $2|2$, $2|3$, $3|3$,
  $4|3$, $5|5$ et $6|2$.
\item Elle rapporte chaque nombre au nombre $8$ de la figure $2|2$ de la plaque carrée
  (figure 68 de Chladni), ce qui donne le rapport théorique $R = \frac N8$, et le compare
  à l'intervalle qu'elle entend entre cette figure et la figure $T|L$ : $\frac{15}{8}$
  pour $2|3$, $\frac83$ pour $3|3$, $\frac{32}{9}$ pour $4|3$, $8$ pour $5|5$, et « un
  peu moindre » que $\frac{50}{9}$ pour $6|2$. Pour chacun de ces cas, calculer le
  rapport de l'intervalle entendu à $R$, puis sa valeur en cents.
\item Pour $2|2$, elle dit que $R$ est plus grand que $\frac98$ « d'une quantité
  inappréciable ». Pour $3|3$, que l'écart est « bien moindre qu'un demi-ton mineur » ;
  pour $4|3$, qu'il est « d'un peu plus d'un demi-ton » ; pour $5|5$, qu'il est d'« un
  ton mineur » ; pour $6|2$, d'« environ un demi-ton ». Juger chaque phrase.
\item Elle conclut que, parmi ses dix-huit comparaisons, « il ne se trouve qu'un seul
  cas » où l'écart excède un demi-ton. Que peut-on dire de cette conclusion au vu des
  cas calculés ?
\end{enumerate}
\end{exercice}

\begin{solution}
\begin{enumerate}
\item $2|2$ : $\frac{256+324}{64} = \frac{145}{16}$ ; $2|3$ : $\frac{256+729}{64} =
  \frac{985}{64}$ ; $3|3$ : $\frac{576+729}{64} = \frac{1305}{64}$ ; $4|3$ :
  $\frac{1024+729}{64} = \frac{1753}{64}$ ; $5|5$ : $\frac{1600+2025}{64} =
  \frac{3625}{64}$ ; $6|2$ : $\frac{2304+324}{64} = \frac{2628}{64} = \frac{657}{16}$.
  Ce sont les nombres du feuillet.
\item $2|3$ : $R = \frac{985}{512}$ ; $\frac{15}{8} : \frac{985}{512} = \frac{960}{985}
  = \frac{192}{197}$, soit $-45$ cents (le son entendu est plus grave). $3|3$ :
  $R = \frac{1305}{512}$ ; $\frac83 : R = \frac{4096}{3915}$, soit $+78$ cents. $4|3$ :
  $R = \frac{1753}{512}$ ; $\frac{32}{9} : R = \frac{16384}{15777}$, soit $+65$ cents.
  $5|5$ : $R = \frac{3625}{512}$ ; $8 : R = \frac{4096}{3625}$, soit $+212$ cents.
  $6|2$ : $R = \frac{657}{128}$ ; $\frac{50}{9} : R = \frac{6400}{5913}$, soit
  $+137$ cents. Les produits en croix du feuillet ($960$ et $985$, $15777$ et $16384$,
  $3625$ et $4096$, $5913$ et $6400$) sont justes.
\item En cents, le demi-ton mineur vaut $71$, le demi-ton majeur $112$, le ton mineur
  $182$ et le ton majeur $204$. Pour $2|2$ : $\frac{145}{128} : \frac98 =
  \frac{145}{144}$, $12$ cents : la phrase est juste. Pour $3|3$ : $78 > 71$, l'écart
  dépasse le demi-ton mineur ; la phrase est fausse. Sa vérification écrite,
  $\frac{3\cdot1305}{8\cdot512} < \frac{25}{24}$, compare au demi-ton un rapport plus
  petit que $1$, ce qui est toujours vrai ; il fallait comparer
  $\frac{4096}{3915}$, et $4096\times24 = 98\,304 > 3915\times25 = 97\,875$. Pour
  $4|3$ : $65$ cents, c'est moins qu'un demi-ton mineur, et non « un peu plus d'un
  demi-ton » ; l'erreur est ici en défaveur de sa théorie. Pour $5|5$ : $212$ cents,
  plus qu'un ton majeur, a fortiori qu'un ton mineur ; elle le compte pourtant, à juste
  titre, comme dépassant le demi-ton. Pour $6|2$ : $137$ cents, plus qu'un demi-ton
  majeur, si l'on prend l'intervalle entendu pour $\frac{50}{9}$.
\item Sur les cas calculés, $5|5$ dépasse le demi-ton sans discussion, et $6|2$ aussi
  si l'on prend l'intervalle entendu pour $\frac{50}{9}$ ; comme elle ne le donne que
  pour « un peu moindre », l'écart réel peut être plus petit, mais le cas n'est pas
  acquis. La conclusion « un seul cas » est donc trop favorable. L'accord reste bon
  pour la plupart des cas, à l'échelle de la tolérance d'un demi-ton qu'elle s'est
  donnée. Les nombres $\frac{T^2}{a^2} + \frac{L^2}{b^2}$ sont exacts pour une plaque
  appuyée sur son contour ; pour la plaque libre de Chladni, ils
  ne sont qu'approchés, ce que le feuillet ne pouvait savoir.
\end{enumerate}
\end{solution}

\begin{exercice}{La cire sur les bords : une erreur d'unité et sa correction}{Lycée}{2}
\manuscrit{fr-9115}{416}{416}
\manuscrit{fr-9115}{539}{546}
Au paragraphe 6 du long mémoire, elle charge de cire les côtés d'une plaque carrée de verre de
$106$ mm de côté, pesant $1$ once $18$ grains, et en rend compte par une règle : une
dose de cire répartie le long d'un côté équivaut à un prolongement de la plaque de
même poids. On rappelle $1$ once $= 8$ gros $= 576$ grains. Pour la figure de deux
lignes nodales croisées, le son d'une plaque rectangulaire de côtés $a$ et $b$ est,
dans sa théorie, proportionnel à $\frac1{a^2} + \frac1{b^2}$.
\begin{enumerate}
\item Chaque dose pèse $1$ gros $2$ grains. Vérifier qu'elle est le huitième du poids
  de la plaque, et qu'elle équivaut à environ $13$ mm de plaque.
\item Une dose sur un côté : la plaque devient fictivement un rectangle de $119$ mm
  sur $106$ mm, « à très peu près $9 : 8$ ». Le feuillet écrit que le son est
  proportionnel à $\frac{9^2+8^2}{64} = \frac{145}{64}$, celui de la plaque non chargée
  à $2$, et que l'intervalle est $\frac{145}{128}$. Calculer l'intervalle juste, en
  prenant $8$ pour le côté du carré. D'où vient l'erreur ? Un petit feuillet relié plus
  haut dans le mémoire (vue 416) la corrige : « l'intervalle que j'ai exprimé par
  $\frac{N^2+n^2}{2n^2}$ doit l'être par $\frac{2N^2}{N^2+n^2}$ ». Vérifier, et
  comparer les deux valeurs en cents. Elle entend le son passer d'« un peu plus bas que
  sol 1 » à « entre mi 1 et fa 1 ».
\item Deux doses sur deux côtés contigus : la plaque devient un carré de $119$ mm.
  Quel intervalle prévoit-on ? L'erreur du 2 joue-t-elle ici ?
\item Deux doses sur deux côtés opposés : la plaque devient un rectangle « $10 : 8$ ».
  Le feuillet trouve $\frac{41}{32}$, « un peu plus que deux tons majeurs », et entend
  « de ré 1 à sol 1 » plus de deux tons majeurs. Calculer l'intervalle corrigé, et
  juger l'accord avec l'expérience (deux tons majeurs : $\frac{81}{64}$).
\item Une lame de verre de $200$ mm, pesant $4$ gros $12$ grains, montre deux lignes
  nodales transversales, l'intervalle qui les sépare étant double de celui qui les
  sépare des bouts. Chargée à un bout de $1$ gros $3$ grains, elle les montre à $12$ mm
  du bout chargé et à $121$ mm l'une de l'autre, $61$ mm restant jusqu'à l'autre bout.
  Appliquer la règle et comparer.
\end{enumerate}
\end{exercice}

\begin{solution}
\begin{enumerate}
\item La plaque pèse $576 + 18 = 594$ grains et la dose $72 + 2 = 74$ grains ; or
  $\frac{594}{8} = 74{,}25$. Une dose vaut donc le huitième de la plaque, soit une bande
  de $\frac{106}{8} = 13{,}25$ mm sur toute la longueur d'un côté.
\item En prenant pour unité le huitième du côté, le carré a pour nombre
  $\frac1{8^2} + \frac1{8^2} = \frac{2}{64}$ et le rectangle $\frac1{9^2} + \frac1{8^2} =
  \frac{145}{5184}$ ; le rapport du son de la plaque non chargée à celui de la plaque
  chargée est $\frac{2}{64}\cdot\frac{5184}{145} = \frac{162}{145} = 1{,}117$, la
  plaque chargée étant la plus grave. Le nombre $\frac{145}{64} = 81\bigl(\frac1{81} +
  \frac1{64}\bigr)$ est rapporté au grand côté du rectangle, le nombre
  $2 = 64\bigl(\frac1{64} + \frac1{64}\bigr)$ au côté du carré : les deux nombres ne sont
  pas dans la même unité, et, pris tels quels, ils feraient la plaque chargée plus
  aiguë. Avec $N = 9$ et $n = 8$, $\frac{N^2+n^2}{2n^2} = \frac{145}{128}$ et
  $\frac{2N^2}{N^2+n^2} = \frac{162}{145}$ : la correction de la vue 416 est exacte.
  En cents, $216$ contre $192$ ; l'une et l'autre valeurs sont voisines d'un ton
  ($200$ cents), comme l'intervalle entendu, qui va d'un peu moins d'un ton à un peu
  moins d'un ton et demi selon la note, entre mi et fa : la correction ne change presque rien, et le feuillet a raison d'écrire que les deux « ne different pas
  sensiblement ».
\item Le rapport de deux carrés de côtés $9$ et $8$ vaut $\frac{2/64}{2/81} =
  \frac{81}{64}$, deux tons majeurs ($408$ cents). Les deux nombres étant rapportés
  chacun à son propre côté, sans mélange d'unités, l'erreur du 2 ne joue pas.
\item Rectangle $10 \times 8$ : $\frac1{100} + \frac1{64} = \frac{164}{6400}$, et le
  carré $\frac{2}{64} = \frac{200}{6400}$ ; l'intervalle corrigé est
  $\frac{200}{164} = \frac{50}{41} = 1{,}220$, soit $344$ cents, moins que deux tons
  majeurs ($408$ cents), alors que le calcul du feuillet donnait $\frac{41}{32}$, soit
  $429$ cents. Avec les longueurs mêmes, $132$ et $106$ mm,
  $\frac{2\times132^2}{132^2 + 106^2} = 1{,}216$, soit $338$ cents. Elle entend plus de
  deux tons majeurs, de ré à un peu moins que sol, près d'une quarte ($498$ cents) : la
  correction abaisse la prévision de $85$ cents, au-dessous de deux tons majeurs, et
  l'éloigne de l'expérience ; la phrase de la vue 416 ne vaut que pour le premier cas. La lecture modernisée ajoute que
  l'équivalence d'une masse à une longueur n'est exacte qu'au premier ordre, pour de
  petites charges.
\item La lame pèse $4\times72 + 12 = 300$ grains, la cire $72 + 3 = 75$ grains, le
  quart ; elle équivaut à $\frac{200}{4} = 50$ mm de lame. La lame fictive a $250$ mm,
  partagée en $\frac14$, $\frac12$, $\frac14$ : lignes à $62{,}5$ mm et à $187{,}5$ mm de
  son bout fictif, soit à $12{,}5$ mm du bout chargé, à $125$ mm l'une de l'autre, et à
  $62{,}5$ mm de l'autre bout. On observe $12$, $121$ et $61$ mm, qui font $194$ mm, les
  $6$ mm restants étant pris, selon le feuillet, par la cire et les lignes : l'accord
  est bon, et une portion de $12$ mm rend le même son qu'une portion dix fois plus
  longue.
\end{enumerate}
\end{solution}

\begin{exercice}{L'équation de la plaque s'abaisse au second degré}{L2}{2}
\manuscrit{fr-9115}{399}{409}
\manuscrit{fr-9115}{469}{485}
Le paragraphe 1 du long mémoire sur les surfaces élastiques établit l'équation
$\frac{\partial^2z}{\partial t^2} + b\tau^2\Delta^2z = 0$ (A), où $\Delta$ est le
laplacien en $(x,y)$, $\Delta^2 = \Delta\Delta$, $\tau$ l'épaisseur et $b$ une
constante de la matière. Pour un « mouvement régulier »,
$z = Z(x,y)\sin\bigl(\frac{t}{\sqrt K} + \zeta\bigr)$, et, en posant $Kb\tau^2 = h^4$,
elle obtient $h^4\Delta^2Z = Z$ (B).
\begin{enumerate}
\item Montrer que $Z$ satisfait (B) dès qu'il satisfait $h^2\Delta Z = Z$ ou
  $h^2\Delta Z = -Z$ (l'équation « peut s'abaisser au second degré »), et que toute
  solution de (B) s'écrit $Z' + Z''$ avec $h^2\Delta Z' = Z'$ et $h^2\Delta Z'' = -Z''$.
\item Plaque carrée de côté $A$, origine en un coin. Montrer que
  $Z = \sin\frac{n\pi x}{A}\sin\frac{m\pi y}{A}$ satisfait (B) pour une valeur de $K$ à
  déterminer, et que le son, mesuré par $\sqrt{1/K}$, est proportionnel à $m^2 + n^2$.
  Vérifier que $Z = 0$ et $\Delta Z = 0$ sur les quatre côtés. Une membrane tendue de
  même forme, $\frac{\partial^2z}{\partial t^2} = c^2\Delta z$, a les mêmes figures :
  comment ses sons dépendent-ils de $m$ et de $n$ ?
\item Au N\textsuperscript{o} 15, avec $u = \frac{\pi x}{A}$ et $v = \frac{\pi y}{A}$,
  elle forme
  \[
  Z = \sin u\,\sin v\Bigl(\beta\cos^2u + \beta'\cos^2v - \frac{\beta+\beta'}{4}\Bigr),
  \]
  $\beta$ et $\beta'$ étant deux constantes (ses $b$ et $B$). Montrer que
  $Z = \frac\beta4\sin3u\,\sin v + \frac{\beta'}{4}\sin u\,\sin3v$, que $Z$ satisfait
  (B), et que son son est proportionnel à $10$.
\item Pour $\beta' = -\beta$, montrer que les deux diagonales du carré sont des lignes
  nodales. Pour $\beta' = 0$, décrire les lignes nodales.
\item On pose $r = \frac{\beta'}{\beta}$ ($\beta \neq 0$). Sur la médiane
  $y = \frac A2$, où $\cos v = 0$, la ligne nodale du facteur entre parenthèses passe
  aux points où $\cos^2u = \frac{\beta+\beta'}{4\beta}$ ; sur la médiane
  $x = \frac A2$, aux points où $\cos^2v = \frac{\beta+\beta'}{4\beta'}$. Montrer que
  ces deux sortes de points existent ensemble si et seulement si
  $\frac13 \le r \le 3$ ou $r = -1$.
\end{enumerate}
\end{exercice}

\begin{solution}
\begin{enumerate}
\item Les opérateurs $h^2\Delta - 1$ et $h^2\Delta + 1$ commutent, et
  $(h^2\Delta - 1)(h^2\Delta + 1)Z = h^4\Delta^2Z - Z$. Si $h^2\Delta Z = \pm Z$, alors
  $h^4\Delta^2Z = h^2\Delta(\pm Z) = \pm h^2\Delta Z = Z$. Réciproquement, soit $Z$ une
  solution de (B), $Z' = \frac12(Z + h^2\Delta Z)$ et $Z'' = \frac12(Z - h^2\Delta Z)$ :
  $Z' + Z'' = Z$, $h^2\Delta Z' = \frac12(h^2\Delta Z + h^4\Delta^2Z) =
  \frac12(h^2\Delta Z + Z) = Z'$, et de même $h^2\Delta Z'' = -Z''$. La seconde de ces
  équations est celle des mouvements réguliers des surfaces tendues, qu'on nomme
  aujourd'hui équation de Helmholtz.
\item $\Delta Z = -\frac{\pi^2}{A^2}(m^2+n^2)Z$ et
  $\Delta^2Z = \frac{\pi^4}{A^4}(m^2+n^2)^2Z$ ; (B) a lieu si
  $Kb\tau^2\frac{\pi^4}{A^4}(m^2+n^2)^2 = 1$, c'est-à-dire
  $\sqrt{1/K} = \frac{\pi^2(m^2+n^2)\tau\sqrt b}{A^2}$, proportionnel à $m^2 + n^2$ :
  c'est la formule du feuillet. Sur les côtés $x = 0$ et $x = A$,
  $\sin\frac{n\pi x}{A} = 0$, donc $Z = 0$, et $\Delta Z$, proportionnel à $Z$, s'annule
  aussi ; de même sur $y = 0$ et $y = A$. Ce sont, sur un bord droit, les conditions
  d'un appui simple dans la théorie actuelle : ces produits de sinus sont exactement
  les modes de la plaque carrée appuyée sur son contour (solution de Navier). Pour la
  membrane, $\Omega^2 = c^2\frac{\pi^2}{A^2}(m^2+n^2)$ : le son est proportionnel à
  $\sqrt{m^2+n^2}$. Pour un même nombre d'onde, la plaque a pour sons les carrés de
  ceux de la membrane, ce que le mémoire dit au N\textsuperscript{o} 19.
\item $\sin3u = 3\sin u - 4\sin^3u = \sin u(4\cos^2u - 1)$, donc
  $\frac\beta4\sin3u\sin v + \frac{\beta'}4\sin u\sin3v =
  \sin u\sin v\bigl(\beta\cos^2u - \frac\beta4 + \beta'\cos^2v - \frac{\beta'}4\bigr)$,
  qui est $Z$. Chacun des deux termes, modes $(3,1)$ et $(1,3)$ du 2, a
  $\Delta = -\frac{\pi^2}{A^2}\cdot10$ : ils ont le même $K$, leur somme satisfait (B),
  et le son est proportionnel à $3^2 + 1^2 = 10$, comme elle l'écrit.
\item Pour $\beta' = -\beta$, $Z = \beta\sin u\sin v(\cos^2u - \cos^2v) =
  \beta\sin u\sin v(\cos u - \cos v)(\cos u + \cos v)$. Sur $[0,\pi]$, $\cos u = \cos v$
  équivaut à $u = v$, soit $y = x$, et $\cos u = -\cos v$ à $u = \pi - v$, soit
  $x + y = A$ : les deux diagonales sont nodales. Pour $\beta' = 0$,
  $Z = \frac\beta4\sin3u\sin v$ ne s'annule à l'intérieur que pour $\sin3u = 0$, soit
  $x = \frac A3$ et $x = \frac{2A}{3}$ : deux droites parallèles à un côté, au tiers de
  la largeur.
\item Le premier point existe si et seulement si $0 \le \frac{1+r}{4} \le 1$, soit
  $-1 \le r \le 3$. Le second, avec $\frac{\beta+\beta'}{4\beta'} = \frac{1+r}{4r}$
  ($r \neq 0$) : pour $r > 0$, il faut $1 + r \le 4r$, soit $r \ge \frac13$ ; pour
  $r < 0$, $\frac{1+r}{4r} \ge 0$ exige $1 + r \le 0$, soit $r \le -1$, et
  $\frac{1+r}{4r} \le 1$ a alors lieu. Les deux ensemble : $\frac13 \le r \le 3$, ou
  $r = -1$, cas où les deux points se réduisent au centre du carré. Si $r$ est positif
  et hors de $[\frac13, 3]$, ou négatif et différent de $-1$, une seule des médianes
  est coupée, ce que le feuillet dit en termes de « lignes sinueuses ».
\end{enumerate}
\end{solution}

\begin{exercice}{$p^4 + 4$, $p^4 + q^4$ et les nombres de Fermat}{Lycée}{1}
\manuscrit{fr-9118}{91}{93}
Un feuillet détaché de sa main dit : « Aucun nombre de la forme $p^4 + 4$ excepté 5
n'est un nombre premier, car $p^4 + 4 = (p^2 - 2)^2 + 4p^2$ et par conséquent ces
nombres sont, de plusieurs manières, la somme de deux quarrés. » Puis :
$p^4 + q^4 = (p^2 - q^2)^2 + 2(pq)^2 = (p^2+q^2)^2 - 2(pq)^2$, et « les nombres
$2^{2^i} + 1$ ne sont qu'un cas particulier de la forme $p^4 + q^4$ ».
\begin{enumerate}
\item Vérifier les identités du feuillet.
\item Démontrer l'identité $a^4 + 4b^4 = (a^2 + 2b^2 - 2ab)(a^2 + 2b^2 + 2ab)$, et en
  déduire directement que $p^4 + 4$ ($p \ge 1$ entier) n'est premier que pour $p = 1$.
\item L'argument du feuillet par deux sommes de deux carrés demande que les deux
  décompositions $p^4 + 4 = (p^2)^2 + 2^2 = (p^2 - 2)^2 + (2p)^2$ soient
  différentes. Pour quelles valeurs de $p$ coïncident-elles ? Le feuillet en
  exclut-il assez ?
\item Pour quels $i$ le nombre $2^{2^i} + 1$ est-il de la forme $p^4 + q^4$ avec
  $p, q$ entiers positifs ?
\end{enumerate}
\end{exercice}

\begin{solution}
\begin{enumerate}
\item $(p^2 - 2)^2 + 4p^2 = p^4 - 4p^2 + 4 + 4p^2 = p^4 + 4$ ;
  $(p^2 - q^2)^2 + 2p^2q^2 = p^4 + q^4 = (p^2 + q^2)^2 - 2p^2q^2$.
\item $(a^2 + 2b^2)^2 - (2ab)^2 = a^4 + 4a^2b^2 + 4b^4 - 4a^2b^2 = a^4 + 4b^4$. Avec
  $b = 1$ : $p^4 + 4 = (p^2 - 2p + 2)(p^2 + 2p + 2)$. Le second facteur vaut au moins
  $5$ ; le premier, $(p-1)^2 + 1$, vaut $1$ seulement pour $p = 1$. Pour $p \ge 2$
  les deux facteurs sont $\ge 2$ et le nombre est composé ; pour $p = 1$ on a $5$,
  premier. (L'identité porte aujourd'hui le nom de Sophie Germain ; le feuillet
  donne l'autre argument.)
\item Les paires $\{p^2, 2\}$ et $\{|p^2 - 2|, 2p\}$ coïncident si $p^2 = 2p$ et
  $|p^2 - 2| = 2$, soit $p = 2$, ou si $p^2 = |p^2 - 2|$ et $2 = 2p$, soit $p = 1$.
  Le feuillet exclut $p = 1$ (« excepté 5 ») mais pas $p = 2$ ; ce cas ne menace pas
  la conclusion puisque $2^4 + 4 = 20$ est pair, mais la restriction manque.
\item $2^{2^i} + 1 = \bigl(2^{2^{i-2}}\bigr)^4 + 1^4$ pour $i \ge 2$ : $17 = 2^4 + 1$,
  $257 = 4^4 + 1$, $65\,537 = 16^4 + 1$. Pour $i = 0$ et $i = 1$ ($3$ et $5$), il
  faudrait $p^4 + q^4 \le 5$ avec $p, q \ge 1$, soit $2$ : impossible. L'hypothèse
  $i \ge 2$ est implicite sur le feuillet.
\end{enumerate}
\end{solution}

\begin{exercice}{Le nombre 2 et les formes $8k \pm 3$ ; une extension fausse}{L2}{2}
\manuscrit{fr-9115}{704}{704}
\manuscrit{fr-9115}{720}{723}
On note $\bigl(\frac ap\bigr)$ le symbole de Legendre ($p$ premier impair, $p \nmid a$),
qui vaut $1$ si $a$ est un carré modulo $p$ (« résidu ») et $-1$ sinon. On admet
$\bigl(\frac{-1}{p}\bigr) = (-1)^{\frac{p-1}{2}}$ et la loi de réciprocité quadratique :
pour $p$, $q$ premiers impairs distincts,
$\bigl(\frac qp\bigr)\bigl(\frac pq\bigr) = (-1)^{\frac{p-1}{2}\cdot\frac{q-1}{2}}$.
\begin{enumerate}
\item Une note (vue 704) donne « une autre démonstration de la propriété du nombre 2
  d'être résidu ou non résidu des nombres premiers suivant leurs formes par rapport à
  8 ». Elle part de $-1 \equiv 2\cdot\frac{p-1}{2} \pmod p$. Soit $q$ un facteur premier
  de $\frac{p-1}{2}$ : montrer que $\bigl(\frac pq\bigr) = 1$, puis que
  $\bigl(\frac qp\bigr) = 1$ si $q \equiv 1 \pmod 4$ et $\bigl(\frac{-q}{p}\bigr) = 1$ si
  $q \equiv 3 \pmod 4$.
\item En déduire $\bigl(\frac2p\bigr) = -1$ pour $p \equiv 3 \pmod 8$ et
  $\bigl(\frac2p\bigr) = 1$ pour $p \equiv 7 \pmod 8$. Vérifier sur $p = 11$, $19$ et
  $23$ en dressant la liste des carrés.
\item Soit $p \equiv 1 \pmod 4$, et $2^e$ la plus grande puissance de $2$ qui divise
  $p - 1$. Montrer que la méthode donne $\bigl(\frac2p\bigr) = 1$ si $e$ est impair, et
  qu'elle « n'apprend rien » si $e$ est pair. Dresser la liste des nombres premiers
  inférieurs à $60$ qu'elle ne couvre pas.
\item Un billet (vue 722) démontre que tout nombre de la forme $4n+2$ est non-résidu
  d'un nombre premier. Justifier : $4n + 3$ a un facteur premier $q \equiv 3 \pmod 4$,
  et $4n + 2 \equiv -1 \pmod q$.
\item Un autre feuillet (vue 720) étend une proposition qu'il lit dans « l'auteur » :
  « tout nombre de la forme $4n+1$ pris négativement est non résidu de quelque nombre
  premier moindre que lui et de la forme $4n+3$ ». Montrer que $-5$, $-17$ et $-21$
  sont résidus de tous les nombres premiers de la forme $4k+3$ plus petits qu'eux et
  qui ne les divisent pas.
\end{enumerate}
\end{exercice}

\begin{solution}
\begin{enumerate}
\item $q$ divise $p - 1$, donc $p \equiv 1 \pmod q$ et $\bigl(\frac pq\bigr) =
  \bigl(\frac1q\bigr) = 1$. Par la réciprocité, $\bigl(\frac qp\bigr) =
  (-1)^{\frac{p-1}{2}\cdot\frac{q-1}{2}}$. Si $q \equiv 1 \pmod 4$, l'exposant est pair
  et $\bigl(\frac qp\bigr) = 1$. Si $q \equiv 3 \pmod 4$, $\frac{q-1}{2}$ est impair et
  $\bigl(\frac qp\bigr) = (-1)^{\frac{p-1}{2}} = \bigl(\frac{-1}{p}\bigr)$, d'où
  $\bigl(\frac{-q}{p}\bigr) = \bigl(\frac{-1}{p}\bigr)\bigl(\frac qp\bigr) = 1$. La note
  emploie ces faits sans les justifier : sa démonstration suppose la réciprocité pour
  deux nombres premiers impairs.
\item Si $p \equiv 3 \pmod 4$, le 1 donne $\bigl(\frac qp\bigr) = 1$ pour
  $q \equiv 1$ et $\bigl(\frac qp\bigr) = -1$ pour $q \equiv 3 \pmod 4$. Pour
  $p = 8K + 3$, $\frac{p-1}{2} = 4K + 1 \equiv 1 \pmod 4$ a un nombre pair de facteurs
  premiers (comptés avec multiplicité) congrus à $3$ modulo $4$ : il est résidu. Comme
  $-1$ est non-résidu, $-1 \equiv 2(4K+1)$ donne $\bigl(\frac2p\bigr) = -1$. Pour
  $p = 8K + 7$, $\frac{p-1}{2} = 4K + 3$ en a un nombre impair : il est non-résidu, et
  $-1 = \bigl(\frac2p\bigr)\cdot(-1)$ donne $\bigl(\frac2p\bigr) = 1$. Modulo $11$, les
  carrés non nuls sont $1, 4, 9, 5, 3$ : $2$ n'y est pas ($11 \equiv 3$). Modulo $19$ :
  $1, 4, 9, 16, 6, 17, 11, 7, 5$ : $2$ n'y est pas ($19 \equiv 3$). Modulo $23$,
  $5^2 = 25 \equiv 2$ ($23 \equiv 7$).
\item Écrivons $\frac{p-1}{2} = 2^{e-1}m$, $m$ impair ; ici $-1$ est résidu. Par le 1,
  tout facteur premier $q$ de $m$ a $\bigl(\frac qp\bigr) = 1$, puisque
  $\bigl(\frac{-1}{p}\bigr) = 1$ ; donc $\bigl(\frac mp\bigr) = 1$. Si $e - 1$ est pair,
  $2^{e-1}$ est un carré, $\frac{p-1}{2}$ est résidu, et $-1 \equiv 2\cdot\frac{p-1}{2}$
  donne $\bigl(\frac2p\bigr) = 1$. Si $e - 1$ est impair,
  $\bigl(\frac{(p-1)/2}{p}\bigr) = \bigl(\frac2p\bigr)$, et la relation devient
  $1 = \bigl(\frac2p\bigr)^2$, qui n'apprend rien : c'est ce que dit la dernière ligne
  de la note (« cette méthode n'apprend rien pour le cas $p = 2^{2m}n+1$ »). Sont dans
  ce cas les $p \equiv 5 \pmod 8$ ($e = 2$) et ceux où $e \ge 4$ est pair. En dessous
  de $60$ : $5$, $13$, $17$, $29$, $37$, $53$ ($41$, où $e = 3$, est couvert).
\item $4n + 3$ est impair et congru à $3$ modulo $4$ ; s'il n'avait que des facteurs
  premiers $\equiv 1 \pmod 4$, il serait $\equiv 1$. Il a donc un facteur premier
  $q \equiv 3 \pmod 4$, qui ne divise pas $4n+2$ (sinon il diviserait $1$). Alors
  $4n + 2 = (4n + 3) - 1 \equiv -1 \pmod q$, et $\bigl(\frac{-1}{q}\bigr) = -1$ :
  $4n + 2$ est non-résidu de $q$. (Le billet prend $4n + 2 + (2k+1)^2$ ; c'est ici le
  cas $k = 0$.)
\item $-5 \equiv 1 = 1^2 \pmod 3$, et $3$ est le seul nombre premier $4k+3$ inférieur
  à $5$. Pour $-17$ : $-17 \equiv 1 \pmod 3$, $\equiv 4 = 2^2 \pmod 7$,
  $\equiv 5 \equiv 4^2 \pmod{11}$. Pour $-21$ : $3$ et $7$ divisent $21$ ; $-21 \equiv
  1 \pmod{11}$ et $-21 \equiv 17 \equiv 6^2 \pmod{19}$. Dans les trois cas, aucun
  nombre premier de la forme voulue ne convient : l'extension est fausse. Selon la
  lecture modernisée, ce sont, avec $3$, les seules exceptions en dessous de
  $200\,000$, et c'est la restriction à la forme $4k+3$ qui rend l'énoncé faux.
\end{enumerate}
\end{solution}

\begin{exercice}{Le manuscrit C : deux cas et trois lacunes}{L3}{3}
\manuscrit{fr-9115}{708}{710}
Les trois pages du « manuscrit C » énoncent : « $n$ étant un nombre impair plus grand
que l'unité, il est impossible de satisfaire en nombres entiers à l'équation
$2z^{2n} = y^{2n} + x^{2n}$ » ; elles en déduisent le théorème de Fermat pour les
exposants pairs, $Z^{2n} = X^{2n} + Y^{2n}$ impossible pour $n > 1$, et démontrent le
premier énoncé en supposant $x$, $y$, $z$ premiers entre eux.
\begin{enumerate}
\item Le cas $n = 1$. Vérifier l'identité $(P^2 - Q^2 - 2PQ)^2 + (P^2 - Q^2 + 2PQ)^2 =
  2(P^2 + Q^2)^2$, qui ouvre la démonstration, et en tirer deux solutions de
  $x^2 + y^2 = 2z^2$ en entiers premiers entre eux avec $x^2 \neq y^2$. Pourquoi
  l'énoncé exclut-il $n = 1$ ?
\item Montrer que $x = y = z = 1$ est solution pour tout $n$. La démonstration finit
  ainsi : dans le dernier cas, $x^n = x$ et $y^n = y$, et « ce cas est donc le seul
  dans lequel il soit possible de satisfaire à l'équation », c'est-à-dire $n = 1$. Où
  est la faute ? Corriger l'énoncé.
\item Montrer que si $x$, $y$, $z$ sont premiers entre eux et $x^{2n} + y^{2n} =
  2z^{2n}$, ils sont impairs.
\item La page écrit $x^{2n} + y^{2n} = (x^2 + y^2)\Phi$, avec
  $\Phi = x^{2n-2} - x^{2n-4}y^2 + \cdots + y^{2n-2}$ ($n$ impair), puis $z = z'z''$,
  $x^2 + y^2 = 2z'^{2n}$ et $\Phi = z''^{2n}$, ce qui suppose $x^2 + y^2$ et $\Phi$
  premiers entre eux, au facteur $2$ près. Montrer que $\Phi \equiv ny^{2n-2}
  \pmod{x^2+y^2}$ et que le plus grand commun diviseur de $x^2+y^2$ et de $\Phi$ divise
  $n$. Calculer $\Phi$ et ce diviseur pour $n = 5$, $x = 1$, $y = 3$. Conclure.
\item Plus loin, avec $x = p^2 - q^2 - 2pq$ et $y = p^2 - q^2 + 2pq$, la page affirme
  $2(y^2 - x^2) > (x+y)^2$, puis emploie $(x+y)^2 \ge y^2$. Montrer que
  $2(y^2 - x^2) - (x+y)^2 = (x+y)(y - 3x)$. Tester les deux inégalités pour
  $(p,q) = (3,2)$, puis pour la solution obtenue en changeant le signe de $x$.
  Conclure.
\item Le cas $n = 2$, que l'énoncé laisse de côté. On admet le théorème de Fermat :
  $X^4 - Y^4 = W^2$ n'a pas de solution en entiers avec $XYW \neq 0$.
  Montrer que si $x^4 + y^4 = 2z^4$, alors $(z^2)^4 - (xy)^4 =
  \bigl(\frac{x^4-y^4}{2}\bigr)^2$, et en déduire les solutions en entiers premiers
  entre eux. Pour le second théorème, la page passe d'une égalité $Z^n \pm Y^n =
  2x'^{2n}$ à une égalité entre les bases en divisant par $Z \pm Y$ : montrer que ce
  pas ne vaut pas pour $n = 2$, et démontrer directement que $Z^4 = X^4 + Y^4$ est
  impossible avec $XYZ \neq 0$.
\end{enumerate}
\end{exercice}

\begin{solution}
\begin{enumerate}
\item Avec $a = P^2 - Q^2$ et $b = 2PQ$ : $(a-b)^2 + (a+b)^2 = 2(a^2 + b^2)$ et
  $a^2 + b^2 = (P^2 + Q^2)^2$. Pour $(P,Q) = (2,1)$ : $x = -1$, $y = 7$, $z = 5$, et
  $1 + 49 = 2\cdot25$. Pour $(3,2)$ : $x = -7$, $y = 17$, $z = 13$, et
  $49 + 289 = 2\cdot169$. L'équation de degré $2$ a une infinité de solutions
  primitives non triviales : il faut exclure $n = 1$.
\item $1 + 1 = 2\cdot1$. L'égalité $x^n = x$ n'entraîne pas $n = 1$ : elle a lieu pour
  tout $n$ si $x = \pm1$. La démonstration aboutit donc, dans son dernier cas, à
  $n = 1$ ou $x^2 = y^2 = 1$, et la conclusion oublie la seconde possibilité. L'énoncé
  correct est : pour $n$ impair $> 1$, les seules solutions en entiers premiers entre
  eux sont $x^2 = y^2 = z^2 = 1$. Sous cette forme, il est vrai ; il n'a été démontré
  qu'en 1997, par Darmon et Merel, par de tout autres moyens.
\item Si $x$ et $y$ étaient pairs tous deux, $2^{2n}$ diviserait $2z^{2n}$, donc $z$
  serait pair, et $x$, $y$, $z$ ne seraient pas premiers entre eux. Si l'un est pair et
  l'autre impair, $x^{2n} + y^{2n}$ est impair, alors que $2z^{2n}$ est pair. Donc $x$
  et $y$ sont impairs, leurs puissances paires valent $1$ modulo $8$, et
  $2z^{2n} \equiv 2 \pmod 8$ donne $z^{2n} \equiv 1 \pmod 4$ : $z$ est impair.
\item Modulo $x^2 + y^2$, $x^2 \equiv -y^2$, et le terme
  $(-1)^jx^{2(n-1-j)}y^{2j}$ de $\Phi$ est congru à
  $(-1)^j(-1)^{n-1-j}y^{2n-2} = y^{2n-2}$ puisque $n - 1$ est pair ; les $n$ termes
  donnent $\Phi \equiv ny^{2n-2}$. Un diviseur commun $d$ de $x^2+y^2$ et de $\Phi$
  divise donc $ny^{2n-2}$ ; il est premier avec $y$ (un facteur premier de $y$ et de
  $x^2 + y^2$ diviserait $x$), donc $d \mid n$. Pour $n = 5$, $x = 1$, $y = 3$ :
  $x^2 + y^2 = 10$ et $\Phi = 1 - 9 + 81 - 729 + 6561 = 5905 = 5\times1181$ ; le plus
  grand commun diviseur est $5$. Le couple $(1,3)$ n'est pas une solution, mais il
  montre que le pgcd peut valoir un facteur premier de $n$ de la forme $4k+1$ (les
  facteurs premiers impairs de $x^2 + y^2$ sont de cette forme) : le partage de la page
  n'est justifié que si aucun facteur premier de $n$ n'est $\equiv 1 \pmod 4$. C'est
  une première lacune.
\item $2(y^2 - x^2) - (x+y)^2 = (x+y)\bigl(2(y-x) - (x+y)\bigr) = (x+y)(y - 3x)$. Pour
  $(p,q) = (3,2)$ : $x = -7$, $y = 17$. Alors $2(y^2 - x^2) = 480 > (x+y)^2 = 100$, mais
  $(x+y)^2 = 100 < y^2 = 289$ : la seconde inégalité tombe. Pour $x = 7$, $y = 17$,
  autre solution de $x^2 + y^2 = 2\cdot13^2$ : $480 < (x+y)^2 = 576$, c'est la première
  qui tombe. La chaîne ne conclut que si $x$ et $y$ sont de même signe et
  $|y| > 3|x|$ : c'est une deuxième lacune ; l'oubli de la solution triviale (2) en
  est une troisième, et la lecture modernisée en relève une quatrième (deux valeurs
  d'un même produit, l'une par signe, supposées divisibles ensemble par le même
  nombre).
\item $\bigl(\frac{x^4+y^4}{2}\bigr)^2 - (x^2y^2)^2 = \bigl(\frac{x^4-y^4}{2}\bigr)^2$,
  et $\frac{x^4+y^4}{2} = z^4$ : c'est l'identité. Pour une solution primitive, $x$ et
  $y$ sont impairs (par le 3, qui vaut pour $n = 2$), donc non nuls, et
  $W = \frac{x^4-y^4}{2}$ est entier. Par le théorème admis, $W = 0$, d'où
  $x^4 = y^4 = z^4$ et, les nombres étant premiers entre eux, $x^2 = y^2 = z^2 = 1$ :
  pour $n = 2$ aussi, il n'y a que la solution triviale. Pour le second théorème, la
  page utilise que $Z + Y$ divise $Z^n + Y^n$, vrai pour $n$ impair seulement : pour
  $n = 2$, $Z = 3$, $Y = 1$, $Z + Y = 4$ ne divise pas $Z^2 + Y^2 = 10$. Mais
  $Z^4 = X^4 + Y^4$ s'écrit $Z^4 - Y^4 = (X^2)^2$ avec $X^2 \neq 0$, ce que le théorème
  admis exclut : le cas $n = 2$ est dû à Fermat. Pour $n$ impair $> 1$, le second
  théorème contient un cas du théorème démontré par Wiles en 1995.
\end{enumerate}
\end{solution}

\section*{Guillaume Libri (1803--1869)}
\chapitre{libri}

Libri paraît dans ces volumes comme correspondant de Sophie Germain. Une note
signée « G. Libri » (1824), conservée en Français 9118, porte trois congruences
pour un nombre premier $p = 4k + 1$ ; une lettre qu'il lui écrit de Florence en
1825 (NAF 4073) annonce deux résultats sur la théorie de la chaleur. C'est aussi
lui qui décrit, dans l'album NAF 9544, les manuscrits de Roberval d'où vient un
exercice du chapitre de Roberval.

\begin{exercice}{« Un nombre quelconque de constantes arbitraires »}{L2}{2}
\manuscrit{naf-4073}{39}{39}
Libri écrit de Florence en 1825 : « J'ai envoyé un mémoire à l'Institut sur la
théorie de la chaleur, où je crois avoir démontré qu'on peut donner à l'intégrale
d'une équation différentielle linéaire un nombre quelconque de constantes
arbitraires. »
\begin{enumerate}
\item Soit l'équation $y^{(n)} + a_{n-1}(t)y^{(n-1)} + \cdots + a_0(t)y = 0$, à
  coefficients continus sur un intervalle $I$, et $t_0 \in I$. On admet le théorème
  de Cauchy--Lipschitz linéaire (existence et unicité de la solution de conditions
  initiales données en $t_0$). Montrer que l'application
  $y \mapsto (y(t_0), \dots, y^{(n-1)}(t_0))$ est un isomorphisme de l'espace des
  solutions sur $\mathbf{R}^n$. Combien de constantes arbitraires « essentielles » la
  solution générale contient-elle ?
\item Pour $y'' + y = 0$, on écrit $y = A\cos t + B\sin t + C\cos(t - \gamma)$, avec
  quatre paramètres $A$, $B$, $C$, $\gamma$. L'ensemble des fonctions obtenues est-il
  plus grand que $\{A\cos t + B\sin t\}$ ? Expliquer en quel sens une écriture peut
  contenir « un nombre quelconque de constantes ».
\item Conclure sur l'annonce de Libri.
\end{enumerate}
\end{exercice}

\begin{solution}
\begin{enumerate}
\item L'application est linéaire (la dérivation l'est). Elle est injective par
  l'unicité : une solution de conditions initiales nulles est nulle. Elle est
  surjective par l'existence : tout vecteur de $\mathbf{R}^n$ est la donnée initiale
  d'une solution. C'est donc un isomorphisme, et l'espace des solutions est de
  dimension $n$ : la solution générale s'écrit $c_1y_1 + \cdots + c_ny_n$ pour une base
  $(y_i)$, avec exactement $n$ constantes essentielles, où $n$ est l'ordre.
\item $\cos(t - \gamma) = \cos\gamma\cos t + \sin\gamma\sin t$, donc
  $y = (A + C\cos\gamma)\cos t + (B + C\sin\gamma)\sin t$ : on n'obtient rien de plus.
  Les paramètres $C$ et $\gamma$ ne sont pas essentiels, puisque des valeurs
  différentes donnent la même fonction. Une écriture peut contenir autant de
  paramètres qu'on veut (par exemple dans une solution écrite comme intégrale
  définie, où des bornes et un contour sont libres), sans que l'ensemble décrit
  grandisse.
\item Telle quelle, l'annonce est fausse : le nombre des constantes essentielles est un
  invariant de l'équation, égal à son ordre. Elle ne peut être vraie que de constantes
  non essentielles ; le mémoire n'est pas dans le volume et la lettre ne permet pas de
  dire ce qu'il démontrait.
\end{enumerate}
\end{solution}

\begin{exercice}{La chaleur comme vibration : équation parabolique ou hyperbolique ?}{L3}{2}
\manuscrit{naf-4073}{39}{39}
Libri ajoute : « J'ai aussi déduit les équations du mouvement de la chaleur de
l'hypothèse des vibrations calorifiques. » L'équation de Fourier
$u_t = k\,u_{xx}$ ($k > 0$) est parabolique ; une vibration conduit à une équation
hyperbolique. On étudie un intermédiaire, l'équation « du télégraphiste »
$\tau u_{tt} + u_t = k\,u_{xx}$, avec un temps de relaxation $\tau > 0$.
\begin{enumerate}
\item Chercher des solutions $u = e^{iqx + st}$ ($q$ réel). Pour l'équation de Fourier,
  montrer que $s = -kq^2$. Pour l'équation avec $\tau$, trouver les deux valeurs de $s$
  et montrer que l'une tend vers $-kq^2$ quand $\tau \to 0$, l'autre vers $-\infty$.
\item Pour l'équation de Fourier, la solution de donnée initiale $u(x,0) = \delta_0$ est
  $\frac{1}{\sqrt{4\pi kt}}e^{-x^2/(4kt)}$. Montrer qu'elle est non nulle en tout point
  dès $t > 0$ : la vitesse de propagation est infinie. Est-ce compatible avec une
  propagation par vibrations ?
\item Pour $\tau > 0$, écrire l'équation sous la forme $u_{tt} - c^2u_{xx} = -\frac1\tau u_t$
  et donner $c$. Que se passe-t-il pour les grandes fréquences ($|q| \to \infty$) ?
  Pourquoi aucun calcul exact ne fait-il passer d'une équation des ondes à celle de
  Fourier ?
\end{enumerate}
\end{exercice}

\begin{solution}
\begin{enumerate}
\item Fourier : $s = k(iq)^2 = -kq^2$. Avec $\tau$ : $\tau s^2 + s + kq^2 = 0$, d'où
  $s_\pm = \frac{-1 \pm\sqrt{1 - 4\tau kq^2}}{2\tau}$. Pour $\tau \to 0$,
  $\sqrt{1 - 4\tau kq^2} = 1 - 2\tau kq^2 + O(\tau^2)$, donc
  $s_+ = -kq^2 + O(\tau)$ et $s_- \sim -\frac{1}{\tau} \to -\infty$ : le second mode
  s'amortit instantanément à la limite, le premier redonne Fourier.
\item $e^{-x^2/(4kt)} > 0$ pour tout $x$ et tout $t > 0$ : une source ponctuelle échauffe
  aussitôt toute la droite. Une propagation par vibrations d'un milieu, elle, se fait à
  vitesse finie ; l'équation de Fourier ne peut donc pas en être la conséquence exacte.
\item En divisant par $\tau$ : $u_{tt} - \frac{k}{\tau}u_{xx} = -\frac1\tau u_t$, équation
  des ondes amortie de célérité $c = \sqrt{k/\tau}$ : les perturbations se propagent à
  vitesse finie $c$. Pour $4\tau kq^2 > 1$, $s_\pm = -\frac{1}{2\tau} \pm
  i\frac{\sqrt{4\tau kq^2 - 1}}{2\tau}$ : les grandes fréquences oscillent (se propagent
  à la vitesse $\approx c$) en s'amortissant comme $e^{-t/2\tau}$, au lieu de décroître
  comme $e^{-kq^2t}$. Les deux équations sont de types différents (hyperbolique et
  parabolique) ; on ne passe de l'une à l'autre que par une limite $\tau \to 0$, où la
  vitesse $c$ devient infinie. Une dérivation « à partir de vibrations » ne peut donc
  être exacte ; ce que valait le mémoire de Libri, la lettre ne permet pas de le dire.
\end{enumerate}
\end{solution}

\begin{exercice}{Le coefficient binomial central et la décomposition en deux carrés}{L3}{3}
\manuscrit{fr-9118}{65}{65}
Une note signée G. Libri (1824) porte trois congruences pour un nombre premier
$p = 4k + 1 = a^2 + b^2$, $a$ impair, $b$ pair. La dernière est
$\frac21\cdot\frac62\cdot\frac{10}{3}\cdots\frac{4k-2}{k} \equiv \pm2a \pmod p$, le signe
étant $+$ si $a \equiv 1 \pmod 4$ ; les deux autres portent le même second membre
$\pm 2b$ : $(k+1)(k+2)\cdots 2k \equiv \pm 2b$ et
$\bigl[(k+1)(k+2)\cdots2k\bigr]^2 \equiv \pm 2b \pmod p$.
\begin{enumerate}
\item Montrer que le produit $\prod_{j=1}^k\frac{4j-2}{j}$ vaut $\binom{2k}{k}$.
\item Vérifier la congruence $\binom{2k}{k} \equiv 2a \pmod p$ ($a \equiv 1 \pmod 4$) pour
  $p = 13, 17, 29, 37, 41$ (on admet ce théorème de Gauss dans le cas général).
\item Montrer que la connaissance de $\binom{2k}{k} \bmod p$ détermine $a$ dès que
  $p > 16$, et en déduire un procédé pour écrire $p$ comme somme de deux carrés.
\item On admet le théorème de Wilson, $(p-1)! \equiv -1 \pmod p$. Montrer que
  $\bigl((2k)!\bigr)^2 \equiv -1 \pmod p$, que $\bigl(\frac ab\bigr)^2 \equiv -1$, et
  en déduire, à partir de la congruence de Gauss, que
  $\bigl[(k+1)\cdots 2k\bigr]^2 \equiv \pm 2b \pmod p$.
\item Montrer que la première congruence, sans le carré, est fausse.
\end{enumerate}
\end{exercice}

\begin{solution}
\begin{enumerate}
\item $\prod_{j=1}^k(4j - 2) = 2^k\prod(2j-1) = 2^k\cdot\frac{(2k)!}{2^kk!} =
  \frac{(2k)!}{k!}$, et en divisant par $k!$ on obtient $\binom{2k}{k}$.
\item $p = 13$ ($k = 3$) : $\binom63 = 20 \equiv 7$ ; $13 = 9 + 4$, $a = -3 \equiv 1
  \pmod 4$, $2a = -6 \equiv 7$. $p = 17$ ($k = 4$) : $\binom84 = 70 \equiv 2$, $a = 1$.
  $p = 29$ ($k = 7$) : $\binom{14}7 = 3432 \equiv 10$, $a = 5$. $p = 37$ ($k = 9$) :
  $\binom{18}{9} = 48\,620 \equiv 2$, $a = 1$. $p = 41$ ($k = 10$) :
  $\binom{20}{10} = 184\,756 \equiv 10$, $a = 5$.
\item $a^2 < p$, donc $|2a| < 2\sqrt p$. Si $p > 16$, $4\sqrt p < p$ : deux valeurs
  distinctes de $2a$ dans $]-2\sqrt p, 2\sqrt p[$ ne sont pas congrues modulo $p$, et
  le résidu de $\binom{2k}k$ (ramené entre $-p/2$ et $p/2$) donne $2a$, puis
  $b = \sqrt{p - a^2}$. C'est un algorithme direct de décomposition.
\item $(p-1)! = \prod_{j=1}^{2k}j\,(p - j) \equiv (-1)^{2k}\bigl((2k)!\bigr)^2 =
  \bigl((2k)!\bigr)^2$, donc $((2k)!)^2 \equiv -1$. Et $a^2 \equiv -b^2$, $b$ inversible
  modulo $p$ : $(a/b)^2 \equiv -1$. Les racines carrées de $-1$ modulo $p$ étant
  $\pm$ l'une d'elles, $(2k)! \equiv \varepsilon\frac ab$ avec $\varepsilon = \pm1$. Alors
  $\bigl[(k+1)\cdots2k\bigr]^2 = \frac{((2k)!)^2}{(k!)^2} \equiv \frac{-1}{(k!)^2}$ ;
  comme $\binom{2k}{k} = \frac{(2k)!}{(k!)^2} \equiv 2a$,
  $\frac{1}{(k!)^2} \equiv \frac{2a}{(2k)!} \equiv \frac{2a b}{\varepsilon a} = 2\varepsilon b$,
  et le carré vaut $-2\varepsilon b \equiv \pm2b$. Vérification : $p = 13$, $4\cdot5\cdot6 =
  120 \equiv 3$, $3^2 = 9 \equiv -4 = -2b$.
\item Pour $p = 5$ ($k = 1$, $b = 2$) : $(k+1)\cdots2k = 2$, alors que $\pm2b = \pm4$,
  soit $4$ ou $1$ modulo $5$. Pour $p = 13$ : $120 \equiv 3 \not\equiv \pm 4$. La
  première congruence est fausse ; la note porte le même second membre pour une
  quantité et son carré, et la faute peut être du feuillet comme de sa lecture.
\end{enumerate}
\end{solution}

\section*{Bibliographie}

Ouvrages cités dans les énoncés, les solutions et les présentations des chapitres,
tels que les lectures modernisées les nomment ; aucun n'est une source des
exercices, qui ne reposent que sur les feuillets.

\begin{itemize}
\item Autolycus, \emph{De la sphère en mouvement} ; Théodose, \emph{Des habitations} ;
  Euclide, \emph{Phénomènes} --- dans les versions de Maurolico copiées en
  Latin 17859.
\item Ernst Florens Friedrich Chladni, \emph{Traité d'acoustique}, Paris, 1809.
\item Henri Darmon et Loïc Merel, article de 1997 sur l'équation
  $a^\ell + b^\ell = 2c^\ell$, \emph{Journal für die reine und angewandte Mathematik}.
\item René Descartes, \emph{Discours de la méthode} et \emph{La Géométrie},
  1637.
\item Le P. Fabri, \emph{Thèses de géométrie}, discutées dans une lettre de NAF 5161.
\item Joseph Fourier, \emph{Théorie analytique de la chaleur}, 1822.
\item Carl Friedrich Gauss, \emph{Disquisitiones arithmeticae}, 1801
  (article 357).
\item Christiaan Huygens, \emph{De ratiociniis in ludo aleae}, 1657.
\item Joseph-Louis Lagrange, mémoire sur l'équation $x^2 - ay^2 = 1$, \emph{Mémoires}
  de l'Académie de Turin, t. 4 (résumé en Français 9115).
\item Reinhard Laubenbacher et David Pengelley, article de 2010 dans \emph{Historia
  Mathematica} sur les manuscrits de Sophie Germain relatifs au théorème de Fermat
  (désignation « manuscrit C »).
\item Guglielmo Libri, \emph{Description des manuscrits de Roberval conservés à la
  Bibliothèque nationale} (NAF 9544).
\item Francesco Maurolico, volume sur la sphère (\emph{Sphériques} de Théodose, de
  Ménélas et de Maurolico, textes grecs, tables et problèmes), Messine, Petrus
  Spira, 1558, dont Latin 17859 est une copie ; \emph{Opuscula mathematica}, 1575.
\item Ménélas d'Alexandrie, \emph{Sphériques}, livre III.
\item Marin Mersenne, \emph{Harmonie universelle}.
\item Claude-Louis Navier, mémoire sur la flexion des plaques élastiques, 1820.
\item Blaise Pascal, lettre à Pierre de Fermat du 29 juillet 1654.
\item Regiomontanus (Johannes de Monteregio), table des tangentes (« table
  féconde ») et règles de calcul commentées par Maurolico.
\item Gilles Personne de Roberval, traité de mécanique « des plans inclinés », cité par
  une note marginale de Latin 11195.
\item François Viète et son école, \emph{De recognitione et emendatione
  aequationum} (notation des espèces et loi des homogènes).
\item Andrew Wiles, démonstration du théorème de Fermat, \emph{Annals of
  Mathematics}, 1995.
\end{itemize}


\end{document}
